% Lesson 12 — Writing: Writes compositions about anticipating one's second language needs at work — Anglais Tle A — Cours signé OnBuch+
% Programme officiel MINESEC. Compiler avec : tectonic EN12-writing-writes-compositions-about-anticipating-one-s-se.tex
\newcommand{\DOCMATIERE}{Anglais}
\newcommand{\DOCNIVEAU}{Tle A}
\newcommand{\DOCMODULE}{Chapter 1 : Family and Social Life}
\newcommand{\DOCLECON}{EN12}
\newcommand{\DOCTITRE}{Lesson 12 — Writing: Writes compositions about anticipating one's second language needs at work}
% OnBuch+ — préambule « cours signés » ENRICHI (compilé avec tectonic / XeLaTeX)
% Système visuel « Pop » d'OnBuch+ : fond crème (jamais blanc), bordures encre
% épaisses, ombres « dures » décalées, titres Archivo Black en capitales.
% Version MATHÉMATIQUES : graphes (pgfplots/tikz), boîtes pédagogiques
% (définition, propriété, méthode, attention, le savais-tu, exercice résolu,
% exercice, TP, bilan), page de garde et signature OnBuch+ ancrée en bas.
%
% Macros attendues dans main.tex AVANT \input{preamble} :
%   \DOCMATIERE  \DOCNIVEAU  \DOCMODULE  \DOCLECON (numéro)  \DOCTITRE
\documentclass[11pt]{article}

\usepackage{fontspec}
\usepackage[english,french]{babel} % français = langue principale ; l'anglais via \eng{..} / textebox
\usepackage[a4paper,margin=2cm,top=3.4cm,bottom=2.8cm,headheight=1.4cm,footskip=1.3cm]{geometry}
\usepackage{amsmath,amssymb,amsfonts}
\usepackage{mathtools} % \xmapsto, \coloneqq... souvent utilisés par les modèles
\usepackage{cancel} % simplifications barrées (bilans de forces)
% \abs{z} (module, valeur absolue) et \norm{u} (norme), utilisés par les modèles
\providecommand{\abs}{}\let\abs\relax\DeclarePairedDelimiter{\abs}{\lvert}{\rvert}
\providecommand{\norm}{}\let\norm\relax\DeclarePairedDelimiter{\norm}{\lVert}{\rVert}
\usepackage[table]{xcolor} % [table] : fournit aussi \rowcolors (alternance de lignes)
\usepackage{pagecolor}
\usepackage{tikz}
\usetikzlibrary{arrows.meta,positioning,calc,shapes.geometric,shapes.misc,shapes.arrows,decorations.pathmorphing,decorations.pathreplacing,decorations.markings,patterns,shadows,mindmap,trees,fit,backgrounds,matrix,intersections,angles,quotes}
\usepackage{pgfplots}
\usepackage[version=4]{mhchem} % équations nucléaires \ce{^{235}_{92}U}
\usepackage[european,siunitx]{circuitikz} % schémas électriques
\pgfplotsset{compat=1.18}
\usepgfplotslibrary{fillbetween,groupplots} % aires entre courbes (calcul intégral)
\usepackage{enumitem}
\usepackage{fancyhdr}
% [most] charge toutes les bibliothèques tcolorbox (listings, minted, raster,
% documentation...) et gonfle inutilement la mémoire TeX sur un document long
% (~300 boîtes) : on ne charge que ce qui est réellement utilisé.
\usepackage[skins,breakable]{tcolorbox}
\usepackage{graphicx}
\usepackage{colortbl}
\usepackage{booktabs}
\usepackage{tabularx}
\usepackage{diagbox} % case d'en-tête diagonale des tableaux à double entrée
% Colonnes extensibles centrée / gauche / droite (C, L, R), souvent utilisées par les modèles
\newcolumntype{C}{>{\centering\arraybackslash}X}
\newcolumntype{L}{>{\raggedright\arraybackslash}X}
\newcolumntype{R}{>{\raggedleft\arraybackslash}X}
\usepackage{array}
\usepackage{multicol}
\usepackage{multirow}
\usepackage{siunitx}
% parse-numbers=false : accepte aussi \SI{2,5 \times 10^{-3}}{...}
\sisetup{locale=FR,output-decimal-marker={,},group-minimum-digits=4,parse-numbers=false}
\DeclareSIUnit\ohm{\ensuremath{\symup{\Omega}}} % U+2126 absent de la police
\DeclareSIUnit\division{div} % réglages d'oscilloscope (ms/div, V/div)
\DeclareSIUnit\tr{tr} % tours (vitesse de rotation, ex. \SI{1500}{\tr\per\minute})
\DeclareSIUnit\molar{M} % concentration molaire (ex. \SI{3}{\molar}, \SI{1}{\micro\molar})
\DeclareSIUnit\an{an} % année (le modèle confond parfois avec \year, un compteur TeX) — ex. \SI{5}{\kilo\meter\per\an}
\DeclareSIUnit\FCFA{FCFA} % franc CFA (ex. \SI{500}{\FCFA\per\kilo\gram})
\DeclareSIUnit\degreeBrix{\textdegree Brix} % teneur en sucre d'un sirop/jus (réfractométrie)
\DeclareSIUnit\month{mois} % le modèle utilise parfois \month, un compteur TeX, comme unité de durée
\usepackage{qrcode}
% \degree : commande fréquemment utilisée par les modèles pour un angle ou une
% température en dehors de \SI{}{} (non fournie par les paquets chargés).
\providecommand{\degree}{^{\circ}}
% \qed (fin de démonstration), utilisé par les modèles sans amsthm
\providecommand{\qed}{\hfill$\square$}
% \barre : double barre des valeurs interdites dans un tableau de variations (tkz-tab)
\providecommand{\barre}{\ensuremath{\|}}
% environnement proof (amsthm non chargé) utilisé par les modèles
\ifdefined\proof\else\newenvironment{proof}[1][Démonstration]{\par\noindent\textit{#1.}\ }{\qed\par}\fi
\usepackage{lastpage}
\usepackage{hyperref}
\hypersetup{hidelinks}

% ============================================================
% Palette « Pop » OnBuch+ — mêmes hex que la classe P (plus_ui.dart)
% ============================================================
\definecolor{poporange}{HTML}{F59321}
\definecolor{popdark}{HTML}{E07A0C}
\definecolor{popcream}{HTML}{FFF6E8}
\definecolor{popink}{HTML}{1C1714}
\definecolor{poppurple}{HTML}{7A5AE0}
\definecolor{popgreen}{HTML}{2E9E6A}
\definecolor{poppink}{HTML}{E0457B}
\definecolor{popblue}{HTML}{2F7DE1}
\definecolor{popgold}{HTML}{C98A12}
\definecolor{popmuted}{HTML}{8A7F76}
\definecolor{popline}{HTML}{EADFD2}
\definecolor{popgray}{HTML}{8A7F76}
\colorlet{popgrey}{popgray}
% Alias tolérants (le modèle nomme parfois les couleurs différemment)
\colorlet{popwhite}{white}
\colorlet{popblack}{popink}
\colorlet{popred}{poppink}
\colorlet{popyellow}{popgold}
% Teintes claires (fonds de tableaux/encadrés) — le modèle nomme parfois
% les couleurs avec un suffixe L (ex. popblueL) sans les avoir définies.
\colorlet{poporangeL}{poporange!20}
\colorlet{popdarkL}{popdark!20}
\colorlet{poppurpleL}{poppurple!20}
\colorlet{popgreenL}{popgreen!20}
\colorlet{poppinkL}{poppink!20}
\colorlet{popblueL}{popblue!20}
\colorlet{popgoldL}{popgold!20}
\colorlet{popredL}{popred!20}
\colorlet{popgrayL}{popgray!20}
% Teintes claires pour fonds de tableaux / zones
\colorlet{popblueL}{popblue!10!white}
\colorlet{poporangeL}{poporange!14!white}
\colorlet{popgreenL}{popgreen!12!white}
\colorlet{poppurpleL}{poppurple!12!white}
\colorlet{poppinkL}{poppink!12!white}
\colorlet{popgoldL}{popgold!14!white}

\pagecolor{popcream}

% ============================================================
% Typographie
% ============================================================
\defaultfontfeatures{Ligatures=TeX}
\setmainfont{PlusJakartaSans-Medium.ttf}[Path=../../../fonts/,BoldFont=PlusJakartaSans-Bold.ttf]
% Maths en sans-serif (Fira Math) avec chiffres et lettres droites de Plus
% Jakarta, pour que les formules se fondent dans le texte.
\usepackage[math-style=ISO,bold-style=ISO]{unicode-math}
\setmathfont{FiraMath-Regular.otf}[Path=../../../fonts/]
\setmathfont{PlusJakartaSans-Medium.ttf}[Path=../../../fonts/,range={up/num,up/{Latin,latin},\mathup}]
\usepackage{icomma} % virgule décimale sans espace en mode math
% Caractères mathématiques Unicode absents des polices de texte (Plus Jakarta,
% Archivo Black) : rendus par leur équivalent mathématique, en texte comme en
% formule — sinon ils s'affichent en carré vide (ex. « Arithmétique dans ℤ »).
\usepackage{newunicodechar}
\newunicodechar{ℝ}{\ensuremath{\mathbb{R}}}
\newunicodechar{ℤ}{\ensuremath{\mathbb{Z}}}
\newunicodechar{ℂ}{\ensuremath{\mathbb{C}}}
\newunicodechar{ℕ}{\ensuremath{\mathbb{N}}}
\newunicodechar{ℚ}{\ensuremath{\mathbb{Q}}}
\newunicodechar{𝔻}{\ensuremath{\mathbb{D}}}
\newunicodechar{α}{\ensuremath{\alpha}}
\newunicodechar{β}{\ensuremath{\beta}}
\newunicodechar{γ}{\ensuremath{\gamma}}
\newunicodechar{δ}{\ensuremath{\delta}}
\newunicodechar{ε}{\ensuremath{\varepsilon}}
\newunicodechar{θ}{\ensuremath{\theta}}
\newunicodechar{λ}{\ensuremath{\lambda}}
\newunicodechar{μ}{\ensuremath{\mu}}
\newunicodechar{σ}{\ensuremath{\sigma}}
\newunicodechar{τ}{\ensuremath{\tau}}
\newunicodechar{φ}{\ensuremath{\varphi}}
\newunicodechar{ψ}{\ensuremath{\psi}}
\newunicodechar{ω}{\ensuremath{\omega}}
\newunicodechar{Σ}{\ensuremath{\Sigma}}
% majuscules produites par \MakeUppercase dans les titres (α-aminés -> Α)
\newunicodechar{Α}{\ensuremath{\alpha}}
\newunicodechar{Β}{\ensuremath{\beta}}
\newunicodechar{Γ}{\ensuremath{\Gamma}}
\newunicodechar{Δ}{\ensuremath{\Delta}}
\newunicodechar{Ε}{\ensuremath{\varepsilon}}
\newunicodechar{Θ}{\ensuremath{\Theta}}
\newunicodechar{Λ}{\ensuremath{\Lambda}}
\newunicodechar{Μ}{\ensuremath{\mu}}
\newunicodechar{Τ}{\ensuremath{\tau}}
\newunicodechar{Φ}{\ensuremath{\Phi}}
\newunicodechar{Ψ}{\ensuremath{\Psi}}
\newunicodechar{Ω}{\ensuremath{\Omega}}
\newunicodechar{↦}{\ensuremath{\mapsto}}
\newunicodechar{∈}{\ensuremath{\in}}
\newunicodechar{∉}{\ensuremath{\notin}}
\newunicodechar{∧}{\ensuremath{\wedge}}
\newunicodechar{⇔}{\ensuremath{\Leftrightarrow}}
\newunicodechar{⟺}{\ensuremath{\Longleftrightarrow}}
\newunicodechar{⇒}{\ensuremath{\Rightarrow}}
\newunicodechar{⟹}{\ensuremath{\Longrightarrow}}
\newunicodechar{∘}{\ensuremath{\circ}}
\newunicodechar{∪}{\ensuremath{\cup}}
\newunicodechar{∩}{\ensuremath{\cap}}
\newunicodechar{≡}{\ensuremath{\equiv}}
\newunicodechar{⊂}{\ensuremath{\subset}}
\newunicodechar{⊥}{\ensuremath{\perp}}
\newunicodechar{∃}{\ensuremath{\exists}}
\newunicodechar{∀}{\ensuremath{\forall}}
\newunicodechar{∛}{\ensuremath{\sqrt[3]{\,}}}
\newunicodechar{∜}{\ensuremath{\sqrt[4]{\,}}}
\newunicodechar{⃗}{\ensuremath{{}^{\to}}}
\newunicodechar{ⁿ}{\ifmmode^{n}\else\textsuperscript{\ensuremath{n}}\fi}
\newunicodechar{ˣ}{\ifmmode^{x}\else\textsuperscript{\ensuremath{x}}\fi}
\newunicodechar{ᵇ}{\ifmmode^{b}\else\textsuperscript{\ensuremath{b}}\fi}
\newunicodechar{ʳ}{\ifmmode^{r}\else\textsuperscript{\ensuremath{r}}\fi}
\newunicodechar{ᵘ}{\ifmmode^{u}\else\textsuperscript{\ensuremath{u}}\fi}
\newunicodechar{ʸ}{\ifmmode^{y}\else\textsuperscript{\ensuremath{y}}\fi}
\newunicodechar{ᵃ}{\ifmmode^{a}\else\textsuperscript{\ensuremath{a}}\fi}
\newunicodechar{ᵉ}{\ifmmode^{e}\else\textsuperscript{\ensuremath{e}}\fi}
\newunicodechar{ᵏ}{\ifmmode^{k}\else\textsuperscript{\ensuremath{k}}\fi}
\newunicodechar{ᵖ}{\ifmmode^{p}\else\textsuperscript{\ensuremath{p}}\fi}
\newunicodechar{ᴺ}{\ifmmode^{N}\else\textsuperscript{\ensuremath{N}}\fi}
\newunicodechar{ᶜ}{\ifmmode^{c}\else\textsuperscript{\ensuremath{c}}\fi}
\newunicodechar{ʰ}{\ifmmode^{h}\else\textsuperscript{\ensuremath{h}}\fi}
\newunicodechar{ᵠ}{\ifmmode^{q}\else\textsuperscript{\ensuremath{q}}\fi}
\newunicodechar{ₙ}{\ifmmode_{n}\else\textsubscript{\ensuremath{n}}\fi}
\newunicodechar{ₐ}{\ifmmode_{a}\else\textsubscript{\ensuremath{a}}\fi}
\newunicodechar{ᵢ}{\ifmmode_{i}\else\textsubscript{\ensuremath{i}}\fi}
\newunicodechar{ᵣ}{\ifmmode_{r}\else\textsubscript{\ensuremath{r}}\fi}
\newunicodechar{ₖ}{\ifmmode_{k}\else\textsubscript{\ensuremath{k}}\fi}
\newunicodechar{ᵦ}{\ifmmode_{\beta}\else\textsubscript{\ensuremath{\beta}}\fi}
\newunicodechar{ₓ}{\ifmmode_{x}\else\textsubscript{\ensuremath{x}}\fi}
\newfontfamily\popdisplay{ArchivoBlack-Regular.ttf}[Path=../../../fonts/]
\newfontfamily\poplogo{SpaceGrotesk-Bold.ttf}[Path=../../../fonts/]
\newfontfamily\popsemi{PlusJakartaSans-SemiBold.ttf}[Path=../../../fonts/]
\newfontfamily\popxbold{PlusJakartaSans-ExtraBold.ttf}[Path=../../../fonts/]
\newfontfamily\popxboldls{PlusJakartaSans-ExtraBold.ttf}[Path=../../../fonts/,LetterSpace=3.0]

\setlist[enumerate]{leftmargin=1.7em,itemsep=5pt,topsep=4pt}
\setlist[itemize]{leftmargin=1.7em,itemsep=4pt,topsep=4pt,label={\color{poporange}\textbullet}}
\setlength{\parindent}{0pt}
\setlength{\parskip}{5pt}
\renewcommand{\arraystretch}{1.25}

% pgfplots : style Pop par défaut
\pgfplotsset{
  every axis/.append style={axis line style={popink,line width=1pt},tick style={popink},
    grid style={popline,line width=0.5pt},label style={font=\small},tick label style={font=\footnotesize}},
  popcurve/.style={poporange,line width=1.8pt,no markers,smooth},
}
% Même style disponible dans un \draw TikZ simple (hors pgfplots)
\tikzset{popcurve/.style={poporange,line width=1.8pt}}
\tikzset{popgrid/.style={popline,very thin}}
\colorlet{popcurve}{poporange} % le nom du style est parfois utilisé comme couleur

% ============================================================
% Pill / squiggle
% ============================================================
\newcommand{\poppill}[3]{%
  \tikz[baseline=-0.55ex]\node[draw=popink,line width=1.2pt,rounded corners=2.6pt,%
    fill=#2,inner xsep=6.5pt,inner ysep=3pt]{%
    \popxboldls\fontsize{7}{7}\selectfont\color{#3}\MakeUppercase{#1}};%
}
\newcommand{\popsquiggle}[1]{%
  \tikz[baseline=-0.4ex]{
    \draw[#1,line width=2.6pt,line cap=round]
      plot [smooth] coordinates {(0,0.09) (0.34,0.01) (0.68,0.21) (1.02,0.01) (1.36,0.10)};
  }%
}
% Mot-clé surligné (marqueur orange sous le texte)
\newcommand{\cle}[1]{{\popxbold\color{popdark}#1}}

% ============================================================
% En-tête / pied
% ============================================================
\pagestyle{fancy}
\fancyhf{}
\renewcommand{\headrulewidth}{0pt}
\renewcommand{\footrulewidth}{0pt}
\lhead{\raisebox{-0.15ex}{\poplogo\fontsize{13}{13}\selectfont\color{popink}OnBuch\color{poporange}+}}
\rhead{\poppill{\DOCMATIERE\ \textbullet\ \DOCNIVEAU\ \textbullet\ Leçon \DOCLECON}{white}{popink}}
\lfoot{\popxbold\fontsize{7.6}{7.6}\selectfont\color{popmuted}\MakeUppercase{Cours signé OnBuch+}\ \textbullet\ {\popsemi\DOCTITRE}}
\rfoot{\tikz[baseline=-0.6ex]\node[rounded corners=8.5pt,fill=popink,minimum height=17pt,minimum width=17pt,inner xsep=5pt,inner sep=0]{\popxbold\fontsize{8}{8}\selectfont\color{popcream}\thepage\,/\,\pageref*{LastPage}};}

% ============================================================
% Titres
% ============================================================
\newcommand{\chaptitre}[2]{%
  \begin{tcolorbox}[enhanced,colback=poporange,colframe=popink,boxrule=2.6pt,arc=11pt,
      drop shadow=popink,left=15pt,right=15pt,top=12pt,bottom=11pt,before skip=3pt,after skip=18pt]
    \poppill{#2}{popink}{popcream}\par\vspace{9pt}
    {\raggedright\hyphenpenalty=10000\exhyphenpenalty=10000\popdisplay\fontsize{22}{24}\selectfont\color{popcream}\MakeUppercase{#1}\par}
  \end{tcolorbox}
}
% Section numérotée : pastille orange avec le numéro + titre Archivo
\newcounter{popsec}
\newcommand{\coursec}[1]{%
  \stepcounter{popsec}\setcounter{popsub}{0}%
  \par\vspace{16pt}\noindent
  \begin{minipage}{\linewidth}
  \tikz[baseline=-0.7ex]\node[draw=popink,line width=1.6pt,fill=poporange,rounded corners=4pt,minimum size=22pt,inner sep=0]{\popdisplay\fontsize{12}{12}\selectfont\color{popcream}\thepopsec};%
  \hspace{8pt}{\popdisplay\fontsize{15}{17}\selectfont\color{popink}\MakeUppercase{#1}}\par
  \vspace{4pt}\noindent{\color{poporange}\rule{36pt}{3.6pt}}
  \end{minipage}\par\nopagebreak\vspace{8pt}
}
\newcounter{popsub}
\newcommand{\courssub}[1]{%
  \stepcounter{popsub}%
  \par\vspace{9pt}\noindent{\popxbold\fontsize{11.5}{13}\selectfont\color{popink}{\color{poporange}\thepopsec.\thepopsub}\hspace{6pt}#1}\par\nopagebreak\vspace{4pt}
}

% ============================================================
% Boîtes pédagogiques — même recette : fond blanc, bordure épaisse colorée,
% ombre dure encre, badge plein. Titre optionnel [#1].
% ============================================================
\tcbset{popbase/.style={breakable,enhanced,colback=white,boxrule=2.3pt,arc=9pt,
  shadow={3pt}{-3pt}{0pt}{popink},left=14pt,right=14pt,top=11pt,bottom=11pt,before skip=11pt,after skip=15pt,
  fontupper=\popsemi\fontsize{10.6}{14.5}\selectfont\color{popink},
  fontlower=\popsemi\fontsize{10.6}{14.5}\selectfont\color{popink}}}
% Coche dessinée (le glyphe ✓ n'existe pas dans les polices du document)
\newcommand{\popcheck}[1]{\tikz[baseline=-0.5ex]\draw[#1,line width=1.6pt,line cap=round,line join=round](-0.13,0.0)--(-0.04,-0.09)--(0.13,0.1);}
\providecommand{\coche}{\popcheck{popgreen}}
\providecommand{\resultat}[1]{\par\smallskip\noindent\fcolorbox{popgreen}{popgreenL}{\parbox{\dimexpr\linewidth-2\fboxsep-2\fboxrule\relax}{\textbf{Résultat.} #1}}\par\smallskip}
\newcommand{\popboxhead}[3]{\poppill{#1}{#2}{white}\ifx\relax#3\relax\else\hspace{6pt}{\popxbold\fontsize{10.6}{12}\selectfont\color{popink}#3}\fi\par\vspace{7pt}}

\newtcolorbox{aretenir}[1][]{popbase,colframe=popgreen,before upper={\popboxhead{À retenir}{popgreen}{#1}}}
% Texte en anglais (document de lecture, dialogue, extrait) : cadre
% d'étude. Un titre optionnel (source, auteur) s'affiche dans l'en-tête.
\newtcolorbox{textebox}[1][]{popbase,colframe=popblue,before upper={\popboxhead{Text}{popblue}{#1}\begin{otherlanguage}{english}},after upper={\end{otherlanguage}}}
% Marques ✓ / ✗ (forme correcte / fautive) souvent utilisées par les modèles
\providecommand{\cross}{\textcolor{poppink}{\ensuremath{\times}}}
\providecommand{\xmark}{\cross}\providecommand{\crossmark}{\cross}\providecommand{\cmark}{\ensuremath{\checkmark}}
% Ligne de réponse / justification à compléter (inventée par les modèles)
\providecommand{\justification}[1]{\par\noindent\textit{Justification~:} \dotfill\par}
% \textipa (package tipa non chargé) : la police ne couvre pas l'API -> simple passage
\providecommand{\textipa}[1]{#1}
% Soulignement (ulem non chargé) : \uline, \ul
\providecommand{\uline}[1]{\underline{#1}}\providecommand{\ul}[1]{\underline{#1}}
% \glossaire{\item ...} inventée par les modèles : liste de glossaire
\providecommand{\glossaire}[1]{\begin{itemize}#1\end{itemize}}
% \alert (beamer) inventée par les modèles : texte mis en évidence
\providecommand{\alert}[1]{\textbf{\textcolor{poppink}{#1}}}
% variantes ASCII d'ordinaux écrites en macro
\providecommand{\iere}{\textsuperscript{re}}\providecommand{\ieme}{\textsuperscript{e}}\providecommand{\ere}{\textsuperscript{re}}\providecommand{\eme}{\textsuperscript{e}}\providecommand{\er}{\textsuperscript{er}}\providecommand{\ie}{\textsuperscript{e}}
% Ordinaux français écrits en macro par les modèles : 1\ière, 3\ième
\newcommand{\ière}{\textsuperscript{re}}\newcommand{\ième}{\textsuperscript{e}}
% \exemple (inventée par les modèles) : étiquette « Exemple : »
\providecommand{\exemple}{\textbf{Exemple~:}\ }
% \square utilisable aussi hors mode math (cases à cocher)
\AtBeginDocument{\let\squareMath\square\renewcommand{\square}{\ensuremath{\squareMath}}}
% Mot ou phrase en anglais à l'intérieur d'un paragraphe français
\newcommand{\eng}[1]{\ifmmode\text{\textit{#1}}\else\begingroup\selectlanguage{english}\itshape #1\endgroup\fi}
\let\esp\eng % alias toléré
\newtcolorbox{exemplebox}[1][]{popbase,colframe=poporange,before upper={\popboxhead{Exemple}{poporange}{#1}}}
\newtcolorbox{definition}[1][]{popbase,colframe=popblue,before upper={\popboxhead{Définition}{popblue}{#1}}}
\newtcolorbox{propriete}[1][]{popbase,colframe=poppurple,before upper={\popboxhead{Propriété}{poppurple}{#1}}}
\newtcolorbox{methode}[1][]{popbase,colframe=popgold,before upper={\popboxhead{Méthode}{popgold}{#1}}}
\newtcolorbox{attention}[1][]{popbase,colframe=poppink,before upper={\popboxhead{Attention}{poppink}{#1}}}
\newtcolorbox{experience}[1][]{popbase,colframe=popblue,colback=popblueL,before upper={\popboxhead{Expérience}{popblue}{#1}}}
% « Le savais-tu ? » : carte violette pleine (contexte, culture, Cameroun)
\newtcolorbox{savaistu}[1][]{popbase,colback=poppurple,colframe=popink,
  fontupper=\popsemi\fontsize{10.4}{14.2}\selectfont\color{white},
  before upper={\poppill{Le savais-tu\,?}{white}{poppurple}\ifx\relax#1\relax\else\hspace{6pt}{\popxbold\color{white}#1}\fi\par\vspace{7pt}}}
% Exercice résolu : énoncé en haut, \tcblower puis la solution sur fond crème
\newcounter{popexo}\newcounter{popexr}
\newtcolorbox{exoresolu}[1][]{popbase,colframe=poporange,colbacklower=popcream,
  segmentation style={popink,line width=1.2pt,dashed},
  before upper={\stepcounter{popexr}\popboxhead{Exercice résolu \thepopexr}{poporange}{#1}},
  before lower={\poppill{Solution}{popink}{popcream}\par\vspace{6pt}}}
% Exercice (non résolu en place) — la correction est dans la partie Corrigés
\newtcolorbox{exercice}[1][]{popbase,colframe=popink,
  before upper={\stepcounter{popexo}\popboxhead{Exercice \thepopexo}{popink}{#1}}}
% Corrigé
\newtcolorbox{corrige}[1][]{popbase,colframe=popgreen,colback=popgreenL,
  before upper={\popboxhead{Corrigé}{popgreen}{#1}}}
% Objectifs / prérequis / bilan
\newtcolorbox{objectifs}{popbase,colframe=popink,colback=poporangeL,
  before upper={\popboxhead{Objectifs de la leçon}{popink}{}}}
\newtcolorbox{prerequis}{popbase,colframe=popink,colback=popblueL,
  before upper={\popboxhead{Prérequis}{popblue}{}}}
\newtcolorbox{bilan}{popbase,colframe=popink,colback=white,boxrule=2.8pt,
  before upper={\popboxhead{Fiche bilan}{poporange}{L'essentiel à connaître pour l'examen}}}
% Figure encadrée avec légende
\newtcolorbox{popfigure}[1][]{enhanced,breakable=false,colback=white,colframe=popline,boxrule=1.4pt,arc=8pt,
  left=10pt,right=10pt,top=10pt,bottom=8pt,before skip=10pt,after skip=12pt,halign=center,
  lower separated=false,halign lower=center,
  fontlower=\popsemi\fontsize{9}{11.5}\selectfont\color{popmuted},
  #1}
\newcommand{\legende}[1]{\par\vspace{6pt}{\popsemi\fontsize{9}{11.5}\selectfont\color{popmuted}#1}}

% ============================================================
% Signature OnBuch+
% ============================================================
\newcommand{\poplogoinline}[1]{{\poplogo\fontsize{#1}{#1}\selectfont\color{popink}OnBuch\color{poporange}+}}
\newcommand{\signatureonbuch}{%
  \begin{tcolorbox}[enhanced,colback=popink,colframe=popink,boxrule=2.2pt,arc=10pt,
      drop shadow=poporange,left=14pt,right=14pt,top=10pt,bottom=10pt,before skip=0pt,after skip=0pt]
    \tikz[baseline=-0.7ex]\node[circle,fill=popgreen,draw=popcream,line width=1.3pt,minimum size=22pt,inner sep=0]{\popcheck{white}};%
    \hspace{9pt}\begin{minipage}[c]{0.62\linewidth}
      {\popxbold\fontsize{12}{13}\selectfont\color{popcream}Signé\ \textbullet\ {\poplogo OnBuch\color{poporange}+}}\par\vspace{2pt}
      {\raggedright\popsemi\fontsize{8}{10.5}\selectfont\color{popline}\MakeUppercase{Équipe pédagogique OnBuch+}\\\MakeUppercase{Conforme au programme officiel MINESEC}\par}
    \end{minipage}\hfill
    {\poplogo\fontsize{22}{22}\selectfont\color{popcream}OnBuch\color{poporange}+}
  \end{tcolorbox}%
}

% ============================================================
% Page de garde
% #1 titre · #2 matière · #3 niveau · #4 module · #5 n° leçon · #6 durée ·
% #7 description / objectifs (texte ou itemize)
% ============================================================
\newcommand{\pagegarde}[7]{%
  \thispagestyle{empty}
  \newgeometry{a4paper,margin=1.7cm,top=1.6cm,bottom=1.4cm}%
  \begin{tikzpicture}[remember picture,overlay]
    \fill[poporange!18] ([xshift=-3.2cm,yshift=-3.6cm]current page.north east) circle (4.6cm);
    \fill[poppurple!14] ([xshift=2.4cm,yshift=7.6cm]current page.south west) circle (3.8cm);
  \end{tikzpicture}%
  {\poplogo\fontsize{30}{30}\selectfont\color{popink}OnBuch\color{poporange}+}\hfill
  \raisebox{0.6ex}{\poppill{Cours signé \textbullet\ Premium}{popink}{popcream}}\par
  \vspace{1pt}\popsquiggle{poppurple}\par
  \vspace{8pt}
  \begin{tcolorbox}[enhanced,colback=poporange,colframe=popink,boxrule=2.8pt,arc=13pt,
      drop shadow=popink,left=18pt,right=18pt,top=12pt,bottom=13pt,after skip=10pt]
    \poppill{#2 \textbullet\ #3}{popink}{popcream}\hspace{5pt}\poppill{#4}{white}{popink}\par\vspace{8pt}
    {\popdisplay\fontsize{14}{14}\selectfont\color{popink}LEÇON #5}\par\vspace{4pt}
    {\raggedright\hyphenpenalty=10000\exhyphenpenalty=10000\popdisplay\fontsize{25}{27}\selectfont\color{popcream}\MakeUppercase{#1}\par}
    \vspace{8pt}
    {\popxbold\fontsize{9.5}{9.5}\selectfont\color{popink}\MakeUppercase{Durée conseillée : #6}}
  \end{tcolorbox}
  \begin{tcolorbox}[enhanced,colback=white,colframe=popink,boxrule=2.2pt,arc=10pt,
      drop shadow=popink,left=16pt,right=16pt,top=10pt,bottom=10pt,after skip=8pt]
    \poppill{Dans cette leçon}{poppurple}{white}\par\vspace{5pt}
    \popsemi\fontsize{9.3}{12.6}\selectfont\color{popink}#7
  \end{tcolorbox}
  \noindent\begin{minipage}[t]{0.487\linewidth}
    \begin{tcolorbox}[enhanced,colback=popgreen,colframe=popink,boxrule=2.2pt,arc=10pt,
        drop shadow=popink,left=11pt,right=11pt,top=10pt,bottom=10pt,height=3.1cm,valign=center]
      \tcbox[colback=white,colframe=popink,boxrule=1.3pt,arc=5pt,boxsep=0pt,nobeforeafter,
        left=3pt,right=3pt,top=3pt,bottom=3pt,valign=center]{\qrcode[height=1.9cm]{https://play.google.com/store/apps/details?id=com.luvvix.onbuch&hl=fr}}%
      \hspace{8pt}\begin{minipage}[c]{0.5\linewidth}
        {\popdisplay\fontsize{11}{12.5}\selectfont\color{white}\MakeUppercase{Télécharge OnBuch}}\par\vspace{4pt}
        {\raggedright\popsemi\fontsize{8.4}{11}\selectfont\color{white}Gratuit \textbullet\ Google Play\\Tuteur IA, annales, concours, résultats\par}
      \end{minipage}
    \end{tcolorbox}
  \end{minipage}\hfill
  \begin{minipage}[t]{0.487\linewidth}
    \begin{tcolorbox}[enhanced,colback=poppurple,colframe=popink,boxrule=2.2pt,arc=10pt,
        drop shadow=popink,left=11pt,right=11pt,top=10pt,bottom=10pt,height=3.1cm,valign=center]
      \tikz[baseline=-0.7ex]\node[circle,fill=white,draw=popink,line width=1.3pt,minimum size=1.6cm,inner sep=0]{\popdisplay\fontsize{24}{24}\selectfont\color{poppurple}+};%
      \hspace{8pt}\begin{minipage}[c]{0.6\linewidth}
        {\popdisplay\fontsize{11}{12.5}\selectfont\color{white}\MakeUppercase{Passe à OnBuch+}}\par\vspace{4pt}
        {\raggedright\popsemi\fontsize{8.4}{11}\selectfont\color{white}Tous les cours signés, Tuteur IA illimité, fiches PDF — onglet OnBuch+ de l'app\par}
      \end{minipage}
    \end{tcolorbox}
  \end{minipage}
  \vspace{8pt}
  \popcontenu
  \vfill
  \signatureonbuch
  \restoregeometry
  \newpage
}

% Rangée « Ce cours contient » (page de garde)
\newcommand{\poptile}[3]{% #1 couleur #2 gros texte #3 légende
  \begin{tcolorbox}[enhanced,colback=#1,colframe=popink,boxrule=2pt,arc=9pt,shadow={2.5pt}{-2.5pt}{0pt}{popink},
      left=6pt,right=6pt,top=9pt,bottom=9pt,halign=center,valign=center,height=1.75cm,width=\linewidth,nobeforeafter]
    {\popdisplay\fontsize{12.5}{13}\selectfont\color{white}\MakeUppercase{#2}}\par\vspace{3pt}
    {\centering\popsemi\fontsize{7.8}{9.5}\selectfont\color{white}#3\par}
  \end{tcolorbox}}
\newcommand{\popcontenu}{%
  \par\noindent{\popxboldls\fontsize{7.5}{7.5}\selectfont\color{popmuted}CE COURS SIGNÉ CONTIENT}\par\vspace{7pt}
  \noindent\begin{minipage}[t]{0.235\linewidth}\poptile{popblue}{Cours}{complet, illustré, pas à pas}\end{minipage}\hfill
  \begin{minipage}[t]{0.235\linewidth}\poptile{poporange}{Méthodes}{et exercices résolus}\end{minipage}\hfill
  \begin{minipage}[t]{0.235\linewidth}\poptile{poppink}{Exercices}{type examen + corrigés}\end{minipage}\hfill
  \begin{minipage}[t]{0.235\linewidth}\poptile{popgreen}{Fiche bilan}{l'essentiel à retenir}\end{minipage}\par}

% Sommaire
\newtcolorbox{sommaire}{enhanced,colback=white,colframe=popink,boxrule=2.2pt,arc=10pt,
  drop shadow=popink,left=18pt,right=18pt,top=13pt,bottom=13pt,before skip=6pt,after skip=20pt,
  before upper={\poppill{Sommaire}{poppurple}{white}\par\vspace{9pt}\popsemi\fontsize{10.6}{14.5}\selectfont\color{popink}}}

% Signature de fin de cours — poussée tout en bas de la dernière page
\newcommand{\findecours}{%
  \par\vfill\nopagebreak
  \begin{center}\popsquiggle{poporange}\end{center}\vspace{4pt}
  \signatureonbuch
}
\AtBeginDocument{\def\llbracket{\mathopen{[\mkern-3.5mu[}}\def\rrbracket{\mathclose{]\mkern-3.5mu]}}} % intervalles d'entiers (glyphe ⟦ absent de la police)
% Glyphes absents de Fira Math : redéfinis à partir de glyphes disponibles
\AtBeginDocument{%
  \def\setminus{\mathbin{\backslash}}\let\smallsetminus\setminus
  \def\vdots{\mathord{\vbox{\baselineskip3.2pt\lineskiplimit0pt\kern4pt\hbox{.}\hbox{.}\hbox{.}}}}%
  \def\ddots{\mathinner{\mkern1mu\raise6.5pt\hbox{.}\mkern2mu\raise3.6pt\hbox{.}\mkern2mu\raise0.7pt\hbox{.}\mkern1mu}}%
  \def\triangle{\mathord{\scalebox{0.9}{$\symup{\Delta}$}}}%
  \def\nabla{\mathord{\raisebox{\depth}{\scalebox{1}[-1]{$\symup{\Delta}$}}}}%
  \def\checkmark{\popcheck{popdark}}%
  \def\star{\mathbin{\ast}}\def\bigstar{\mathord{\ast}}%
  \def\diamond{\mathbin{\scalebox{0.6}{$\Diamond$}}}%
  \def\bigcup{\mathop{\mathchoice{\raisebox{-0.3em}{\scalebox{1.7}{$\cup$}}}{\scalebox{1.25}{$\cup$}}{\cup}{\cup}}\displaylimits}%
  \def\bigcap{\mathop{\mathchoice{\raisebox{-0.3em}{\scalebox{1.7}{$\cap$}}}{\scalebox{1.25}{$\cap$}}{\cap}{\cap}}\displaylimits}%
  \def\lesseqqgtr{\mathrel{\lessgtr}}\def\bigoplus{\mathop{\oplus}\displaylimits}%
}
\providecommand{\ssi}{si et seulement si} % abréviation fréquente
\AtBeginDocument{\providecommand{\wideparen}{\overparen}} % arc de cercle

%% FIGFIX-BEGIN v3 — correctifs globaux des figures (docs/onbuch-generation/tools/figfix)
\usepackage{adjustbox}\usepackage{etoolbox}
% toute figure TikZ/pgfplots plus large que la ligne est réduite à la largeur disponible
\BeforeBeginEnvironment{tikzpicture}{\begin{adjustbox}{max width=\linewidth}}
\AfterEndEnvironment{tikzpicture}{\end{adjustbox}}
% légendes pgfplots lisibles par-dessus les courbes
\pgfplotsset{every axis/.append style={legend style={fill=white,fill opacity=0.92,text opacity=1,draw=black!15,inner sep=3pt}}}
% étiquettes d'arêtes (syntaxe edge["..."]) avec fond blanc
\tikzset{every edge quotes/.append style={fill=white,inner sep=1.2pt}}
% bulles de cartes mentales : texte à la ligne dans la bulle, bulle un peu plus grande
\tikzset{every concept/.append style={text width=2.0cm,align=center,minimum size=2.3cm,inner sep=2pt}}
%% FIGFIX-END
\begin{document}

\pagegarde{Lesson 12 — Writing: Writes compositions about anticipating one's second language needs at work}{Anglais}{Tle A}{Chapter 1 : Family and Social Life}{EN12}{2 h}{Cette leçon d'expression écrite apprend aux élèves de Tle A à planifier, rédiger et corriger une composition sur l'anticipation de leurs besoins en anglais au travail. Elle fournit la structure, les connecteurs, le vocabulaire, un modèle commenté et une grille d'évaluation conformes aux attentes du Bac.
\vspace{4pt}
\begin{multicols}{2}\small
\begin{itemize}
\item À la fin de cette leçon, je sais analyser un sujet de composition et dégager les mots-clés
\item À la fin de cette leçon, je sais construire un plan en trois parties : introduction, développement, conclusion
\item À la fin de cette leçon, je sais rédiger une introduction avec accroche, présentation du sujet et annonce de plan
\item À la fin de cette leçon, je sais développer des arguments avec des exemples liés au monde du travail
\item À la fin de cette leçon, je sais utiliser les connecteurs logiques appropriés (addition, cause, conséquence, opposition)
\item À la fin de cette leçon, je sais employer les temps verbaux adaptés (futur, présent, conditionnel)
\end{itemize}\end{multicols}}

\begin{sommaire}
\begin{multicols}{2}
\begin{enumerate}
\item Analyser le sujet et planifier
\item La structure d'une composition
\item Connecteurs, temps et vocabulaire utiles
\item Modèle commenté
\item Erreurs à éviter et relecture
\item Grille d'évaluation et entraînement au Bac
\item Activité d'intégration
\item Méthodes et exercices résolus
\item Exercices
\item Corrigés des exercices
\item Fiche bilan
\end{enumerate}
\end{multicols}
\end{sommaire}

\begin{objectifs}
\begin{itemize}
\item À la fin de cette leçon, je sais analyser un sujet de composition et dégager les mots-clés
\item À la fin de cette leçon, je sais construire un plan en trois parties : introduction, développement, conclusion
\item À la fin de cette leçon, je sais rédiger une introduction avec accroche, présentation du sujet et annonce de plan
\item À la fin de cette leçon, je sais développer des arguments avec des exemples liés au monde du travail
\item À la fin de cette leçon, je sais utiliser les connecteurs logiques appropriés (addition, cause, conséquence, opposition)
\item À la fin de cette leçon, je sais employer les temps verbaux adaptés (futur, présent, conditionnel)
\item À la fin de cette leçon, je sais mobiliser le vocabulaire du travail et des compétences linguistiques
\item À la fin de cette leçon, je sais relire et corriger ma production à l'aide d'une grille d'évaluation
\end{itemize}
\end{objectifs}

\begin{prerequis}
\begin{itemize}
\item Les temps verbaux de base : present simple, future (will / be going to), conditionnel (would)
\item Le vocabulaire des métiers et de la vie professionnelle (vu en Première)
\item La rédaction d'un paragraphe argumentatif simple
\item Les connecteurs de base : and, but, because, so, first, then, finally
\end{itemize}
\end{prerequis}

\newpage

\coursec{Analyser le sujet et planifier}

Imaginez : vous venez d'être embauché par une ONG à Buea. Dès votre premier jour, votre superviseur vous demande de rédiger des courriels en anglais, d'assister à des réunions avec des partenaires du Nigeria et du Ghana, et de rédiger des rapports. Vos compétences en anglais sont-elles prêtes pour le monde du travail ? Dans cette leçon, vous apprendrez à écrire une composition dans laquelle vous anticipez les compétences en anglais dont vous aurez besoin dans votre future carrière — un sujet fréquent au Baccalauréat.

\courssub{Comprendre le sujet : mots-clés et consignes}

Avant toute rédaction, la première étape — et la plus décisive — est l'analyse rigoureuse du sujet. Au Baccalauréat, un hors-sujet ou une analyse superficielle coûte très cher, même si votre anglais est correct. Prenons le sujet type :

\begin{textebox}{Sujet d'examen type}
Write a composition about how you anticipate your second language needs at work.
\end{textebox}

\begin{methode}{Analyser un sujet de composition (3 étapes)}
\begin{enumerate}
    \item \textbf{Soulignez les mots-clés} (mots porteurs de sens, verbes d'action, thèmes).
    \item \textbf{Posez les questions W :} \eng{Who?} (qui ?), \eng{What?} (quoi ?), \eng{Why?} (pourquoi ?), \eng{How?} (comment ?), \eng{When/Where?} (quand/où ?).
    \item \textbf{Reformulez le sujet en français} pour vérifier votre compréhension : « Rédigez une composition sur la manière dont vous anticipez vos besoins en langue seconde au travail. »
\end{enumerate}
\end{methode}

\noindent
\textbf{Repérage des mots-clés :}

\begin{center}
\begin{tabularx}{\linewidth}{|X|X|X|}
\hline
\rowcolor{popblueL}
\textbf{Mot-clé (English)} & \textbf{Nature} & \textbf{Sens / Traduction} \\
\hline
\cle{anticipate} & Verbe & Prévoir, anticiper, se préparer à l'avance (pas seulement « deviner ») \\
\hline
\cle{second language needs} & Groupe nominal & Besoins en langue seconde (compétences : lire, écrire, parler, écouter) \\
\hline
\cle{at work} & Locution prépositive & Dans le cadre professionnel, sur le lieu de travail \\
\hline
\cle{composition} & Nom & Rédaction structurée (introduction, développement, conclusion) — pas une lettre, pas un dialogue \\
\hline
\end{tabularx}
\end{center}

\begin{attention}{Piège fréquent : « Anticipate » vs « Expect »}
Ne confondez pas \eng{anticipate} (prévoir \textbf{et agir en conséquence}, se préparer) et \eng{expect} (s'attendre à, penser que quelque chose va arriver).
\begin{itemize}
    \item \eng{I anticipate needing English for reports.} = Je prévois d'avoir besoin de l'anglais pour les rapports (donc je me forme).
    \item \eng{I expect to use English.} = Je m'attends à utiliser l'anglais (constat passif).
\end{itemize}
Au Bac, le sujet demande une démarche \textbf{active} (stratégies, formation, préparation), pas une simple liste de souhaits.
\end{attention}

\courssub{Le brainstorming : générer et classer ses idées}

Une fois le sujet compris, ne commencez \textbf{jamais} à rédiger directement. La technique du \textbf{brainstorming} (remue-méninges) consiste à noter \textbf{toutes} les idées qui vous viennent, sans les juger, sans les trier, pendant 5 à 10 minutes. Utilisez une \textbf{carte mentale} (mind map) pour visualiser les liens.

\begin{popfigure}
\begin{tikzpicture}[
    mindmap,
    concept color=popblue!80,
    text=white,
    font=\small\bfseries,
    level 1/.style={level distance=3.5cm, sibling angle=90},
    level 2/.style={level distance=2.5cm, sibling angle=45},
    every node/.style={concept, execute at begin node=\hskip0pt},
    popblue/.style={concept color=popblue},
    poporange/.style={concept color=poporange},
    popgreen/.style={concept color=popgreen},
    poppurple/.style={concept color=poppurple},
    poppink/.style={concept color=poppink},
]
\node[concept, scale=1.2] {My English needs at work}
    child[poporange] {node[concept] {Job / Career}
        child {node[concept, scale=0.7] {Teacher}}
        child {node[concept, scale=0.7] {Nurse}}
        child {node[concept, scale=0.7] {Engineer}}
        child {node[concept, scale=0.7] {Business}}
    }
    child[popgreen] {node[concept] {Skills Needed}
        child {node[concept, scale=0.7] {Writing emails}}
        child {node[concept, scale=0.7] {Speaking in meetings}}
        child {node[concept, scale=0.7] {Reading reports}}
        child {node[concept, scale=0.7] {Listening to clients}}
    }
    child[poppurple] {node[concept] {Strategies}
        child {node[concept, scale=0.7] {English courses}}
        child {node[concept, scale=0.7] {Online practice}}
        child {node[concept, scale=0.7] {Media (films, news)}}
        child {node[concept, scale=0.7] {Language exchange}}
    }
    child[poppink] {node[concept] {Benefits}
        child {node[concept, scale=0.7] {Promotion}}
        child {node[concept, scale=0.7] {International mobility}}
        child {node[concept, scale=0.7] {Better salary}}
        child {node[concept, scale=0.7] {Confidence}}
    };
\end{tikzpicture}
\legende{Carte mentale : organisation des idées autour du thème central}
\end{popfigure}

\begin{exemplebox}{Exemples de brainstorming selon le métier visé}
\begin{itemize}
    \item \textbf{Future Teacher (Enseignant) :} \eng{Need English to teach bilingual classes, read pedagogical resources, attend international workshops.} « J'ai besoin d'anglais pour enseigner des classes bilingues, lire des ressources pédagogiques, assister à des ateliers internationaux. » \eng{Strategy: Take a TEFL certification, practice speaking with colleagues.} « Stratégie : passer une certification TEFL, pratiquer l'oral avec des collègues. »
    \item \textbf{Future Nurse (Infirmier) :} \eng{Need English to understand medical prescriptions, communicate with expatriate patients, read WHO guidelines.} « J'ai besoin d'anglais pour comprendre les prescriptions médicales, communiquer avec des patients expatriés, lire les directives de l'OMS. » \eng{Strategy: Medical English MOOC, vocabulary flashcards.} « Stratégie : MOOC d'anglais médical, fiches de vocabulaire. »
    \item \textbf{Future Business Manager (Gestionnaire) :} \eng{Need English for negotiations, contracts, financial reports, emails to partners in Lagos or Dubai.} « J'ai besoin d'anglais pour les négociations, contrats, rapports financiers, courriels aux partenaires à Lagos ou Dubaï. » \eng{Strategy: Business English course, read The Economist, practice presentations.} « Stratégie : cours d'anglais des affaires, lire The Economist, s'entraîner aux présentations. »
\end{itemize}
\end{exemplebox}

\courssub{Choisir un angle personnel et réaliste}

Le correcteur attend une composition \textbf{personnelle} (\eng{first person}: \eng{I, my, me}) et \textbf{authentique}. N'inventez pas un métier qui ne vous intéresse pas. Choisissez \textbf{votre} projet professionnel réel (ou celui que vous préparez sérieusement).

\begin{propriete}{Choix de l'angle personnel}
\begin{itemize}
    \item Utilisez le \textbf{futur simple} (\eng{will need, will have to}) pour les prédictions certaines.
    \item Utilisez les \textbf{modaux de probabilité} (\eng{might need, could be necessary, may have to}) pour les éventualités.
    \item Ancrez votre propos dans le \textbf{contexte camerounais} (Buea, Douala, Yaoundé, Bamenda) ou international si pertinent.
\end{itemize}
\end{propriete}

\begin{center}
\begin{tabularx}{\linewidth}{|X|X|}
\hline
\rowcolor{popblueL}
\textbf{Angle faible (générique)} & \textbf{Angle fort (personnalisé)} \\
\hline
\eng{People need English at work.} & \eng{As a future civil engineer in Douala, I will need English to read international standards.} \\
\hline
\eng{English is important for everyone.} & \eng{My goal is to work for the African Development Bank; therefore, mastering formal writing is crucial.} \\
\hline
\end{tabularx}
\end{center}

\courssub{Élaborer le plan en trois parties}

Une composition au Bac suit impérativement la structure ternaire : \textbf{Introduction} — \textbf{Développement} (2 ou 3 paragraphes) — \textbf{Conclusion}. Chaque paragraphe du développement = \textbf{une idée principale} + \textbf{arguments/exemples}.

\begin{methode}{Construire son plan détaillé (Squelette de la rédaction)}
\begin{enumerate}
    \item \textbf{Introduction} (1 paragraphe) :
    \begin{itemize}
        \item \textit{Hook / Accroche :} Contexte général (mondialisation, bilinguisme au Cameroun).
        \item \textit{Thesis statement / Annonce du plan :} \eng{In this composition, I will explain the English skills I will need as a [métier] and the strategies I plan to adopt.}
    \end{itemize}
    \item \textbf{Développement} (2 à 3 paragraphes) — \textit{Choisissez 2 ou 3 axes parmi :}
    \begin{itemize}
        \item \textbf{Axe 1 : Les compétences linguistiques spécifiques} (Speaking, Writing, Reading, Listening) liées au métier.
        \item \textbf{Axe 2 : Les situations professionnelles concrètes} (réunions, mails, rapports, formations).
        \item \textbf{Axe 3 : Les stratégies d'anticipation / préparation} (cours, auto-formation, certifications, immersion).
    \end{itemize}
    \item \textbf{Conclusion} (1 paragraphe) :
    \begin{itemize}
        \item Résumé des points forts.
        \item Ouverture / engagement personnel : \eng{Mastering English is not an option but a necessity for my career.}
    \end{itemize}
\end{enumerate}
\end{methode}

\begin{exemplebox}{Plan type pour le sujet}
\textbf{Sujet :} \eng{Write a composition about how you anticipate your second language needs at work.}

\textbf{Profil :} Élève de Terminale A, futur infirmier (nurse) à Bamenda.

\begin{enumerate}
    \item \textbf{Introduction :} Le Cameroun est bilingue ; le secteur de la santé accueille des patients anglophones et francophones. En tant que futur infirmier, l'anglais est un outil vital. \eng{Thesis: As a future nurse, I anticipate needing strong medical English for patient care and professional development.}
    \item \textbf{Paragraphe 1 (Compétences) :} \eng{Speaking} et \eng{Listening} pour l'anamnèse (interrogatoire) de patients anglophones ; \eng{Reading} pour les protocoles OMS et notices de médicaments ; \eng{Writing} pour les rapports de garde.
    \item \textbf{Paragraphe 2 (Stratégies) :} Cours d'anglais médical à l'université ; application \eng{Anki} pour le vocabulaire ; stage en zone anglophone (Buea) ; visionnage de séries médicales (\eng{The Good Doctor}) en VO.
    \item \textbf{Conclusion :} L'anticipation aujourd'hui évite l'urgence demain. Mon objectif : valider le \eng{OET} (Occupational English Test) avant l'embauche.
\end{enumerate}
\end{exemplebox}

\courssub{Organigramme du processus de planification}

\begin{popfigure}
\begin{tikzpicture}[
    node distance=1.2cm and 2cm,
    box/.style={draw=popblue, thick, rounded corners, fill=popblue!10, text width=2.8cm, align=center, minimum height=1.2cm, font=\small},
    arrow/.style={-Stealth, thick, popblue}
]
\node[box] (sujet) {Sujet : \eng{Anticipate L2 needs at work}};
\node[box, below=of sujet] (mots) {Repérage des mots-clés : \eng{anticipate, needs, work}};
\node[box, below=of mots] (brain) {Brainstorming / Carte mentale : Job, Skills, Strategies};
\node[box, below=of brain] (plan) {Choix de l'angle + Plan en 3 parties};
\node[box, below=of plan] (redac) {Rédaction guidée par le plan};

\draw[arrow] (sujet) -- (mots);
\draw[arrow] (mots) -- (brain);
\draw[arrow] (brain) -- (plan);
\draw[arrow] (plan) -- (redac);
\end{tikzpicture}
\legende{Organigramme : du sujet à la rédaction via la planification}
\end{popfigure}

\begin{savaistu}[Le bilinguisme professionnel au Cameroun]
Au Cameroun, la maîtrise de l'anglais (L2 pour les francophones) est un atout majeur, voire une exigence, dans de nombreux secteurs :
\begin{itemize}
    \item \textbf{Fonction publique :} Les concours de l'ENAM, de la Police, de la Douane comportent des épreuves d'anglais ; l'avancement exige souvent un niveau B2.
    \item \textbf{Banques et Assurances :} (Afriland First Bank, UBA, Société Générale) — relation clientèle internationale, SWIFT, conformité.
    \item \textbf{ONG et Organisations internationales :} (UNHCR, OMS, Plan International) — langue de travail quotidienne, rapports aux bailleurs de fonds.
    \item \textbf{Télécommunications / Tech :} (MTN, Orange, start-ups de la Silicon Mountain à Buea) — documentation technique, support client, veille en anglais.
    \item \textbf{Secteur pétrolier et portuaire :} (Kribi, Limbe, Douala) — normes HSE, contrats, expatriés.
\end{itemize}
Anticiper ses besoins, c'est donc sécuriser son employabilité dès la sortie du lycée ou de l'université.

\end{savaistu}

\courssub{Exercice résolu : du sujet au plan}

\begin{exoresolu}{Analyser et planifier}
\textbf{Énoncé :} Vous êtes élève en Terminale A, option Sciences de l'Éducation. Vous voulez devenir professeur d'anglais au secondaire. Le sujet est : \eng{``Write a composition about how you anticipate your second language needs at work.''} \\

\textbf{Solution détaillée :}

\textbf{1. Mots-clés soulignés :} \eng{anticipate} (prévoir/préparer), \eng{second language needs} (besoins en L2 = anglais), \eng{at work} (contexte scolaire).

\textbf{2. Questions W :}
\begin{itemize}
    \item \eng{Who?} Me, future English teacher.
    \item \eng{What?} Needs: advanced methodology vocabulary, classroom management language, assessment writing, research reading.
    \item \eng{Why?} Teach effectively, pass CAP/CAPIEMP, attend international training.
    \item \eng{How?} Master's in ELT, British Council webinars, practice teaching in English, reading \eng{ELT Journal}.
\end{itemize}

\textbf{3. Plan détaillé :}

\begin{itemize}
    \item \textbf{Introduction :} Hook sur l'importance de l'anglais comme langue d'enseignement au Cameroun (écoles bilingues). Thesis : \eng{As a future English teacher, I anticipate needing high-level proficiency in pedagogical English, academic writing, and intercultural communication.}
    \item \textbf{Development 1 (Classroom skills) :} \eng{Speaking} : instructions claires, metalanguage (grammar terms), pronunciation model. \eng{Listening} : understanding learners' errors, native speaker trainers.
    \item \textbf{Development 2 (Administrative \& Academic skills) :} \eng{Writing} : lesson plans, exam papers (Probatoire, Bac), official reports, correspondence with hierarchy. \eng{Reading} : new curriculum documents, research articles for professional development.
    \item \textbf{Development 3 (Strategies) :} Enroll in Master ELT at HTTC/ENS ; \eng{TEFL} certification online ; join \eng{Cameroon English Language Teachers Association (CAMELTA)} ; daily reading of \eng{BBC Learning English} / \eng{British Council} resources.
    \item \textbf{Conclusion :} Teaching English \eng{in} English requires constant upgrading. My anticipation strategy is a lifelong learning commitment. \eng{``The best teachers are eternal students.''}
\end{itemize}
\end{exoresolu}

\begin{exercice}{À vous de jouer : Planification}
Choisissez \textbf{un} des profils ci-dessous et produisez : (1) la liste des 5 mots-clés du sujet, (2) les réponses aux 5 questions W, (3) un plan détaillé en 3 parties (Introduction, 2 ou 3 paragraphes de développement, Conclusion) avec l'idée principale de chaque paragraphe.
\begin{enumerate}
    \item Profil A : Futur développeur web (Web developer) à Yaoundé, clients internationaux.
    \item Profil B : Futur agent de douane (Customs officer) à Kribi, zone portuaire anglophone/francophone.
    \item Profil C : Futur entrepreneur agricole (Agribusiness entrepreneur) à Dschang, export cacao/café.
\end{enumerate}
\end{exercice}

\begin{corrige}{Exercice : Planification}
\textbf{Profil A (Web Developer) — Éléments de correction :}
\begin{itemize}
    \item \textbf{Mots-clés :} Anticipate, second language, needs, work, composition.
    \item \textbf{Questions W :} Who: Me, junior dev. What: Reading docs (Stack Overflow, MDN), writing clean comments/commit messages, speaking in stand-up meetings with remote teams. Why: 80\% of tech docs in English; remote work for US/EU startups. How: English for Developers course (Coursera), coding in English, tech podcasts (\eng{Syntax}, \eng{ShopTalk Show}).
    \item \textbf{Plan :} Intro (Tech = English). Dev 1: Technical reading/writing (documentation, code comments). Dev 2: Oral communication (remote meetings, interviews). Dev 3: Strategies (specialized MOOCs, open source contribution). Conclusion: English as core hard skill.
\end{itemize}
\textbf{Profil B (Customs Officer) — Éléments de correction :}
\begin{itemize}
    \item \textbf{Mots-clés :} Identiques.
    \item \textbf{Questions W :} Who: Customs officer. What: Reading manifests/Bill of Lading, speaking with Nigerian/Ghanaian traders, writing seizure reports. Why: CEMAC/ECOWAS cross-border trade, legal liability. How: Customs English training (WCO), bilingual glossary, practice at Kribi port.
    \item \textbf{Plan :} Intro (Kribi deep-sea port = hub). Dev 1: Documentary English (HS codes, Incoterms). Dev 2: Oral interaction (traders, vessel crews). Dev 3: Preparation (WCO courses, legal English). Conclusion: Facilitating trade requires linguistic precision.
\end{itemize}
\textbf{Profil C (Agribusiness) — Éléments de correction :}
\begin{itemize}
    \item \textbf{Mots-clés :} Identiques.
    \item \textbf{Questions W :} Who: Agripreneur. What: Negotiating contracts (FOB/CIF), reading phytosanitary certificates, emailing buyers in Europe/US. Why: Export revenue depends on compliance and trust. How: Business English (British Council), Incoterms 2020 in English, trade fair simulation (SIAO).
    \item \textbf{Plan :} Intro (Cocoa/coffee = forex). Dev 1: Contractual \& regulatory writing. Dev 2: Negotiation \& relationship speaking. Dev 3: Strategies (certifications, trade missions). Conclusion: English opens premium markets.
\end{itemize}
\end{corrige}

\begin{attention}{Erreurs fatales à éviter dès la planification}
\begin{itemize}
    \item \textbf{Pas de plan = composition décousue.} Même avec un bon niveau, l'absence de structure (intro/dev/conclusion) fait perdre 4 à 6 points sur 20 au Bac.
    \item \textbf{Sujet traité au présent uniquement.} Le sujet demande d'\textbf{anticiper} (futur, modaux). Écrire \eng{I need English} (présent) au lieu de \eng{I will need / I anticipate needing} (projection) = hors-sujet partiel.
    \item \textbf{Anglais général vs Anglais professionnel.} Ne listez pas « \eng{I will learn vocabulary, grammar, verbs} ». Soyez spécifique : « \eng{I will master medical terminology / legal drafting / financial reporting} ».
    \item \textbf{Oublier le « How » (stratégies).} Anticiper, c'est \textbf{agir maintenant} pour le futur. Une composition qui ne cite aucune action concrète (cours, appli, stage, certification) est incomplète.
\end{itemize}
\end{attention}

\coursec{La structure d'une composition}

Une composition réussie respecte une architecture rigoureuse. Chaque partie a une fonction précise et des marqueurs linguistiques attendus.

\courssub{L'introduction : capter l'attention et annoncer le plan}

L'introduction (environ 10\% du texte) suit le modèle \textbf{entonnoir} : du général au particulier.

\begin{propriete}{Les 3 composantes de l'introduction}
\begin{enumerate}
    \item \textbf{Hook (Accroche) :} Une phrase d'ouverture large sur le thème (mondialisation, bilinguisme, évolution du marché du travail).
    \item \textbf{Contextualisation :} Lien avec votre situation personnelle (pays, ville, filière, métier visé).
    \item \textbf{Thesis statement (Phrase thèse) :} \textbf{La phrase la plus importante.} Elle annonce \textbf{ce que vous allez démontrer} et \textbf{comment} (les 2-3 axes du développement).
\end{enumerate}
\end{propriete}

\begin{exemplebox}{Introductions : comparer}
\begin{itemize}
    \item \textbf{Faible :} \eng{English is very important today. I want to be a nurse. I will need English.} (Trop simple, pas de thèse, pas de plan).
    \item \textbf{Forte :} \eng{In today's globalised world, bilingualism is a key asset in Cameroon's professional landscape. As a final-year student in Bamenda aiming to become a registered nurse, I am fully aware that English is not merely a school subject but a vital tool for patient safety. In this composition, I will outline the specific medical English skills I will need — particularly listening and reading — and the concrete strategies I am already implementing to acquire them.}
\end{itemize}
\end{exemplebox}

\courssub{Le développement : un paragraphe = une idée principale}

Chaque paragraphe de développement (environ 30-35\% chacun) suit la structure \textbf{PEEL} :
\begin{itemize}
    \item \textbf{P}oint : Phrase sujet (topic sentence) qui annonce l'idée principale.
    \item \textbf{E}vidence : Arguments, exemples concrets, détails professionnels.
    \item \textbf{E}xplanation : Analyse, lien avec l'anticipation (pourquoi cette compétence ? comment l'acquérir ?).
    \item \textbf{L}ink : Transition vers le paragraphe suivant ou retour à la thèse.
\end{itemize}

\begin{propriete}{Exigences du développement}
\begin{itemize}
    \item \textbf{Cohérence interne :} Utilisez des connecteurs logiques (voir section 3).
    \item \textbf{Spécificité :} « \eng{I will need to write emails} » → trop vague. « \eng{I will need to write formal incident reports and referral letters using medical terminology} » → précis.
    \item \textbf{Équilibre :} Les 2-3 paragraphes doivent avoir une longueur comparable.
\end{itemize}
\end{propriete}

\courssub{La conclusion : synthétiser et ouvrir}

La conclusion (environ 10-15\%) ne doit \textbf{pas} introduire de nouvelle idée.

\begin{propriete}{Les 3 mouvements de la conclusion}
\begin{enumerate}
    \item \textbf{Rappel reformulé de la thèse} (pas de copier-coller).
    \item \textbf{Synthèse des axes principaux} (une phrase par axe).
    \item \textbf{Ouverture / Engagement personnel :} Projection future, citation, résolution, appel à l'action.
\end{enumerate}
\end{propriete}

\begin{exemplebox}{Conclusion modèle (profil infirmier)}
\eng{To sum up, anticipating my English needs as a future nurse in Bamenda is not a luxury but a professional duty. Mastering medical terminology will ensure accurate patient assessment, while fluency in writing will guarantee clear handover reports. The strategies I have outlined — enrolling in a Medical English course, practising with Anki, and completing an internship in Buea — are concrete steps I am taking today. Ultimately, passing the OET will be the proof that I am ready to serve both Anglophone and Francophone patients with the same excellence. ``A language is not just a tool; it is a bridge to better care.''}
\end{exemplebox}

\coursec{Connecteurs, temps et vocabulaire utiles}

Cette section rassemble les outils linguistiques indispensables pour une composition cohérente, précise et de niveau B2/C1.

\courssub{Connecteurs logiques (Linking words)}

\begin{center}
\begin{tabularx}{\linewidth}{|X|X|X|}
\hline
\rowcolor{popblueL}
\textbf{Fonction} & \textbf{Connecteurs (English)} & \textbf{Exemples d'utilisation} \\
\hline
Addition & \eng{furthermore, moreover, in addition, additionally, also} & \eng{Furthermore, I will need to write referral letters.} \\
\hline
Contraste / Concession & \eng{however, nevertheless, on the other hand, although, while} & \eng{Although French is my first language, English is mandatory in my field.} \\
\hline
Cause / Raison & \eng{because, since, as, due to, owing to} & \eng{Due to the bilingual nature of Cameroon, English is unavoidable.} \\
\hline
Conséquence / Résultat & \eng{therefore, consequently, as a result, thus, hence} & \eng{I will take a course; therefore, I will be better prepared.} \\
\hline
But / Finalité & \eng{to, in order to, so as to, so that} & \eng{I practise daily so that I can communicate fluently.} \\
\hline
Exemplification & \eng{for example, for instance, such as, namely} & \eng{I will need specific vocabulary, for instance, medical terms.} \\
\hline
Ordre / Séquence & \eng{firstly, secondly, finally, subsequently, eventually} & \eng{Firstly, I will enrol in a course. Secondly, I will practise daily.} \\
\hline
Condition & \eng{if, unless, provided that, as long as} & \eng{Provided that I study regularly, I will reach B2 level.} \\
\hline
\end{tabularx}
\end{center}

\begin{attention}{Pièges de connecteurs (Francophones)}
\begin{itemize}
    \item \eng{Moreover} ne s'utilise pas en début de phrase pour ajouter une idée mineure ; il renforce un argument fort.
    \item \eng{Because of / Due to} + \textbf{nom / gérondif} (pas de clause). \eng{Because} + \textbf{sujet + verbe}.
    \item \eng{However} (adverbe) $\neq$ \eng{But} (conjonction). \eng{However} demande un point-virgule ou un point avant, et une virgule après : \eng{English is hard; however, I will succeed.}
    \item \eng{Also} ne commence généralement pas une phrase en style formel ; préférez \eng{In addition} ou \eng{Furthermore}.
\end{itemize}
\end{attention}

\courssub{Temps verbaux et modaux pour l'anticipation}

Le sujet impose une projection dans le futur. Variez les formes pour montrer votre maîtrise.

\begin{center}
\begin{tabularx}{\linewidth}{|X|X|X|}
\hline
\rowcolor{popblueL}
\textbf{Forme} & \textbf{Emploi} & \textbf{Exemples traduits} \\
\hline
\eng{will + BV} (Future Simple) & Prédictions certaines, engagements fermes & \eng{I will need English for reports.} (J'aurai besoin...) \\
\hline
\eng{be going to + BV} & Intentions, plans déjà décidés & \eng{I am going to take a Business English course.} (Je vais suivre...) \\
\hline
\eng{anticipate + V-ing} & Verbe du sujet : prévoir activement & \eng{I anticipate needing strong writing skills.} (J'anticipe avoir besoin...) \\
\hline
\eng{expect + to + BV} & Attente (plus passif) — \textbf{à éviter comme verbe principal} & \eng{I expect to use English daily.} (Je m'attends à...) \\
\hline
\eng{might / may / could + BV} & Probabilité, éventualité (modaux) & \eng{I might have to negotiate contracts.} (Il se peut que je doive...) \\
\hline
\eng{will have to / will need to} & Obligation future, nécessité & \eng{I will have to write minutes of meetings.} (Je devrai rédiger...) \\
\hline
\eng{present perfect} (\eng{have/has + V-en}) & Expérience passée liée au futur & \eng{I have already started an online course.} (J'ai déjà commencé...) \\
\hline
\end{tabularx}
\end{center}

\courssub{Vocabulaire thématique : compétences et stratégies}

\begin{definition}{Lexique clé pour « Anticipating L2 needs at work »}
\begin{itemize}
    \item \textbf{Compétences (Skills) :} \cle{proficiency} (maîtrise), \cle{fluency} (aisance), \cle{accuracy} (correction), \cle{terminology} (terminologie), \cle{jargon} (jargon), \cle{drafting} (rédaction), \cle{proofreading} (relecture), \cle{minutes} (procès-verbal), \cle{report writing} (rédaction de rapports), \cle{public speaking} (prise de parole en public), \cle{negotiation skills} (compétences en négociation).
    \item \textbf{Stratégies (Strategies) :} \cle{enroll in} (s'inscrire à), \cle{undertake} (entreprendre), \cle{self-study} (auto-formation), \cle{immersion} (immersion), \cle{language exchange} (échange linguistique), \cle{certification} (certification), \cle{mock interview} (entretien simulé), \cle{flashcards} (fiches de vocabulaire), \cle{MOOC} (cours en ligne ouvert), \cle{webinar} (webinaire).
    \item \textbf{Contextes professionnels (Work contexts) :} \cle{workplace} (lieu de travail), \cle{stakeholder} (partie prenante), \cle{deadline} (échéance), \cle{multilingual environment} (environnement multilingue), \cle{cross-cultural communication} (communication interculturelle), \cle{remote work} (télétravail), \cle{on-the-job training} (formation en poste).
\end{itemize}
\end{definition}

\begin{exemplebox}{Collocations utiles (associations de mots naturelles)}
\begin{center}
\begin{tabularx}{\linewidth}{|X|X|}
\hline
\rowcolor{popblueL}
\textbf{Collocation} & \textbf{Traduction} \\
\hline
\eng{bridge the language gap} & combler le fossé linguistique \\
\hline
\eng{enhance my employability} & améliorer mon employabilité \\
\hline
\eng{stay ahead of the curve} & garder une longueur d'avance \\
\hline
\eng{invest in language training} & investir dans la formation linguistique \\
\hline
\eng{master the jargon of...} & maîtriser le jargon de... \\
\hline
\eng{communicate effectively with...} & communiquer efficacement avec... \\
\hline
\eng{draft professional emails} & rédiger des courriels professionnels \\
\hline
\eng{attend international conferences} & assister à des conférences internationales \\
\hline
\end{tabularx}
\end{center}
\end{exemplebox}

\coursec{Modèle commenté}

Voici une composition complète (environ 250 mots) suivie d'une analyse détaillée. Le profil : élève de Terminale A, futur ingénieur civil à Douala.

\begin{textebox}{Modèle de composition : Anticipating my English needs as a civil engineer}
In today's interconnected world, infrastructure projects in Cameroon increasingly involve international partners, making English proficiency a decisive asset for engineers. As a final-year student in Douala preparing to enter the École Nationale Supérieure des Travaux Publics, I have a clear vision of the linguistic demands of my future career. In this composition, I will describe the specific English skills I will need on construction sites and in design offices, and the proactive steps I am taking to acquire them.

Firstly, technical reading and writing will be at the core of my daily work. I will have to interpret international standards such as Eurocodes, read technical specifications from foreign suppliers, and draft site reports and correspondence with stakeholders. Misunderstanding a clause in a contract or a specification sheet could lead to costly errors or safety hazards. Therefore, mastering technical vocabulary — terms like \eng{reinforced concrete}, \eng{load-bearing capacity}, \eng{geotechnical survey} — is not optional.

Secondly, oral communication will be crucial during site meetings with expatriate consultants and during negotiations with contractors from Nigeria or China. I will need to give clear instructions, explain technical solutions, and participate in progress reviews. To prepare, I am currently following a ``Technical English for Engineers'' MOOC on Coursera and I practise speaking with a language partner via Zoom every week.

Finally, I plan to obtain the \eng{IELTS Academic} certification before graduation, as many international firms require it. I also intend to do my final-year internship in a bilingual firm in Douala to gain real-world exposure.

In conclusion, anticipating my English needs today means securing my professional credibility tomorrow. By combining specialised courses, daily practice, and targeted certification, I am building the linguistic foundation that will allow me to contribute effectively to Cameroon's infrastructure development. ``Engineering is not only about structures; it is about communication.''
\end{textebox}

\begin{methode}{Analyse du modèle (pourquoi ça marche)}
\begin{enumerate}
    \item \textbf{Structure respectée :} Introduction (hook + contexte + thèse), 3 paragraphes de développement (axes distincts), conclusion (synthèse + ouverture).
    \item \textbf{Thesis statement claire :} \eng{``I will describe the specific English skills I will need... and the proactive steps I am taking...''} Annonce les 2 axes (compétences + stratégies) mais le développement en fait 3 (lecture/écriture, oral, certification/stage) — c'est acceptable car le 3e axe développe les stratégies.
    \item \textbf{Connecteurs variés :} \eng{Firstly, Therefore, Secondly, Finally, In conclusion}.
    \item \textbf{Temps et modaux appropriés :} \eng{will have to, will need to, am currently following, plan to obtain, intend to do}.
    \item \textbf{Vocabulaire précis :} \eng{Eurocodes, load-bearing capacity, geotechnical survey, expatriate consultants, progress reviews, IELTS Academic}.
    \item \textbf{Ancrage camerounais :} Douala, ENSTP, stages, firmes bilingues, partenaires nigérians/chinois.
    \item \textbf{Ton personnel et engagé :} \eng{I have a clear vision, I am building...}
\end{enumerate}
\end{methode}

\coursec{Erreurs à éviter et relecture}

Au-delà des erreurs de planification déjà vues, voici les pièges linguistiques les plus fréquents chez les francophones. Relisez votre copie en traquant ces points.

\begin{attention}{Faux amis (False friends) — Liste de survie}
\begin{center}
\begin{tabularx}{\linewidth}{|X|X|X|}
\hline
\rowcolor{popblueL}
\textbf{Mot anglais} & \textbf{Piège (sens français ressemblant)} & \textbf{Vrai sens / Traduction correcte} \\
\hline
\cle{actually} & actuellement & en fait, réellement \\
\hline
\cle{eventually} & éventuellement & finalement, à la fin \\
\hline
\cle{opportunity} & opportunité (risque) & occasion, chance favorable \\
\hline
\cle{realise} & réaliser (créer) & se rendre compte, comprendre \\
\hline
\cle{resume} & résumer & reprendre (après pause) ; \eng{CV} = \eng{résumé} \\
\hline
\cle{sensible} & sensible (qui sent) & raisonnable, sensible (émotionnellement) \\
\hline
\cle{attend} & attendre & assister à, être présent \\
\hline
\cle{demand} & demander (poliment) & exiger, réclamer \\
\hline
\cle{report} & reporter (remettre à plus tard) & rapport, compte rendu ; \eng{to report} = signaler \\
\hline
\cle{stage} & stage (période de formation) & scène, étape ; \eng{stage} = \eng{internship} / \eng{work placement} \\
\hline
\end{tabularx}
\end{center}
\end{attention}

\begin{attention}{Erreurs grammaticales récurrentes}
\begin{itemize}
    \item \textbf{Omission du sujet :} \eng{*Is important to learn English.} → \eng{\textbf{It} is important to learn English.} (Sujet \eng{it} impératif).
    \item \textbf{Ordre des mots (SVOMPT) :} \eng{*I need at work English.} → \eng{I need English \textbf{at work}.} (Complément de lieu à la fin).
    \item \textbf{Prépositions après verbes/adjectifs :} \eng{interested \textbf{in}}, \eng{good \textbf{at}}, \eng{responsible \textbf{for}}, \eng{apply \textbf{for}}, \eng{consist \textbf{of}}, \eng{depend \textbf{on}}.
    \item \textbf{Gérondif après préposition :} \eng{anticipate \textbf{needing}}, \eng{plan \textbf{on doing}}, \eng{look forward to \textbf{hearing}}.
    \item \textbf{Accord sujet-verbe :} \eng{The needs \textbf{are} high.} (pluriel) ; \eng{English \textbf{is} essential.} (singulier).
    \item \textbf{Articles :} \eng{at \textbf{work}} (zéro article), \eng{in \textbf{the} workplace}, \eng{\textbf{a} nurse}, \eng{\textbf{the} OET}.
    \item \textbf{Verbes irréguliers (forme participe passé) :} \eng{write/wrote/written}, \eng{speak/spoke/spoken}, \eng{read/read/read}, \eng{begin/began/begun}, \eng{take/took/taken}.
\end{itemize}
\end{attention}

\begin{attention}{Orthographe et ponctuation}
\begin{itemize}
    \item \textbf{Majuscules :} Jours, mois, langues, nationalités, \eng{I} (toujours majuscule), noms propres (\eng{Cameroon, Douala, IELTS, British Council}).
    \item \textbf{Virgule après connecteur en début de phrase :} \eng{However, ...}, \eng{Therefore, ...}, \eng{For example, ...}.
    \item \textbf{Apostrophe :} \eng{it's} = \eng{it is} / \eng{it has} ; \eng{its} = possessif (pas d'apostrophe).
    \item \textbf{Guillemets :} Utilisez “ ... ” pour les citations anglaises (pas `` ... '' ni " ... ").
\end{itemize}
\end{attention}

\begin{methode}{Check-list de relecture (5 minutes avant de rendre)}
\begin{enumerate}
    \item \textbf{Sujet :} Ai-je répondu exactement à la consigne (anticipation + stratégies) ?
    \item \textbf{Structure :} Intro (hook, thèse) / 2-3 dev. (1 idée = 1 §) / Conclusion (synthèse + ouverture) ?
    \item \textbf{Connecteurs :} Au moins 5 connecteurs variés utilisés correctement ?
    \item \textbf{Temps/Modaux :} Futur, modaux de probabilité, \eng{anticipate + V-ing} présents ?
    \item \textbf{Vocabulaire :} Mots techniques du métier, pas de vocabulaire trop général ?
    \item \textbf{Grammaire :} Sujets présents, accords, prépositions, gérondifs, articles ?
    \item \textbf{Orthographe :} Majuscules, virgules, apostrophes, guillemets anglais ?
    \item \textbf{Longueur :} Environ 250-300 mots (selon consignes Bac) ?
\end{enumerate}
\end{methode}

\coursec{Grille d'évaluation et entraînement au Bac}

Comprendre la grille officielle vous permet de cibler les points qui rapportent.

\courssub{Grille type d'évaluation (sur 20 points)}

\begin{center}
\begin{tabularx}{\linewidth}{|X|c|X|}
\hline
\rowcolor{popblueL}
\textbf{Critère} & \textbf{Points} & \textbf{Indicateurs de réussite} \\
\hline
\textbf{Contenu \& Pertinence} & 5 & Sujet respecté, angle personnel, idées développées, exemples concrets, stratégies citées. \\
\hline
\textbf{Organisation \& Cohérence} & 4 & Structure intro/dev/concl, paragraphes logiques, connecteurs fluides, thesis statement clair. \\
\hline
\textbf{Grammaire \& Syntaxe} & 4 & Temps/modaux appropriés, accords, ordre des mots, phrases complexes correctes, peu d'erreurs bloquantes. \\
\hline
\textbf{Vocabulaire \& Registre} & 4 & Lexique professionnel précis, collocations, variété, registre formel/semi-formel, pas de franglais. \\
\hline
\textbf{Mécanique (Orthographe, Ponctuation)} & 3 & Majuscules, virgules, apostrophes, guillemets, orthographe des mots courants. \\
\hline
\textbf{Total} & 20 & \\
\hline
\end{tabularx}
\end{center}

\courssub{Sujet d'entraînement type Bac}

\begin{textebox}{Sujet d'entraînement (2 h)}
\eng{Write a composition of about 250-300 words on the following topic:}

\eng{``As a future professional in Cameroon, how do you anticipate your second language (English) needs at work? Illustrate with specific skills and concrete preparation strategies.''}

\textbf{Consignes :}
\begin{itemize}
    \item Choisissez votre vrai projet professionnel.
    \item Respectez la structure en 3 parties.
    \item Utilisez au moins 6 connecteurs différents.
    \item Intégrez du vocabulaire technique de votre domaine.
    \item Relisez-vous avec la check-list (section 5).
\end{itemize}
\end{textebox}

\begin{experience}{Atelier d'écriture guidée (en classe, 45 min)}
\begin{enumerate}
    \item \textbf{5 min :} Analyse du sujet (mots-clés, questions W) — individuel.
    \item \textbf{10 min :} Brainstorming + carte mentale — individuel.
    \item \textbf{5 min :} Échange en binôme : présentez votre plan (thesis statement + 2 axes) à votre voisin ; il valide la cohérence.
    \item \textbf{20 min :} Rédaction de la composition (brouillon puis propre).
    \item \textbf{5 min :} Relecture croisée avec la check-list (section 5) — binôme corrige le vôtre, vous corrigez le sien.
\end{enumerate}
\end{experience}

\begin{exemplebox}{Exemple de thesis statements selon profils (pour s'inspirer)}
\begin{itemize}
    \item \textbf{Future Accountant (Comptable) :} \eng{As a future chartered accountant in Yaoundé, I anticipate needing advanced financial English to interpret IFRS standards, audit multinational clients, and draft audit reports for international partners.}
    \item \textbf{Future Journalist (Journaliste) :} \eng{In a media landscape where 80\% of global news sources are in English, I anticipate needing strong reading and writing skills to monitor international wires, conduct interviews with Anglophone sources, and publish bilingual articles.}
    \item \textbf{Future Hotel Manager (Directeur d'hôtel) :} \eng{Working in the hospitality industry in Limbe means serving a global clientele; I therefore anticipate needing fluent spoken English for guest relations, written English for online marketing, and cultural awareness for conflict resolution.}
\end{itemize}
\end{exemplebox}

\begin{corrige}{Proposition de corrigé type (profil Ingénieur civil — adapté du modèle)}
\textbf{Note :} Ce corrigé sert de référence ; votre composition doit refléter \textbf{votre} projet.

\eng{In today's interconnected world, infrastructure projects in Cameroon increasingly involve international partners, making English proficiency a decisive asset for engineers. As a final-year student in Douala preparing to enter the École Nationale Supérieure des Travaux Publics, I have a clear vision of the linguistic demands of my future career. In this composition, I will describe the specific English skills I will need on construction sites and in design offices, and the proactive steps I am taking to acquire them.}

\eng{Firstly, technical reading and writing will be at the core of my daily work. I will have to interpret international standards such as Eurocodes, read technical specifications from foreign suppliers, and draft site reports and correspondence with stakeholders. Misunderstanding a clause in a contract or a specification sheet could lead to costly errors or safety hazards. Therefore, mastering technical vocabulary — terms like reinforced concrete, load-bearing capacity, geotechnical survey — is not optional.}

\eng{Secondly, oral communication will be crucial during site meetings with expatriate consultants and during negotiations with contractors from Nigeria or China. I will need to give clear instructions, explain technical solutions, and participate in progress reviews. To prepare, I am currently following a ``Technical English for Engineers'' MOOC on Coursera and I practise speaking with a language partner via Zoom every week.}

\eng{Finally, I plan to obtain the IELTS Academic certification before graduation, as many international firms require it. I also intend to do my final-year internship in a bilingual firm in Douala to gain real-world exposure.}

\eng{In conclusion, anticipating my English needs today means securing my professional credibility tomorrow. By combining specialised courses, daily practice, and targeted certification, I am building the linguistic foundation that will allow me to contribute effectively to Cameroon's infrastructure development. ``Engineering is not only about structures; it is about communication.''}

\textbf{Auto-évaluation (grille) :}
\begin{itemize}
    \item Contenu : 5/5 (spécifique, stratégies concrètes, ancrage local).
    \item Organisation : 4/4 (thesis claire, 3 § dev, connecteurs, conclusion).
    \item Grammaire : 4/4 (futur, modaux, gérondifs, passif, pas d'erreur).
    \item Vocabulaire : 4/4 (termes techniques, collocations, registre formel).
    \item Mécanique : 3/3 (ponctuation, majuscules, guillemets corrects).
    \item \textbf{Total : 20/20}
\end{itemize}
\end{corrige}

\begin{savaistu}[Le Baccalauréat et l'épreuve d'expression écrite]
\begin{itemize}
    \item L'épreuve d'anglais au Bac (série A) dure 2 h et comprend généralement : Compréhension (écrite/orale), Grammaire/Vocabulaire, Expression écrite (composition ou lettre/email).
    \item La composition vaut souvent 8 à 10 points sur 20. C'est la partie où l'écart se creuse le plus entre candidats.
    \item Les correcteurs suivent une grille harmonisée nationale. Les critères ci-dessus en sont une transcription fidèle.
    \item \textbf{Astuce :} Même si votre niveau de langue n'est pas parfait, une \textbf{structure impeccable} (intro/thesis/dev/concl) et un \textbf{contenu pertinent} (métier réel, stratégies réelles) vous assurent la moyenne (10/20) sur la composition. La langue fait la différence pour l'excellence.
\end{itemize}
\end{savaistu}


\coursec{La structure d'une composition}

Dans la section précédente, vous avez appris à analyser le sujet (\eng{anticipating one's second language needs at work}), à chercher des idées et à les organiser en un plan. Il est temps maintenant de donner à ce plan une \textbf{forme} : celle d'une composition anglaise bien construite. En anglais comme en français, un bon essai repose sur une architecture en trois parties. Mais attention : les attentes des correcteurs anglophones ne sont pas exactement les mêmes que celles des correcteurs francophones. Cette section vous donne le plan de construction complet, pièce par pièce.

\courssub{Vue d'ensemble : les trois parties et leurs proportions}

\begin{definition}[Composition / essay]
Une \cle{composition} (en anglais \eng{essay} ou \eng{composition}) est un texte argumentatif ou explicatif organisé en trois parties : une \cle{introduction} qui présente le sujet, un \cle{development} (ou \eng{body}, le « corps » du texte) qui développe les idées, et une \cle{conclusion} qui les synthétise. Au Baccalauréat, la longueur attendue est d'environ \textbf{250 à 300 mots}.
\end{definition}

Voici un extrait de composition rédigé par une élève de Terminale. Observez comment les trois parties s'enchaînent.

\begin{textebox}[A well-structured essay (extract)]
\textbf{[Introduction]} English is often called the language of business, and this is especially true in a bilingual country like Cameroon. As I prepare to enter the world of work, I must think carefully about the English skills I will need. This essay will discuss the English skills required in my future career and explain how I plan to acquire them.

\textbf{[Body, paragraph 1]} First of all, English is essential in international trade. Many companies in Douala do business with partners in Nigeria, Ghana and beyond, and their contracts, emails and meetings are in English. For example, my uncle, who works for an import company, uses English every day to negotiate prices with foreign suppliers.

\textbf{[Body, paragraph 2]} Secondly, I will need spoken English for job interviews and daily communication. Employers in Yaoundé increasingly test candidates in English, even for local positions. A friend of mine lost a good job last year simply because she could not answer the interviewer's questions fluently.

\textbf{[Conclusion]} In conclusion, English will play a central role in my professional life, both in writing and in speaking. That is why I have decided to read the English press every week and to practise speaking with my classmates. If I start preparing now, I am confident that the language will open doors for me instead of closing them.
\end{textebox}

\begin{center}
\begin{tabularx}{\linewidth}{|l|X|X|}\hline
\rowcolor{popblueL} \textbf{Partie} & \textbf{Proportion} & \textbf{Fonction} \\ \hline
Introduction & 10--15 \% (3--4 phrases) & Accrocher le lecteur, présenter le sujet, annoncer le plan \\ \hline
Development (body) & environ 70 \% (2--3 paragraphes) & Développer les arguments : idée + explication + exemple \\ \hline
Conclusion & 10--15 \% (3--4 phrases) & Résumer, donner son opinion, ouvrir vers l'avenir \\ \hline
\end{tabularx}
\end{center}

\begin{popfigure}
\begin{tikzpicture}[font=\small]
  \node[draw=popblue, fill=popblueL, rounded corners, minimum width=10cm, minimum height=1.7cm, align=center] (intro) at (0,0)
    {\textbf{INTRODUCTION} (10--15 \%)\\ hook (accroche) + presentation of the topic + thesis statement};
  \node[draw=poporange, fill=poporangeL, rounded corners, minimum width=10cm, minimum height=2.9cm, align=center] (body) at (0,-2.7)
    {\textbf{BODY / DEVELOPMENT} (70 \%)\\ Paragraph 1 : topic sentence + explanation + example\\ Paragraph 2 : topic sentence + explanation + example\\ Paragraph 3 : topic sentence + explanation + example};
  \node[draw=popgreen, fill=popgreenL, rounded corners, minimum width=10cm, minimum height=1.7cm, align=center] (concl) at (0,-5.4)
    {\textbf{CONCLUSION} (10--15 \%)\\ summary of main ideas + personal opinion + opening to the future};
  \draw[-{Stealth}, thick, popmuted] (intro.south) -- (body.north);
  \draw[-{Stealth}, thick, popmuted] (body.south) -- (concl.north);
\end{tikzpicture}
\legende{L'architecture d'une composition : trois blocs, trois fonctions.}
\end{popfigure}

\begin{attention}[Pas de titre, pas de plan apparent]
Contrairement à la dissertation française, la composition anglaise au Bac ne porte \textbf{pas de titre} et les parties ne sont \textbf{jamais annoncées par des sous-titres} (\eng{Introduction}, \eng{Development}...). La structure doit être visible uniquement grâce aux paragraphes (on saute une ligne ou on met un alinéa) et aux connecteurs.
\end{attention}

\courssub{L'introduction : hook, sujet, thesis statement}

L'introduction se construit en \textbf{trois mouvements}, comme un entonnoir : on part du général pour arriver au précis.

\begin{enumerate}
  \item \textbf{Le hook} (l'accroche) : une phrase qui attire l'attention du lecteur. Ce peut être une question, un constat général, un fait de la vie quotidienne.
  \item \textbf{La présentation du sujet} : une ou deux phrases qui resserrent le propos et le relient au thème imposé.
  \item \textbf{La thesis statement} : la phrase qui annonce clairement ce que la composition va démontrer ou décrire.
\end{enumerate}

\begin{definition}[Thesis statement]
La \cle{thesis statement} (littéralement « phrase-thèse ») est la phrase, généralement placée \textbf{à la fin de l'introduction}, qui annonce clairement le plan ou l'idée directrice de la composition. Elle dit au lecteur \emph{ce que} le texte va traiter et \emph{dans quel ordre}. Exemple : \eng{This essay will discuss the English skills I will need at work and how I plan to acquire them.} « Cette composition examinera les compétences en anglais dont j'aurai besoin au travail et la manière dont je compte les acquérir. »
\end{definition}

\begin{exemplebox}[Trois hooks possibles pour notre sujet]
\begin{itemize}
  \item \textbf{Question :} \eng{Have you ever imagined yourself in a job interview conducted entirely in English?} « Vous êtes-vous déjà imaginé dans un entretien d'embauche mené entièrement en anglais ? »
  \item \textbf{Constat général :} \eng{English is often called the language of business.} « On dit souvent que l'anglais est la langue des affaires. »
  \item \textbf{Fait concret :} \eng{In Douala, more and more job advertisements require “good command of English”.} « À Douala, de plus en plus d'offres d'emploi exigent une “bonne maîtrise de l'anglais”. »
\end{itemize}
\end{exemplebox}

\begin{propriete}[Formules utiles pour la thesis statement]
\begin{itemize}
  \item \eng{This essay will discuss ... and explain ...} « Cette composition examinera ... et expliquera ... »
  \item \eng{This essay will examine the reasons why ...} « Cette composition examinera les raisons pour lesquelles ... »
  \item \eng{In this essay, I will first describe ..., then I will suggest ...} « Dans cette composition, je décrirai d'abord ..., puis je proposerai ... »
\end{itemize}
Notez l'emploi du \textbf{futur} (\eng{will}) dans ces formules : on annonce ce que le texte \emph{va} faire.
\end{propriete}

\begin{attention}[L'introduction-pavé : l'erreur n° 1 des francophones]
Beaucoup d'élèves écrivent \textbf{une seule longue phrase} comme introduction, souvent une simple reformulation du sujet. C'est insuffisant et cela coûte des points. Une bonne introduction compte \textbf{3 à 4 phrases courtes et claires} : une accroche, une présentation du sujet, une thesis statement. De même, ne copiez jamais le sujet mot pour mot : reformulez-le avec vos propres mots (\eng{paraphrase}).
\end{attention}

\courssub{Le développement : la règle d'or du paragraphe}

Le développement (\eng{body}) représente environ 70 \% de la composition, soit \textbf{2 à 3 paragraphes}. Chaque paragraphe obéit à une règle absolue.

\begin{propriete}[Une seule idée par paragraphe]
Chaque paragraphe développe \textbf{une seule idée principale}, introduite par une \cle{topic sentence} (phrase d'ouverture ou « phrase-chapeau ») placée en tête de paragraphe. Toutes les autres phrases du paragraphe doivent expliquer ou illustrer cette idée. Si vous changez d'idée, vous changez de paragraphe.
\end{propriete}

\begin{methode}[Écrire un bon paragraphe : le schéma TEE]
Retenez l'acronyme \cle{TEE} : \textbf{T}opic sentence, \textbf{E}xplanation, \textbf{E}xample.
\begin{enumerate}
  \item \textbf{Topic sentence} : annoncez l'idée principale du paragraphe en une phrase claire. \eng{First of all, English is essential in international trade.}
  \item \textbf{Explanation} : développez cette idée en une ou deux phrases (pourquoi ? comment ? dans quel contexte ?). \eng{Many companies in Douala do business with partners in Nigeria and Ghana, and their contracts and emails are in English.}
  \item \textbf{Example} : donnez un exemple concret, précis, si possible tiré du contexte camerounais ou de votre expérience. \eng{For example, my uncle uses English every day to negotiate with foreign suppliers.}
\end{enumerate}
Un paragraphe complet compte donc en général \textbf{3 à 5 phrases}.
\end{methode}

\begin{exemplebox}[Topic sentences et exemples traduits]
\begin{itemize}
  \item \eng{First of all, English is essential in international trade.} « Tout d'abord, l'anglais est essentiel dans le commerce international. »
  \item \eng{Secondly, I will need spoken English for job interviews.} « Ensuite, j'aurai besoin de l'anglais oral pour les entretiens d'embauche. »
  \item \eng{Finally, I plan to improve my English before leaving school.} « Enfin, je compte améliorer mon anglais avant de quitter le lycée. »
  \item \eng{For example, many banks in Yaoundé require bilingual secretaries.} « Par exemple, de nombreuses banques à Yaoundé exigent des secrétaires bilingues. »
\end{itemize}
\end{exemplebox}

\begin{center}
\begin{tabularx}{\linewidth}{|X|X|X|}\hline
\rowcolor{popblueL} \textbf{Français} & \textbf{English} & \textbf{Remarque} \\ \hline
phrase d'ouverture & topic sentence & placée en tête de paragraphe \\ \hline
d'abord / tout d'abord & first of all / firstly & jamais \eng{at first} (= « au début ») \\ \hline
ensuite, deuxièmement & secondly / then & \\ \hline
enfin, finalement & finally / lastly & \\ \hline
par exemple & for example / for instance & abréviation \eng{e.g.} à éviter en rédaction \\ \hline
en conclusion & in conclusion / to conclude & jamais \eng{in the conclusion} \\ \hline
c'est pourquoi & that is why & \\ \hline
\end{tabularx}
\end{center}

\begin{attention}[Faux amis et pièges du paragraphe]
\begin{itemize}
  \item \eng{Actually} ne signifie pas « actuellement » mais « en fait » (prononciation : « èk-chou-eu-li »). Pour « actuellement », dites \eng{currently} ou \eng{at present}.
  \item \eng{An example} se prononce « ène eg-zam-peul » ; ne confondez pas avec \eng{a sample} (un échantillon).
  \item Ne traduisez pas « dans un premier temps » par \eng{in a first time} : dites \eng{first of all} ou \eng{firstly}.
  \item Chaque nouveau paragraphe commence par un \textbf{alinéa} ou une ligne sautée : sans cela, le correcteur ne voit pas votre plan.
\end{itemize}
\end{attention}

\courssub{La conclusion : résumer, opiner, ouvrir}

La conclusion est le miroir de l'introduction. Elle se construit elle aussi en trois mouvements, mais dans l'ordre inverse : du précis vers le général.

\begin{enumerate}
  \item \textbf{Le rappel des idées principales} (\eng{summary}) : reformulez en une phrase les deux ou trois points du développement, sans les répéter mot pour mot.
  \item \textbf{L'opinion personnelle} (\eng{personal opinion}) : dites clairement ce que vous pensez, avec \eng{I think}, \eng{I believe}, \eng{in my opinion}.
  \item \textbf{L'ouverture vers l'avenir} (\eng{opening}) : terminez par une phrase qui projette le lecteur vers l'avenir — une résolution, un espoir, une recommandation.
\end{enumerate}

\begin{exemplebox}[Une conclusion modèle, phrase par phrase]
\eng{In conclusion, English will play a central role in my professional life.} (rappel) « En conclusion, l'anglais jouera un rôle central dans ma vie professionnelle. »\par
\eng{That is why I believe that every student should start preparing for it now.} (opinion) « C'est pourquoi je pense que chaque élève devrait commencer à s'y préparer dès maintenant. »\par
\eng{If we master English today, it will open doors for us tomorrow.} (ouverture) « Si nous maîtrisons l'anglais aujourd'hui, il nous ouvrira des portes demain. »
\end{exemplebox}

\begin{attention}[Ce qu'une conclusion ne doit jamais faire]
\begin{itemize}
  \item Ne introduisez \textbf{aucune idée nouvelle} dans la conclusion : elle ne fait que synthétiser.
  \item Ne terminez pas brutalement après le dernier paragraphe du développement : une composition sans conclusion est pénalisée.
  \item Évitez les formules creuses comme \eng{That is all I have to say.} (« C'est tout ce que j'ai à dire »), qui donnent une impression de paresse.
\end{itemize}
\end{attention}

\courssub{Entraînement : reconstruire une composition}

\begin{exoresolu}[Remettre les phrases dans l'ordre]
Voici les huit phrases d'une composition mélangées. Classez-les dans l'ordre correct (introduction, développement en deux paragraphes, conclusion).\par
a) \eng{For instance, a friend of mine was hired by a bank in Buea mainly because of her fluent English.}\par
b) \eng{This essay will discuss why English matters for my career and how I intend to improve it.}\par
c) \eng{In conclusion, English is a key that I must obtain before entering the job market.}\par
d) \eng{First of all, most well-paid jobs in Cameroon now require a good command of English.}\par
e) \eng{In today's Cameroon, speaking English is no longer a luxury but a necessity.}\par
f) \eng{Secondly, I plan to read English newspapers and watch English films every week.}\par
g) \eng{Employers want workers who can write emails and speak to foreign clients.}\par
h) \eng{If I work hard now, I am sure the language will serve me well in the future.}
\tcblower
\textbf{Solution détaillée.}\par
\textbf{Introduction :} e) est le hook (constat général), b) est la thesis statement (elle annonce le plan : \eng{why English matters ... and how I intend to improve it}).\par
\textbf{Paragraphe 1 :} d) est la topic sentence (annoncée par \eng{First of all}), g) est l'explication (\eng{Employers want workers who...}), a) est l'exemple (signalé par \eng{For instance}).\par
\textbf{Paragraphe 2 :} f) est la topic sentence (annoncée par \eng{Secondly}) ; elle développe la deuxième partie du plan (\eng{how I intend to improve it}).\par
\textbf{Conclusion :} c) résume (\eng{In conclusion}), h) ouvre vers l'avenir (\eng{in the future}).\par
Ordre final : \textbf{e -- b -- d -- g -- a -- f -- c -- h}. Remarquez que la thesis statement (b) annonce exactement les deux paragraphes du développement : c'est le signe d'une composition bien construite.
\end{exoresolu}

\begin{exercice}[À vous de jouer]
Rédigez uniquement l'\textbf{introduction} (3 à 4 phrases : hook, sujet, thesis statement) d'une composition sur le sujet suivant : \eng{“Write a composition about the English skills you will need in your future job and how you plan to acquire them.”} Puis écrivez les \textbf{trois topic sentences} de votre développement. Vous comparerez votre production au modèle commenté de la section 4.
\end{exercice}

\begin{aretenir}[L'essentiel de la section]
\begin{itemize}
  \item Une composition = \textbf{introduction} (10--15 \%) + \textbf{développement} (70 \%, 2 à 3 paragraphes) + \textbf{conclusion} (10--15 \%), pour un total de 250 à 300 mots au Bac.
  \item L'introduction se termine par une \textbf{thesis statement} qui annonce le plan.
  \item Chaque paragraphe suit le schéma \textbf{TEE} : topic sentence, explanation, example — une seule idée par paragraphe.
  \item La conclusion résume, donne une opinion et ouvre vers l'avenir, sans idée nouvelle.
  \item La structure se voit grâce aux \textbf{alinéas} et aux \textbf{connecteurs}, jamais grâce à des sous-titres.
\end{itemize}
\end{aretenir}

\coursec{Connecteurs, temps et vocabulaire utiles}

Dans la section précédente, nous avons appris qu'une composition bien construite repose sur trois parties : une introduction qui pose le sujet, un développement organisé en paragraphes, et une conclusion qui récapitule. Mais une architecture solide ne suffit pas : il faut encore des \cle{mortiers} pour lier les pierres entre elles. En anglais, ce mortier, ce sont les \cle{linking words} (connecteurs logiques), les temps verbaux bien choisis et un vocabulaire précis. C'est exactement ce que le correcteur du Bac évalue dans la rubrique « coherence and cohesion ». Cette section vous donne la boîte à outils complète pour rédiger une composition sur le thème : \eng{anticipating one's second language needs at work} (anticiper ses besoins en langue seconde au travail).

\courssub{Les connecteurs logiques (linking words)}

\courssub{Observer avant de généraliser}

Lisez d'abord ce mini-modèle, extrait d'une composition d'élève, et observez les mots en gras : ils guident le lecteur comme des panneaux routiers.

\begin{textebox}[A mini-model to memorise]
\eng{\textbf{First of all}, English is the language of international business, so every ambitious worker should learn it. \textbf{Moreover}, many companies in Douala and Yaoundé now require English during job interviews. \textbf{However}, learning a language takes time and effort. \textbf{Therefore}, I have decided to attend evening classes and to practise every day. \textbf{To sum up}, anticipating my English needs today means securing my career tomorrow.}
\end{textebox}

« Tout d'abord, l'anglais est la langue des affaires internationales, donc tout travailleur ambitieux devrait l'apprendre. De plus, de nombreuses entreprises à Douala et Yaoundé exigent désormais l'anglais lors des entretiens d'embauche. Cependant, apprendre une langue demande du temps et des efforts. C'est pourquoi j'ai décidé de suivre des cours du soir et de pratiquer chaque jour. En résumé, anticiper mes besoins en anglais aujourd'hui, c'est sécuriser ma carrière demain. »

Retenez cette séquence-type : \eng{First of all... Moreover... However... Therefore... To sum up...} Elle suffit à structurer n'importe quel paragraphe argumentatif.

\courssub{Les quatre familles de connecteurs}

\begin{propriete}[Classer les connecteurs par fonction]
Un connecteur se choisit selon la \cle{relation logique} entre deux idées, jamais au hasard. Quatre fonctions dominent dans une composition argumentative :
\begin{itemize}
\item \textbf{Addition} (ajouter une idée) : \eng{first of all} (tout d'abord), \eng{moreover} (de plus), \eng{furthermore} (en outre), \eng{in addition} (en plus), \eng{besides} (d'ailleurs).
\item \textbf{Cause et conséquence} : \eng{because} (parce que), \eng{since} (puisque), \eng{therefore} (donc), \eng{as a result} (par conséquent), \eng{consequently} (en conséquence), \eng{that is why} (c'est pourquoi).
\item \textbf{Opposition et nuance} : \eng{however} (cependant), \eng{although} (bien que), \eng{whereas} (alors que), \eng{on the other hand} (d'un autre côté).
\item \textbf{Conclusion} : \eng{to sum up} (en résumé), \eng{in conclusion} (en conclusion), \eng{all in all} (tout bien considéré).
\end{itemize}
\end{propriete}

\begin{popfigure}
\begin{tikzpicture}[
  box/.style={draw=popblue, fill=popblueL, rounded corners=2pt, align=center, font=\small, inner sep=4pt},
  head/.style={draw=popdark, fill=popdark, text=white, rounded corners=2pt, align=center, font=\small\bfseries, inner sep=4pt}]
\node[head] (titre) at (0,0) {Linking words by function};
\node[head, fill=popblue, draw=popblue] (add) at (-4.6,-1.4) {Addition};
\node[box, align=center] (addl) at (-4.6,-3.1){first of all\\moreover, furthermore\\in addition, besides};
\node[head, fill=popgreen, draw=popgreen] (cau) at (-1.5,-1.4) {Cause / Result};
\node[box, draw=popgreen, fill=popgreenL, align=center] (caul) at (-1.5,-3.1){because, since\\therefore, as a result\\consequently, that is why};
\node[head, fill=poporange, draw=poporange] (opp) at (1.6,-1.4) {Contrast};
\node[box, draw=poporange, fill=poporangeL, align=center] (oppl) at (1.6,-3.1){however, although\\whereas\\on the other hand};
\node[head, fill=poppurple, draw=poppurple] (con) at (4.7,-1.4) {Conclusion};
\node[box, draw=poppurple, fill=poppurpleL, align=center] (conl) at (4.7,-3.1){to sum up\\in conclusion\\all in all};
\draw[popline] (titre) -- (add); \draw[popline] (titre) -- (cau);
\draw[popline] (titre) -- (opp); \draw[popline] (titre) -- (con);
\draw[popline] (add) -- (addl); \draw[popline] (cau) -- (caul);
\draw[popline] (opp) -- (oppl); \draw[popline] (con) -- (conl);
\end{tikzpicture}
\legende{Carte mentale : les connecteurs classés par fonction logique}
\end{popfigure}

\begin{exemplebox}[Les connecteurs en contexte]
\begin{itemize}
\item \eng{First of all, English is essential in the tourism industry.} « Tout d'abord, l'anglais est essentiel dans l'industrie du tourisme. »
\item \eng{Moreover, speaking English increases your chances of promotion.} « De plus, parler anglais augmente vos chances de promotion. »
\item \eng{Because my company trades with Nigeria, I must improve my English.} « Parce que mon entreprise commerce avec le Nigeria, je dois améliorer mon anglais. »
\item \eng{Since English is required for the job, I am taking evening classes.} « Puisque l'anglais est exigé pour ce poste, je suis des cours du soir. »
\item \eng{I failed the interview; as a result, I decided to work harder.} « J'ai échoué à l'entretien ; par conséquent, j'ai décidé de travailler davantage. »
\item \eng{Although French is my first language, I use English every day at work.} « Bien que le français soit ma première langue, j'utilise l'anglais tous les jours au travail. »
\item \eng{My colleague speaks fluently, whereas I still hesitate.} « Mon collègue parle couramment, alors que j'hésite encore. »
\item \eng{All in all, investing in English is investing in your future.} « Tout bien considéré, investir dans l'anglais, c'est investir dans son avenir. »
\end{itemize}
\end{exemplebox}

\begin{attention}[Ponctuation et place des connecteurs]
\begin{itemize}
\item \eng{However}, \eng{moreover}, \eng{therefore} et \eng{as a result} ouvrent généralement une phrase et sont suivis d'une \textbf{virgule} : \eng{Therefore, I decided to act.} — Jamais : \eng{I was tired, however I continued.} (il faut un point-virgule ou deux phrases).
\item \eng{Because} et \eng{since} introduisent une \textbf{proposition subordonnée} : on ne peut pas les laisser seuls. Faux : \eng{Because I need English.} Correct : \eng{Because I need English, I attend evening classes.}
\item Ne traduisez pas « c'est pourquoi » par \eng{that is because} : la bonne forme est \eng{that is why}. \eng{That is why I am learning English.} « C'est pourquoi j'apprends l'anglais. »
\item \eng{Although} et \eng{but} ne se combinent \textbf{jamais} dans la même phrase. Faux : \eng{Although it is difficult, but I will succeed.} Correct : \eng{Although it is difficult, I will succeed.}
\end{itemize}
\end{attention}

\courssub{Les temps verbaux de la composition prospective}

Une composition sur l'anticipation des besoins linguistiques parle de vérités générales, de prédictions, de projets et d'hypothèses. Quatre temps suffisent, à condition de les distinguer nettement.

\begin{propriete}[Will ou be going to ? La règle du futur]
\begin{itemize}
\item \eng{Will} + base verbale : \textbf{prédiction}, opinion sur l'avenir, promesse ou décision spontanée. \eng{English will open many doors for me.} « L'anglais m'ouvrira de nombreuses portes. »
\item \eng{Be going to} + base verbale : \textbf{projet déjà décidé}, intention formée avant le moment de parole. \eng{I am going to take the TOEFL exam next year.} « Je vais passer l'examen du TOEFL l'année prochaine. »
\item \eng{Would} + base verbale : \textbf{hypothèse} (souvent après \eng{if} + prétérit). \eng{If I were fluent, I would get a promotion.} « Si j'étais bilingue, j'obtiendrais une promotion. »
\item \textbf{Present simple} : vérités générales et faits permanents. \eng{English is the language of international trade.} « L'anglais est la langue du commerce international. »
\end{itemize}
Exception à connaître : après \eng{if}, \eng{when}, \eng{as soon as}, on n'emploie jamais \eng{will} pour parler du futur : \eng{When I finish my course, I will apply for the job.} (et non \eng{when I will finish}).
\end{propriete}

\begin{center}
\begin{tabularx}{\linewidth}{|l|X|X|}
\hline
\rowcolor{popblueL} \textbf{Temps} & \textbf{Emploi} & \textbf{Exemple (traduction)} \\
\hline
Present simple & vérité générale & \eng{Employers value language skills.} « Les employeurs apprécient les compétences linguistiques. » \\
\hline
will + BV & prédiction, promesse & \eng{English will boost my career.} « L'anglais donnera un essor à ma carrière. » \\
\hline
be going to + BV & projet décidé & \eng{I am going to attend evening English classes.} « Je vais suivre des cours d'anglais du soir. » \\
\hline
would + BV & hypothèse, conditionnel & \eng{If I were fluent, I would get a promotion.} « Si j'étais bilingue, j'obtiendrais une promotion. » \\
\hline
\end{tabularx}
\end{center}

\begin{exemplebox}[Projets et hypothèses dans la rédaction]
\begin{itemize}
\item \eng{Next year, I am going to enrol in a language school in Buea.} « L'année prochaine, je vais m'inscrire dans une école de langues à Buea. »
\item \eng{I believe technology will make learning easier.} « Je crois que la technologie rendra l'apprentissage plus facile. »
\item \eng{If my employer paid for my training, I would study every weekend.} « Si mon employeur finançait ma formation, j'étudierais chaque week-end. »
\item \eng{As soon as I master English, I will apply for a better position.} « Dès que je maîtriserai l'anglais, je postulerai à un meilleur poste. »
\end{itemize}
\end{exemplebox}

\begin{attention}[Faux amis et pièges de temps]
\begin{itemize}
\item \eng{Actually} ne veut \textbf{pas} dire « actuellement » mais « en fait, en réalité ». « Actuellement » se dit \eng{currently} ou \eng{at present}. \eng{I am currently learning English.} « J'apprends actuellement l'anglais. »
\item « Je vais suivre des cours » : le français « aller + infinitif » se traduit bien par \eng{be going to}, mais attention à conjuguer \eng{be} : \eng{I \textbf{am} going to}, \eng{she \textbf{is} going to}, \eng{they \textbf{are} going to}.
\item Pas de futur après \eng{if} : \eng{If I \textbf{have} time, I \textbf{will} practise.} (jamais \eng{if I will have}).
\item \eng{To assist} signifie « aider », pas « assister à ». « Assister à un cours » se dit \eng{to attend a course}.
\end{itemize}
\end{attention}

\courssub{Le vocabulaire thématique}

Trois champs lexicaux sont indispensables pour ce sujet : le monde du travail, les compétences linguistiques et les stratégies d'apprentissage.

\begin{definition}[Work vocabulary — le monde du travail]
\begin{itemize}
\item \eng{employer} (n.) : employeur ; \eng{employee} (n.) : employé(e) ; \eng{colleague} (n.) : collègue.
\item \eng{career} (n.) : carrière ; \eng{promotion} (n.) : promotion, avancement.
\item \eng{job interview} (n.) : entretien d'embauche ; \eng{to apply for a job} (v.) : postuler à un emploi.
\end{itemize}
\end{definition}

\begin{definition}[Language skills — les compétences linguistiques]
\begin{itemize}
\item \eng{fluent} (adj.) : courant (à l'oral) ; \eng{fluency} (n.) : aisance. \eng{She is fluent in English.}
\item \eng{proficiency} (n.) : maîtrise, compétence ; \eng{language skills} (n.) : compétences linguistiques.
\item \eng{to improve} (v.) : améliorer ; \eng{to master} (v.) : maîtriser ; \eng{communication skills} (n.) : compétences en communication.
\end{itemize}
\end{definition}

\begin{definition}[Learning strategies — les stratégies d'apprentissage]
\begin{itemize}
\item \eng{to attend courses} : suivre des cours ; \eng{to practise daily} : pratiquer quotidiennement.
\item \eng{to watch English programmes} : regarder des émissions en anglais ; \eng{to read professional documents} : lire des documents professionnels.
\end{itemize}
\end{definition}

\begin{center}
\begin{tabularx}{\linewidth}{|X|X|X|}
\hline
\rowcolor{popblueL} \textbf{Français} & \textbf{English} & \textbf{Remarque} \\
\hline
postuler à un emploi & to apply for a job & préposition \eng{for} obligatoire \\
\hline
suivre des cours & to attend courses & pas \eng{to assist} (faux ami) \\
\hline
courant (en anglais) & fluent (in English) & préposition \eng{in} \\
\hline
améliorer son niveau & to improve one's level & pas de réflexif en anglais \\
\hline
un entretien d'embauche & a job interview & ordre des mots inversé \\
\hline
actuellement & currently & pas \eng{actually} (= en fait) \\
\hline
\end{tabularx}
\end{center}

\begin{exoresolu}[Connecteurs et temps dans un paragraphe]
Énoncé — Complétez ce paragraphe avec les mots suivants : \eng{first of all}, \eng{however}, \eng{therefore}, \eng{am going to}, \eng{would}, \eng{to sum up}.

\eng{\ldots\ldots, English is essential for my career because my company exports cocoa to Ghana. \ldots\ldots, my level is still too low for international meetings. \ldots\ldots, I \ldots\ldots attend evening classes in Yaoundé. If I were fluent, I \ldots\ldots ask for a promotion. \ldots\ldots, learning English now is the best investment for my future.}

\tcblower
Solution détaillée :
\begin{itemize}
\item \eng{\textbf{First of all}, English is essential...} : on ouvre l'argumentation par le connecteur d'introduction d'idée.
\item \eng{\textbf{However}, my level is still too low...} : opposition entre l'importance de l'anglais et le niveau actuel.
\item \eng{\textbf{Therefore}, I \textbf{am going to} attend evening classes...} : conséquence logique ; \eng{be going to} car il s'agit d'un projet déjà décidé (pas \eng{will}).
\item \eng{If I were fluent, I \textbf{would} ask for a promotion.} : hypothèse ; structure \eng{if} + prétérit, \eng{would} + base verbale.
\item \eng{\textbf{To sum up}, learning English now is the best investment...} : connecteur de conclusion.
\end{itemize}
Traduction : « Tout d'abord, l'anglais est essentiel pour ma carrière car mon entreprise exporte du cacao vers le Ghana. Cependant, mon niveau est encore trop faible pour les réunions internationales. C'est pourquoi je vais suivre des cours du soir à Yaoundé. Si j'étais bilingue, je demanderais une promotion. En résumé, apprendre l'anglais maintenant est le meilleur investissement pour mon avenir. »
\end{exoresolu}

\begin{methode}[Bien utiliser cette boîte à outils dans votre composition]
\begin{enumerate}
\item Avant de rédiger, choisissez \textbf{un connecteur par paragraphe} : un pour ajouter (\eng{moreover}), un pour opposer (\eng{however}), un pour conclure (\eng{as a result}).
\item Réservez \eng{will} aux prédictions générales et \eng{be going to} à vos projets personnels : le correcteur vérifie cette distinction.
\item Placez au moins une phrase hypothétique avec \eng{if ... would} : c'est la structure qui rapporte le plus de points au Bac.
\item Glissez cinq à huit mots du vocabulaire thématique (\eng{career}, \eng{promotion}, \eng{fluent}, \eng{to apply for}, \eng{to attend courses}) : ils prouvent que vous maîtrisez le sujet.
\item Relisez en vérifiant la ponctuation après chaque connecteur et l'accord de \eng{be} dans \eng{be going to}.
\end{enumerate}
\end{methode}

\begin{exercice}[À vous de jouer]
\begin{enumerate}
\item Traduisez : « Je vais postuler à un nouvel emploi parce que je veux améliorer ma carrière. »
\item Choisissez le bon connecteur : \eng{My English is weak. (Moreover / However), I am going to practise daily.}
\item Corrigez la faute : \eng{If I will be fluent, I would get the job.}
\item Rédigez trois phrases sur vos propres projets linguistiques avec \eng{be going to}, \eng{will} et \eng{would}.
\end{enumerate}
\end{exercice}

\begin{corrige}[Exercice « À vous de jouer »]
\begin{enumerate}
\item \eng{I am going to apply for a new job because I want to improve my career.} (projet décidé : \eng{be going to} ; \eng{apply for} avec \eng{for}).
\item \eng{However} : il y a opposition entre la faiblesse actuelle et la décision de pratiquer.
\item \eng{If I \textbf{were} fluent, I would get the job.} : jamais de \eng{will} après \eng{if} dans une hypothèse.
\item Modèle : \eng{I am going to watch English programmes every evening. English will help me at work. If I had more time, I would attend courses in Douala.}
\end{enumerate}
\end{corrige}

Vous disposez désormais des connecteurs, des temps et du lexique nécessaires. La section suivante mettra ces outils en action dans un modèle de composition entièrement commenté.

\coursec{Modèle commenté}

Dans la section précédente, nous avons constitué notre « boîte à outils » : connecteurs logiques (\eng{first of all}, \eng{moreover}, \eng{to sum up}…), temps verbaux adaptés au projet professionnel (\eng{will}, \eng{be going to}, présent de vérité générale) et vocabulaire du monde du travail. Il est temps maintenant de voir comment un bon candidat assemble tous ces éléments dans une composition réelle. Nous allons lire un modèle complet, puis le disséquer paragraphe par paragraphe, comme un mécanicien démonte un moteur pour comprendre comment il fonctionne.

\courssub{Lecture du modèle}

\begin{textebox}[Model composition — “My English Needs at Work” (about 280 words)]
English is becoming the language of business worldwide, and Cameroon is no exception. As a future accountant, I know that my success will depend partly on my English skills. This essay will discuss the English I will need at work and how I plan to acquire it.

First of all, I will need English to communicate with foreign clients and partners. Many companies in Douala work with Nigerian and Ghanaian firms; therefore, speaking fluent English will help me negotiate and write professional emails. Moreover, most accounting software and financial reports are written in English, so I must be able to understand technical documents.

To meet these needs, I am going to attend evening classes and practise English every day by watching the news and reading business articles. In addition, I intend to take an international English test before graduation.

To sum up, English is not just a school subject for me: it is a professional tool. If I invest in my English now, I will certainly open many doors in my future career.
\end{textebox}

\textbf{Glossaire}

\begin{center}
\begin{tabularx}{\linewidth}{|l|l|X|}
\hline
\rowcolor{popblueL} \textbf{English} & \textbf{Nature} & \textbf{Français} \\
\hline
worldwide & adverb & dans le monde entier \\
\hline
no exception & expression & ne fait pas exception \\
\hline
accountant & noun & comptable \\
\hline
skills & noun (pl.) & compétences \\
\hline
to acquire & verb & acquérir \\
\hline
foreign clients & noun phrase & clients étrangers \\
\hline
firm & noun & entreprise, cabinet \\
\hline
fluent & adjective & courant (parler un anglais courant) \\
\hline
to negotiate & verb & négocier \\
\hline
accounting software & noun phrase & logiciel de comptabilité \\
\hline
to meet these needs & verb phrase & répondre à ces besoins \\
\hline
evening classes & noun (pl.) & cours du soir \\
\hline
graduation & noun & obtention du diplôme \\
\hline
to open doors & idiom & ouvrir des portes, créer des opportunités \\
\hline
\end{tabularx}
\end{center}

\textbf{Questions de compréhension (reading check)}

\begin{enumerate}
\item What is the writer's future profession?
\item Why does the writer mention Douala?
\item What are the two main English needs identified in the essay?
\item What three actions does the writer plan to take?
\item What is the writer's personal opinion about English?
\end{enumerate}

\begin{corrige}[Questions de compréhension]
\begin{enumerate}
\item The writer wants to become an accountant.
\item Because many companies in Douala work with Nigerian and Ghanaian firms — a concrete, local example.
\item Communicating with foreign clients and partners, and understanding technical documents (software, financial reports).
\item Attend evening classes, practise English every day (news, business articles), and take an international English test.
\item English is not just a school subject: it is a professional tool that will open doors in his or her career.
\end{enumerate}
\end{corrige}

\courssub{Analyse de l'introduction : hook et thesis statement}

Relisons le premier paragraphe phrase par phrase.

\begin{exemplebox}[Découpage de l'introduction]
\begin{itemize}
\item \textbf{Phrase 1 :} \eng{English is becoming the language of business worldwide, and Cameroon is no exception.} → c'est le \cle{hook} (l'accroche) : une affirmation générale et actuelle qui capte l'attention du lecteur.
\item \textbf{Phrase 2 :} \eng{As a future accountant, I know that my success will depend partly on my English skills.} → phrase de \cle{contexte personnel} : le candidat se présente et relie le sujet à sa propre vie.
\item \textbf{Phrase 3 :} \eng{This essay will discuss the English I will need at work and how I plan to acquire it.} → c'est la \cle{thesis statement} (la phrase d'annonce du plan) : elle annonce les deux parties du développement — (1) les besoins en anglais, (2) la stratégie pour les acquérir.
\end{itemize}
\end{exemplebox}

\begin{propriete}[La thesis statement]
La \cle{thesis statement} est la phrase la plus importante de l'introduction. Elle se place généralement en \textbf{dernière position} du premier paragraphe et annonce clairement le contenu du développement. Formule type : \eng{This essay will discuss / examine / argue...} suivi des deux ou trois idées principales. Elle fonctionne comme un contrat avec le lecteur : tout ce qui suit doit correspondre à ce qui est annoncé.
\end{propriete}

\begin{attention}[Piège fréquent des francophones]
Ne traduis pas mot à mot « Cette dissertation va parler de... » par \eng{This dissertation will talk about...}. En anglais, on dit \eng{essay} (pas \eng{dissertation}, qui désigne une thèse universitaire) et le verbe \eng{to talk about} est trop oral ; préfère \eng{to discuss}, \eng{to examine} ou \eng{to deal with}. Écris : \eng{This essay will discuss my English needs at work.} « Cette rédaction examinera mes besoins en anglais au travail. »
\end{attention}

\courssub{Analyse du corps : topic sentences et exemples}

Chaque paragraphe du corps s'ouvre sur une \cle{topic sentence} (phrase d'idée principale), suivie d'explications et d'exemples concrets.

\begin{center}
\begin{tabularx}{\linewidth}{|l|X|X|}
\hline
\rowcolor{popblueL} \textbf{Paragraphe} & \textbf{Topic sentence} & \textbf{Exemples / preuves} \\
\hline
Corps 1 & \eng{I will need English to communicate with foreign clients and partners.} & companies in Douala, Nigerian and Ghanaian firms, negotiating, professional emails \\
\hline
Corps 1 (suite) & \eng{Moreover, most accounting software and financial reports are written in English.} & technical documents \\
\hline
Corps 2 & \eng{To meet these needs, I am going to attend evening classes...} & watching the news, reading business articles, taking an international test \\
\hline
\end{tabularx}
\end{center}

\begin{aretenir}[La règle d'or du paragraphe]
Un paragraphe = \textbf{une idée} (topic sentence) + \textbf{une explication} + \textbf{un exemple concret}. Le modèle illustre parfaitement cette règle : l'idée « communiquer avec des clients étrangers » est immédiatement rendue crédible par l'exemple local des entreprises de Douala qui travaillent avec le Nigeria et le Ghana.
\end{aretenir}

Remarque la progression logique : le corps 1 décrit les \textbf{besoins} (\eng{I will need...}), le corps 2 présente les \textbf{solutions} (\eng{I am going to...}). La charnière \eng{To meet these needs} (« Pour répondre à ces besoins ») relie élégamment les deux parties.

\courssub{Analyse de la conclusion : rappel et opinion personnelle}

\begin{exemplebox}[Découpage de la conclusion]
\begin{itemize}
\item \textbf{Rappel de la thèse (reformulé) :} \eng{To sum up, English is not just a school subject for me: it is a professional tool.} — « En résumé, l'anglais n'est pas seulement une matière scolaire pour moi : c'est un outil professionnel. »
\item \textbf{Opinion personnelle + ouverture :} \eng{If I invest in my English now, I will certainly open many doors in my future career.} — « Si j'investis dans mon anglais maintenant, j'ouvrirai certainement de nombreuses portes dans ma future carrière. »
\end{itemize}
\end{exemplebox}

\begin{propriete}[La bonne conclusion]
Une conclusion efficace : (1) commence par un connecteur de bilan (\eng{to sum up}, \eng{in conclusion}, \eng{all things considered}) ; (2) \textbf{reformule} l'idée principale sans la recopier mot pour mot ; (3) ajoute une opinion personnelle ou une ouverture vers l'avenir. Elle n'introduit \textbf{jamais} d'idée nouvelle.
\end{propriete}

\courssub{Repérage des connecteurs et des temps verbaux}

\begin{center}
\begin{tabularx}{\linewidth}{|l|X|l|}
\hline
\rowcolor{popblueL} \textbf{Connecteur} & \textbf{Fonction} & \textbf{Position dans le modèle} \\
\hline
\eng{first of all} & introduire le premier argument & début du corps 1 \\
\hline
\eng{therefore} & conséquence logique & milieu du corps 1 \\
\hline
\eng{moreover} & ajout d'un argument & corps 1 \\
\hline
\eng{so} & conséquence & corps 1 \\
\hline
\eng{to meet these needs} & transition besoins → solutions & début du corps 2 \\
\hline
\eng{in addition} & ajout & corps 2 \\
\hline
\eng{to sum up} & bilan & début de la conclusion \\
\hline
\end{tabularx}
\end{center}

\begin{center}
\begin{tabularx}{\linewidth}{|l|X|X|}
\hline
\rowcolor{popblueL} \textbf{Temps / forme} & \textbf{Exemple dans le modèle} & \textbf{Emploi} \\
\hline
present continuous & \eng{English is becoming...} & tendance en cours \\
\hline
simple present & \eng{most reports are written in English} & vérité générale \\
\hline
\eng{will} + BV & \eng{my success will depend...} & prédiction \\
\hline
\eng{be going to} + BV & \eng{I am going to attend evening classes} & projet décidé \\
\hline
\eng{intend to} + BV & \eng{I intend to take a test} & intention formelle \\
\hline
1\textsuperscript{re} conditionnelle & \eng{If I invest..., I will open...} & hypothèse réelle \\
\hline
\end{tabularx}
\end{center}

\begin{attention}[Concordance des temps]
Dans la phrase \eng{If I invest in my English now, I will certainly open many doors}, observe bien : \textbf{if + présent}, puis \textbf{will} dans la proposition principale. N'écris jamais \eng{If I will invest...} — c'est l'erreur de conditionnel la plus fréquente chez les francophones, par calque du français « si j'investirai » (incorrect aussi en français d'ailleurs : « si j'investis »).
\end{attention}

\courssub{Schéma annoté du modèle}

\begin{popfigure}
\begin{center}
\begin{tikzpicture}[
  box/.style={draw=popblue, thick, rounded corners, fill=popblueL, text width=4.6cm, align=center, font=\small, inner sep=6pt},
  func/.style={draw=poporange, thick, rounded corners, fill=poporangeL, text width=3.4cm, align=center, font=\small, inner sep=5pt},
  fl/.style={-{Stealth[length=3mm]}, thick, popdark}
]
\node[box, align=center] (intro) at (0,0){\textbf{Introduction}\\ “English is becoming the language of business...”};
\node[box, align=center] (body1) at (0,-2.6){\textbf{Body 1}\\ “First of all, I will need English to communicate...”};
\node[box, align=center] (body2) at (0,-5.2){\textbf{Body 2}\\ “To meet these needs, I am going to attend...”};
\node[box, align=center] (concl) at (0,-7.8){\textbf{Conclusion}\\ “To sum up, English is not just a school subject...”};
\node[func] (f1) at (6.4,0) {Hook + thesis statement (annonce du plan)};
\node[func] (f2) at (6.4,-2.6) {Argument 1 : les besoins (clients, documents) + exemples locaux};
\node[func] (f3) at (6.4,-5.2) {Argument 2 : les solutions (cours, pratique, test)};
\node[func] (f4) at (6.4,-7.8) {Rappel reformulé + opinion personnelle};
\draw[fl] (intro) -- (f1);
\draw[fl] (body1) -- (f2);
\draw[fl] (body2) -- (f3);
\draw[fl] (concl) -- (f4);
\draw[fl, popgreen] (intro.south) -- (body1.north);
\draw[fl, popgreen] (body1.south) -- (body2.north);
\draw[fl, popgreen] (body2.south) -- (concl.north);
\end{tikzpicture}
\end{center}
\legende{Chaque paragraphe du modèle remplit une fonction précise ; les flèches vertes montrent la progression logique de l'ensemble.}
\end{popfigure}

\courssub{Activité de soulignage collectif}

\begin{experience}[Highlight the model!]
En classe entière ou en binômes, reprends le texte du modèle et utilise trois couleurs :
\begin{enumerate}
\item \textbf{Surligne en jaune} tous les connecteurs (\eng{first of all}, \eng{therefore}, \eng{moreover}, \eng{so}, \eng{in addition}, \eng{to sum up}).
\item \textbf{Surligne en vert} tous les marqueurs de futur et de projet (\eng{will depend}, \eng{will need}, \eng{will help}, \eng{am going to attend}, \eng{intend to take}, \eng{will open}).
\item \textbf{Surligne en bleu} le vocabulaire thématique du travail (\eng{business}, \eng{accountant}, \eng{clients}, \eng{negotiate}, \eng{software}, \eng{financial reports}, \eng{career}).
\end{enumerate}
Compare ensuite ton résultat avec celui de ton voisin : avez-vous repéré les mêmes éléments ? Compte-les : le modèle contient sept connecteurs, six marqueurs de futur et plus de dix mots du champ lexical professionnel.
\end{experience}

\begin{methode}[Comment exploiter un modèle sans le recopier]
\begin{enumerate}
\item \textbf{Lis} le modèle une première fois pour le sens global, puis une deuxième fois en repérant la structure (hook, thesis, topic sentences, conclusion).
\item \textbf{Extrais} le « squelette » réutilisable : les connecteurs, la formule de thesis statement (\eng{This essay will discuss...}), le schéma besoins → solutions.
\item \textbf{Remplace} le contenu par tes propres idées : ton futur métier (infirmier, ingénieur, commerçant à Bamenda, hôtelier à Kribi...), tes besoins précis, ton plan d'action.
\item \textbf{Vérifie} que ta composition garde la même architecture : introduction en trois temps, deux paragraphes de corps avec exemples locaux, conclusion avec opinion personnelle.
\item \textbf{Ne recopie jamais} le modèle mot pour mot : au Bac, les correcteurs reconnaissent immédiatement les textes appris par cœur et les pénalisent.
\end{enumerate}
\end{methode}

\begin{exoresolu}[Réutiliser la structure du modèle]
Énoncé : en t'inspirant de la structure du modèle, rédige la thesis statement et la topic sentence du premier paragraphe du corps pour un futur \textbf{infirmier} (nurse) à Buea.
\tcblower
Solution détaillée. On conserve le squelette du modèle et on change le contenu :
\begin{itemize}
\item Thesis statement : \eng{This essay will discuss the English I will need as a nurse and how I plan to improve it.} (« Cette rédaction examinera l'anglais dont j'aurai besoin comme infirmier et comment je compte le perfectionner. ») — on réutilise la formule \eng{This essay will discuss... and how I plan to...}.
\item Topic sentence : \eng{First of all, I will need English to communicate with patients and doctors from English-speaking regions.} (« Tout d'abord, j'aurai besoin de l'anglais pour communiquer avec des patients et des médecins des régions anglophones. ») — on réutilise \eng{First of all, I will need English to...} avec un exemple ancré dans le contexte camerounais (Buea, ville bilingue).
\end{itemize}
La structure est identique ; les idées sont personnelles. C'est exactement ce qu'attend le correcteur.
\end{exoresolu}

\begin{savaistu}[Ce que les correcteurs du Bac récompensent]
Les correcteurs du baccalauréat camerounais lisent des centaines de copies. Ce qui fait la différence ? Les compositions \textbf{personnelles et ancrées dans le réel} : mentionner les entreprises de Douala qui commercent avec le Nigeria, les hôpitaux de Buea, les ONG de Bamenda, le port de Kribi ou le marché de cacao montre que tu sais relier l'anglais à ta vie concrète. À l'inverse, les textes génériques appris par cœur, sans aucun exemple local, obtiennent rarement la moyenne en expression écrite. Ton expérience camerounaise est un atout : utilise-la !
\end{savaistu}

\begin{aretenir}[L'essentiel du modèle]
Le modèle « My English Needs at Work » illustre la composition idéale : une introduction en trois temps (hook + contexte personnel + thesis statement), deux paragraphes de corps construits sur le schéma besoins → solutions avec des exemples locaux, et une conclusion qui reformule la thèse avant de donner une opinion personnelle. Retiens le squelette, pas les phrases : c'est ta capacité à le remplir avec tes propres idées qui sera évaluée.
\end{aretenir}

\coursec{Erreurs à éviter et relecture}

Dans la section précédente, nous avons lu et commenté un modèle de composition sur le thème \eng{Anticipating my second language needs at work}. Vous savez désormais à quoi ressemble une bonne copie : une introduction qui annonce le plan, des paragraphes construits autour d'une \eng{topic sentence}, une conclusion qui ouvre sur l'avenir. Mais entre le modèle et votre propre production, il reste un obstacle de taille : les erreurs. Or, au Baccalauréat, une composition pleine de bonnes idées mais truffée de fautes de grammaire, de faux amis et de majuscules oubliées perd une grande partie de ses points. Cette section vous donne les outils pour repérer, corriger et surtout éviter les erreurs les plus fréquentes des candidats francophones, puis une méthode de relecture en trois passes, simple et redoutablement efficace.

\courssub{Les erreurs de structure}

\begin{definition}[Les trois erreurs de structure à bannir]
Une composition peut être grammaticalement correcte et pourtant médiocre si sa \cle{structure} est défaillante. Trois fautes de construction reviennent sans cesse dans les copies :
\begin{enumerate}
\item \textbf{L'absence de plan} : le candidat écrit au fil de la plume, sans introduction ni conclusion. Le correcteur ne voit aucune progression.
\item \textbf{Des paragraphes sans \eng{topic sentence}} : chaque paragraphe doit s'ouvrir sur une phrase qui annonce l'idée principale. Sans elle, le paragraphe devient une liste de phrases décousues.
\item \textbf{Une conclusion copiée de l'introduction} : répéter mot pour mot sa première phrase en conclusion est une faute de méthode. La conclusion doit être un \emph{bilan suivi d'une ouverture}, pas une \emph{répétition}.
\end{enumerate}
\end{definition}

\begin{exemplebox}[Topic sentence : avant / après]
\textbf{Paragraphe faible} (sans phrase d'annonce) :\\
\eng{English is used in many companies. I want to work in a bank. My sister speaks English well.} — Trois phrases vraies, mais aucun fil conducteur.

\textbf{Paragraphe fort} (avec \eng{topic sentence}) :\\
\eng{First of all, English will be essential in my future career as a bank clerk. Most international banks in Douala use English in their daily communication, and customers often write their complaints in English. That is why I must improve my level before I start working.} — « Tout d'abord, l'anglais sera indispensable dans ma future carrière d'employé de banque. La plupart des banques internationales de Douala utilisent l'anglais dans leur communication quotidienne, et les clients rédigent souvent leurs réclamations en anglais. C'est pourquoi je dois améliorer mon niveau avant de commencer à travailler. » La première phrase annonce l'idée ; les suivantes la développent.
\end{exemplebox}

\begin{attention}[La conclusion n'est pas une photocopie]
Introduction : \eng{In this composition, I will explain how I plan to improve my English for my future job.}\\
Conclusion fautive : \eng{In this composition, I will explain how I plan to improve my English for my future job.} (copiée-collée !)\\
Conclusion correcte : \eng{To sum up, by taking evening courses and practising every day, I am confident that I will master English before I enter the job market.} — « Pour conclure, en suivant des cours du soir et en m'entraînant chaque jour, je suis sûr que je maîtriserai l'anglais avant d'entrer dans la vie active. »
\end{attention}

\courssub{Les erreurs de grammaire fréquentes}

\courssub{Observation d'abord : repérez la faute}

Lisez ces quatre phrases extraites de copies réelles type. Chacune contient exactement une erreur. Saurez-vous la trouver avant de lire la correction ?

\begin{enumerate}
\item \eng{My brother work in a travel agency in Buea.}
\item \eng{I will to attend English courses next year.}
\item \eng{I am agree with this idea.}
\item \eng{Next year, I am going to improving my english.}
\end{enumerate}

Corrections : (1) \eng{My brother \textbf{works}...} — oubli du \textbf{-s} de la 3\textsuperscript{e} personne ; (2) \eng{I will attend...} — jamais de \eng{to} après \eng{will} ; (3) \eng{I agree...} — \eng{agree} est un verbe, pas un adjectif ; (4) \eng{...going to improve my \textbf{English}} — infinitif sans \eng{-ing} après \eng{be going to}, et majuscule obligatoire à \eng{English}.

\begin{propriete}[Règle 1 : le « s » de la 3e personne du singulier]
Au présent simple, toute la conjugaison est identique \textbf{sauf} à la 3\textsuperscript{e} personne du singulier (\eng{he, she, it}), qui prend un \textbf{-s} (ou \textbf{-es}) : \eng{I work} mais \eng{she work\textbf{s}} ; \eng{they improve} mais \eng{he improve\textbf{s}}.
\begin{itemize}
\item \eng{My father works in Yaoundé.} — « Mon père travaille à Yaoundé. »
\item \eng{English helps me every day.} — « L'anglais m'aide chaque jour. »
\end{itemize}
\textbf{Exception de forme} : les verbes en \eng{-ss, -sh, -ch, -x, -o} prennent \textbf{-es} : \eng{she watches, he goes, it finishes}. Les verbes en « consonne + \eng{-y} » transforment le \eng{y} en \eng{-ies} : \eng{she studies, he carries} (mais \eng{he plays}, car \eng{-ay} contient une voyelle).
\textbf{Attention} : le \eng{-s} disparaît avec les auxiliaires : \eng{she \textbf{does} not work}, \eng{does she work?}
\end{propriete}

\begin{propriete}[Règle 2 : will ou be going to ?]
Les deux expriment le futur, mais pas le même :
\begin{itemize}
\item \eng{\textbf{will} + base verbale} : décision spontanée, promesse, prédiction. \eng{I will work hard, I promise.} — « Je travaillerai dur, je te le promets. »
\item \eng{\textbf{be going to} + base verbale} : projet déjà décidé, intention planifiée. \eng{I am going to take evening classes next term.} — « Je vais suivre des cours du soir le trimestre prochain » (projet mûri).
\end{itemize}
Dans une composition sur l'anticipation de ses besoins en anglais, \eng{be going to} domine naturellement, car vous décrivez des \emph{projets}. \textbf{Jamais de \eng{to} après \eng{will}} : \eng{I will attend} (et non \eng{I will to attend}). Et après \eng{be going to}, on met la \textbf{base verbale}, jamais la forme en \eng{-ing} : \eng{I am going to improve} (et non \eng{I am going to improving}).
\end{propriete}

\begin{exemplebox}[Exemples traduits : la phrase fautive et sa correction]
\begin{itemize}
\item Fautif : \eng{I will to attend courses.} Correct : \eng{I will attend courses.} — « Je suivrai des cours. » (pas de \eng{to} après \eng{will})
\item Fautif : \eng{I am going to improving my English.} Correct : \eng{I am going to improve my English.} — « Je vais améliorer mon anglais. » (base verbale, sans \eng{-ing})
\item Fautif : \eng{He don't speak English.} Correct : \eng{He doesn't speak English.} — « Il ne parle pas anglais. »
\item Fautif : \eng{My sister want to become a nurse.} Correct : \eng{My sister wants to become a nurse.} — « Ma sœur veut devenir infirmière. »
\end{itemize}
\end{exemplebox}

\begin{attention}[Erreur fréquente des francophones : « I am agree »]
En français, on dit « je suis d'accord », avec le verbe \emph{être}. En anglais, \eng{agree} est un \textbf{verbe}, pas un adjectif : on dit donc \eng{I agree} — « je suis d'accord » — et \eng{I disagree} — « je ne suis pas d'accord ». La forme \eng{I am agree} n'existe pas. Même vigilance avec \eng{I am think} : on dit \eng{I think}.
\end{attention}

\begin{attention}[Majuscules et réflexifs fantômes]
Deux pièges réunis dans une seule phrase fautive : \eng{I am going to improve me english.}
\begin{enumerate}
\item \eng{English} prend \textbf{toujours} une majuscule, comme toutes les langues, nationalités et noms de pays : \eng{English, French, Cameroon, Cameroonian}. On écrit donc \eng{my English}, jamais \eng{my english}.
\item Le français « m'améliorer » tente les élèves d'ajouter un réflexif : on dit \eng{to improve \textbf{my English}} (améliorer \emph{mon anglais}), jamais \eng{to improve me}.
\end{enumerate}
Phrase modèle : \eng{I am going to improve my English before I leave school.} — « Je vais améliorer mon anglais avant de quitter le lycée. »
\end{attention}

\courssub{Les erreurs de vocabulaire : faux amis et traductions mot à mot}

\begin{definition}[Faux ami]
Un \cle{faux ami} (\eng{false friend}) est un mot anglais qui ressemble à un mot français mais a un sens différent. Le candidat pressé traduit « actuellement » par \eng{actually}... et dit en réalité « en fait » !
\end{definition}

\begin{center}
\begin{tabularx}{\linewidth}{|X|X|X|}
\hline
\rowcolor{popblueL} Français visé & English correct & Piège (faux ami) \\
\hline
actuellement & currently, at present & \eng{actually} = en fait, en réalité \\
\hline
en fait & actually, in fact & \eng{in fact} est correct ici \\
\hline
assister à (un cours) & to attend (a course) & \eng{to assist} = aider, secourir \\
\hline
aider & to help, to assist & \eng{to help} est le plus courant \\
\hline
une formation & training, a course & \eng{a formation} n'existe pas en anglais \\
\hline
améliorer & to improve & pas de réflexif : \eng{improve my English} \\
\hline
un emploi / un travail & a job, work & \eng{a work} n'existe pas (\eng{work} est indénombrable) \\
\hline
avoir besoin de & to need & pas \eng{to have need of} (littéraire) \\
\hline
réussir (un examen) & to pass (an exam) & \eng{to succeed} = réussir en général, suivi de \eng{in} \\
\hline
\end{tabularx}
\end{center}

\begin{exemplebox}[Traductions mot à mot à proscrire]
\begin{itemize}
\item « J'assiste à des cours du soir » : fautif \eng{I assist evening courses} ; correct \eng{I attend evening classes.}
\item « Actuellement, je travaille mon vocabulaire » : fautif \eng{Actually, I work my vocabulary} ; correct \eng{Currently, I am working on my vocabulary.}
\item « J'aurai besoin de l'anglais dans mon travail » : fautif \eng{I will have need of English in my work} ; correct \eng{I will need English in my job.}
\end{itemize}
\end{exemplebox}

\courssub{Ponctuation et majuscules : les détails qui coûtent cher}

\begin{propriete}[Majuscules obligatoires en anglais]
Contrairement au français, l'anglais exige la majuscule pour :
\begin{itemize}
\item le pronom \eng{I} (« je »), \textbf{toujours}, même au milieu d'une phrase : \eng{My brother and I study English.} ;
\item les langues et nationalités : \eng{English, French, Cameroonian} ;
\item les noms de pays et de villes : \eng{Cameroon, Nigeria, Douala, Bamenda} ;
\item les jours et les mois : \eng{Monday, January} (alors que le français écrit « lundi, janvier » en minuscules !).
\end{itemize}
Côté ponctuation : en anglais, \textbf{pas d'espace avant} les deux-points, le point-virgule, le point d'exclamation ou d'interrogation : \eng{I need English: it is the language of business!}
\end{propriete}

\begin{center}
\begin{tabularx}{\linewidth}{|X|X|X|}
\hline
\rowcolor{popblueL} Français & English & Remarque \\
\hline
je parle anglais & I speak English & \eng{I} et \eng{English} : majuscules \\
\hline
le lundi, en janvier & on Monday, in January & majuscules aux jours et mois \\
\hline
tous les lundis & every Monday, on Mondays & les deux formes sont correctes \\
\hline
le camerounais & Cameroonian (a Cameroonian) & nationalité : majuscule \\
\hline
j'habite à Yaoundé ! & I live in Yaoundé! & pas d'espace avant « ! » \\
\hline
\end{tabularx}
\end{center}

\courssub{La relecture en trois passes}

Relire, ce n'est pas « jeter un coup d'œil ». Une relecture efficace se fait en \textbf{trois passes séparées}, chacune avec un objectif unique. Vouloir tout corriger en même temps, c'est le meilleur moyen de ne rien voir.

\begin{popfigure}
\begin{tikzpicture}[font=\small]
\node[draw=popblue, fill=popblueL, rounded corners, align=center, minimum width=3.6cm, minimum height=1.6cm] (p1) at (0,0) {\textbf{Passage 1 : idées}\\ contenu, plan,\\ topic sentences};
\node[draw=poporange, fill=poporangeL, rounded corners, align=center, minimum width=3.6cm, minimum height=1.6cm] (p2) at (5.2,0) {\textbf{Passage 2 : grammaire}\\ temps, -s de la 3e pers.,\\ will / be going to};
\node[draw=popgreen, fill=popgreenL, rounded corners, align=center, minimum width=3.6cm, minimum height=1.6cm] (p3) at (10.4,0) {\textbf{Passage 3 : forme}\\ orthographe,\\ majuscules, ponctuation};
\draw[-{Stealth}, very thick, popdark] (p1) -- (p2);
\draw[-{Stealth}, very thick, popdark] (p2) -- (p3);
\node[align=center, popmuted] at (5.2,-1.6) {Une seule cible par passage : le cerveau voit enfin les fautes.};
\end{tikzpicture}
\legende{Frise de la relecture en trois passes : des idées vers la forme.}
\end{popfigure}

\begin{methode}[Relire efficacement sa composition]
\begin{enumerate}
\item \textbf{Passage 1 — Contenu et organisation.} Vérifiez : ai-je une introduction, deux ou trois paragraphes avec chacun une \eng{topic sentence}, une conclusion qui ne copie pas l'introduction ? Ai-je vraiment répondu au sujet ?
\item \textbf{Passage 2 — Grammaire et temps.} Soulignez chaque verbe : le \eng{-s} de la 3\textsuperscript{e} personne est-il là ? \eng{will} est-il suivi de la base verbale sans \eng{to} ? \eng{be going to} est-il suivi d'un infinitif ?
\item \textbf{Passage 3 — Orthographe et ponctuation.} Traquez les majuscules (\eng{I, English, Cameroon}), les faux amis (\eng{actually}, \eng{assist}) et les espaces avant la ponctuation.
\item \textbf{Astuce finale} : lisez votre composition \emph{à voix basse}. Toute phrase qui « sonne faux » à l'oreille doit être simplifiée : une phrase courte et correcte vaut toujours mieux qu'une phrase longue et bancale.
\end{enumerate}
\end{methode}

\courssub{Auto-correction guidée : les symboles du correcteur}

Pour s'entraîner, corrigez les copies de vos camarades (et les vôtres) avec les symboles officiels : \textbf{Gr} = grammaire, \textbf{V} = vocabulaire, \textbf{Sp} = orthographe (\eng{spelling}), \textbf{P} = ponctuation, \textbf{WO} = ordre des mots (\eng{word order}). Inscrire le symbole dans la marge sans donner la correction oblige l'auteur à trouver lui-même la faute : c'est ainsi qu'on progresse.

\begin{exoresolu}[Copie fautive à corriger avec les symboles]
Énoncé — Voici un extrait de composition. Annotez chaque faute avec le symbole adapté (Gr, V, Sp, P), puis réécrivez le passage correctement :\\
\eng{Actually, i am going to assist evening courses because I will to improve me english. My brother work in Douala and he is agree : english is very important for find a job.}
\tcblower
Solution détaillée :
\begin{itemize}
\item \eng{Actually} → \textbf{V} : faux ami ; on veut dire « actuellement » → \eng{Currently}.
\item \eng{i} → \textbf{Sp} : le pronom \eng{I} prend toujours une majuscule.
\item \eng{assist evening courses} → \textbf{V} : \eng{to assist} = aider ; « assister à » = \eng{to attend}.
\item \eng{I will to improve} → \textbf{Gr} : jamais de \eng{to} après \eng{will} → \eng{I will improve}.
\item \eng{improve me english} → \textbf{Gr + Sp} : pas de réflexif, et majuscule → \eng{improve my English}.
\item \eng{My brother work} → \textbf{Gr} : 3\textsuperscript{e} personne → \eng{works}.
\item \eng{he is agree} → \textbf{Gr} : \eng{agree} est un verbe → \eng{he agrees}.
\item \eng{agree : english} → \textbf{P + Sp} : pas d'espace avant les deux-points, et majuscule → \eng{agrees: English}.
\item \eng{for find a job} → \textbf{Gr} : après la préposition \eng{for}, le verbe prend \eng{-ing} → \eng{for finding a job}.
\end{itemize}
Version corrigée : \eng{Currently, I am going to attend evening courses because I will improve my English. My brother works in Douala and he agrees: English is very important for finding a job.} — « Actuellement, je vais suivre des cours du soir parce que je veux améliorer mon anglais. Mon frère travaille à Douala et il est d'accord : l'anglais est très important pour trouver un emploi. »
\end{exoresolu}

\begin{exercice}[À vous de jouer]
Corrigez ce paragraphe en utilisant les symboles Gr, V, Sp et P, puis réécrivez-le entièrement :\\
\eng{Next year, i will to work in a bank, so actually I am going to improve me english. My teacher say that english is the language of business. I am agree with her, and I will attend courses every monday.}
\end{exercice}

\begin{corrige}[Exercice « À vous de jouer »]
Fautes : \eng{i} → \eng{I} (Sp) ; \eng{will to work} → \eng{will work} (Gr) ; \eng{so actually} → \eng{so} suffit, ou \eng{that is why} (V : \eng{actually} est un faux ami qui signifie « en fait ») ; \eng{improve me english} → \eng{improve my English} (Gr + Sp) ; \eng{My teacher say} → \eng{My teacher says} (Gr) ; \eng{english} → \eng{English} (Sp) ; \eng{I am agree} → \eng{I agree} (Gr) ; \eng{every monday} → \eng{every Monday} (Sp).\\
Version corrigée : \eng{Next year, I will work in a bank, so I am going to improve my English. My teacher says that English is the language of business. I agree with her, and I will attend courses every Monday.} — « L'année prochaine, je travaillerai dans une banque, alors je vais améliorer mon anglais. Mon professeur dit que l'anglais est la langue des affaires. Je suis d'accord avec elle, et je suivrai des cours tous les lundis. »
\end{corrige}

\begin{aretenir}[L'essentiel de la section]
\begin{itemize}
\item Structure : un plan visible, une \eng{topic sentence} par paragraphe, une conclusion qui conclut sans copier l'introduction.
\item Grammaire : \eng{-s} à la 3\textsuperscript{e} personne ; \eng{will} + base verbale (sans \eng{to}) ; \eng{be going to} + infinitif pour les projets ; \eng{I agree}, jamais \eng{I am agree} ; verbe en \eng{-ing} après une préposition (\eng{for finding}).
\item Vocabulaire : \eng{currently} = actuellement, \eng{actually} = en fait ; \eng{to attend} = assister à, \eng{to assist} = aider.
\item Forme : majuscules à \eng{I}, \eng{English}, aux pays, jours et mois ; pas d'espace avant la ponctuation.
\item Relecture en trois passes : idées, grammaire, orthographe — et lecture à voix basse pour traquer les phrases bancales.
\end{itemize}
Vous voilà prêt pour la dernière étape : la grille d'évaluation et l'entraînement en conditions d'examen.
\end{aretenir}

\coursec{Grille d'évaluation et entraînement au Bac}

Dans la section précédente, nous avons appris à relire notre composition et à corriger les erreurs les plus fréquentes des francophones : oubli du \cle{-s} de la troisième personne, confusion entre \eng{job} et \eng{work}, temps verbaux incohérents. Mais relire ne suffit pas : il faut aussi savoir \emph{comment le correcteur va nous noter}. C'est exactement l'objet de cette dernière section. Comprendre la grille d'évaluation, c'est comprendre ce que l'examinateur attend de vous — et donc écrire en conséquence. Terminons par un entraînement chronométré, dans les conditions réelles du Baccalauréat A.

\courssub{La grille d'évaluation : comment le correcteur vous note}

Au Baccalauréat, la composition n'est pas notée « au feeling ». Le correcteur applique une grille précise, qui répartit les points selon cinq critères. Voici la grille que nous utiliserons en classe, sur 10 points, fidèle à l'esprit de l'évaluation officielle.

\begin{popfigure}
\begin{tikzpicture}[
  crit/.style={draw=popblue, fill=popblueL, rounded corners=2pt, minimum height=0.9cm, align=left, text width=5.6cm, font=\small},
  pts/.style={draw=poporange, fill=poporangeL, rounded corners=2pt, minimum height=0.9cm, align=center, text width=1.6cm, font=\small\bfseries},
  desc/.style={draw=popline, fill=popcream, rounded corners=2pt, minimum height=0.9cm, align=left, text width=7.2cm, font=\small},
  head/.style={fill=popblue, text=white, rounded corners=2pt, minimum height=0.8cm, align=center, font=\small\bfseries}
]
\node[head, text width=5.6cm] (h1) at (0,0) {Critère};
\node[head, text width=1.6cm] (h2) at (3.8,0) {Barème};
\node[head, text width=7.2cm] (h3) at (8.4,0) {Ce que le correcteur vérifie};

\node[crit] (c1) at (0,-1.0) {1. Respect du sujet (\eng{task achievement})};
\node[pts] (p1) at (3.8,-1.0) {2 pts};
\node[desc] (d1) at (8.4,-1.0) {Chaque paragraphe répond au sujet posé ; aucun hors sujet.};

\node[crit] (c2) at (0,-2.1) {2. Organisation (\eng{structure})};
\node[pts] (p2) at (3.8,-2.1) {2 pts};
\node[desc] (d2) at (8.4,-2.1) {Introduction, développement en paragraphes, conclusion ; alinéas visibles.};

\node[crit] (c3) at (0,-3.2) {3. Connecteurs logiques (\eng{linking})};
\node[pts] (p3) at (3.8,-3.2) {2 pts};
\node[desc] (d3) at (8.4,-3.2) {\eng{first, moreover, however, therefore, in conclusion} utilisés à bon escient.};

\node[crit] (c4) at (0,-4.3) {4. Correction grammaticale (\eng{accuracy})};
\node[pts] (p4) at (3.8,-4.3) {2 pts};
\node[desc] (d4) at (8.4,-4.3) {Temps verbaux, accords, \cle{-s} de la 3\textsuperscript{e} personne, ordre des mots.};

\node[crit] (c5) at (0,-5.4) {5. Vocabulaire (\eng{range})};
\node[pts] (p5) at (3.8,-5.4) {2 pts};
\node[desc] (d5) at (8.4,-5.4) {Lexique du travail et des langues : \eng{to apply for, skills, fluent, training}\dots};

\node[draw=popdark, fill=popgoldL, rounded corners=2pt, minimum height=0.8cm, align=center, text width=15.4cm, font=\small\bfseries] at (4.2,-6.5) {TOTAL : 10 points};
\end{tikzpicture}
\legende{La grille d'évaluation de la composition : cinq critères, deux points chacun.}
\end{popfigure}

\begin{propriete}[Ce que le correcteur récompense vraiment]
Une composition claire, bien organisée, écrite dans un anglais \emph{simple mais correct}, obtient toujours une meilleure note qu'une composition ambitieuse pleine de fautes. Le correcteur valorise :
\begin{itemize}
\item la \cle{clarté} : des phrases courtes et maîtrisées (\eng{English is essential in my future job.}) plutôt que des phrases longues et bancales ;
\item la \cle{cohérence} : chaque idée est reliée à la précédente par un connecteur ;
\item la \cle{pertinence} : chaque paragraphe répond directement au sujet.
\end{itemize}
Règle d'or : \eng{Keep it simple, keep it correct.} (« Restez simple, restez correct. ») N'écrivez jamais en anglais une phrase que vous ne seriez pas capable de corriger vous-même.
\end{propriete}

\begin{attention}[Le hors sujet : la faute qui coûte le plus cher]
Le \cle{hors sujet} est sanctionné très sévèrement : une composition entièrement hors sujet peut perdre la quasi-totalité des points du critère 1 et pénaliser tous les autres critères. Avant de rédiger chaque paragraphe, posez-vous la question : \emph{ce paragraphe répond-il au sujet posé ?} Par exemple, si le sujet est \eng{“Discuss the English language skills you will need in your future career”}, un long paragraphe sur vos loisirs ou votre famille est un hors sujet, même s'il est écrit dans un anglais parfait. Soulignez les mots-clés du sujet (\eng{skills}, \eng{future career}) et revenez-y sans cesse.
\end{attention}

\courssub{L'évaluation par les pairs : apprendre à noter pour apprendre à écrire}

Une excellente façon de progresser est de se mettre dans la peau du correcteur. En classe, nous pratiquons l'\cle{évaluation par les pairs} (\eng{peer assessment}) : vous échangez votre brouillon avec un camarade, et vous notez sa composition avec la grille officielle.

\begin{methode}[Comment évaluer la composition d'un camarade]
\begin{enumerate}
\item \textbf{Lisez} la composition une première fois sans rien noter, pour en saisir le sens général.
\item \textbf{Vérifiez le sujet} : soulignez en vert les phrases qui répondent au sujet, en rouge celles qui s'en écartent.
\item \textbf{Comptez les connecteurs} : entourez \eng{first, moreover, however, therefore, in conclusion}. Y en a-t-il au moins quatre, bien placés ?
\item \textbf{Repérez trois fautes de grammaire} maximum (accord sujet-verbe, temps, préposition) et proposez une correction.
\item \textbf{Attribuez les points} critère par critère, puis justifiez chaque note par une phrase en anglais : \eng{“I gave 2 points for linking because the writer used five connectors correctly.”}
\item \textbf{Rendez la copie} avec un conseil constructif : \eng{“To improve, add a clearer conclusion.”}
\end{enumerate}
\end{methode}

\begin{exemplebox}[Une évaluation par les pairs réussie]
Voici le commentaire rédigé par une élève de Tle A après avoir noté la copie de son voisin :

\eng{“Your composition answers the topic well and your introduction is clear. I gave you 2 points for task achievement and 2 for vocabulary. However, you forgot the -s in “she need” and you used “actually” with the wrong meaning. I gave you 1 point for accuracy. To improve, check your verbs and use “nowadays” instead of “actually”.”}

« Ta composition répond bien au sujet et ton introduction est claire. Je t'ai donné 2 points pour le respect du sujet et 2 pour le vocabulaire. Cependant, tu as oublié le -s dans “she need” et tu as utilisé “actually” avec un mauvais sens. Je t'ai donné 1 point pour la correction grammaticale. Pour progresser, vérifie tes verbes et utilise “nowadays” à la place de “actually”. »
\end{exemplebox}

\begin{attention}[Le faux ami \eng{actually}]
L'exemple ci-dessus illustre l'un des faux amis les plus fréquents des francophones. \eng{Actually} ne signifie pas « actuellement » mais « en fait, en réalité ». Pour dire « actuellement », utilisez \eng{nowadays}, \eng{at present} ou \eng{currently}. Comparez : \eng{I actually passed the exam.} (« J'ai en fait réussi l'examen. ») et \eng{I am currently attending evening classes.} (« Je suis actuellement des cours du soir. »)
\end{attention}

\courssub{Gérer son temps à l'examen : la règle des 55 minutes}

À l'épreuve, le temps est votre adversaire principal. Beaucoup d'élèves se lancent dans la rédaction dès la lecture du sujet — et se retrouvent bloqués au milieu, sans conclusion. Voici le découpage que nous vous recommandons pour une épreuve de 55 minutes consacrée à la composition.

\begin{center}
\begin{tabularx}{\linewidth}{|l|X|X|}
\hline
\rowcolor{popblueL} \textbf{Étape} & \textbf{Durée} & \textbf{Ce que vous faites} \\
\hline
1. Analyse et plan & 10 min & Souligner les mots-clés du sujet, noter les idées au brouillon, rédiger le plan (introduction, 2 ou 3 paragraphes, conclusion). \\
\hline
2. Rédaction & 35 min & Écrire la composition en suivant le plan, sans chercher la perfection : phrases simples et correctes, connecteurs à chaque étape. \\
\hline
3. Relecture & 10 min & Vérifier les verbes (accord, temps), les connecteurs, l'orthographe ; corriger proprement. \\
\hline
\end{tabularx}
\end{center}

\begin{methode}[Gérer son temps au Bac : ne jamais sauter l'étape du plan]
\begin{enumerate}
\item \textbf{Ne sautez jamais le plan.} Dix minutes de planification font gagner vingt points de clarté : une composition planifiée ne part jamais en hors sujet.
\item \textbf{Écrivez le plan au brouillon} en anglais, en mots-clés : \eng{Intro: English = key skill. P1: needs at work (emails, meetings). P2: how to prepare (courses, practice). Concl.: my decision.}
\item \textbf{Surveillez la montre} : si après 35 minutes votre conclusion n'est pas écrite, écrivez-la immédiatement, même en deux phrases. Une composition sans conclusion perd des points d'organisation.
\item \textbf{Gardez 10 minutes de relecture} quoi qu'il arrive : c'est là que se gagnent les points de grammaire.
\end{enumerate}
\end{methode}

\courssub{La présentation matérielle : votre copie est votre première impression}

Avant même de lire la première phrase, le correcteur voit votre copie. Une présentation soignée crée une impression favorable ; une copie illisible fatigue le correcteur et peut vous coûter des points.

\begin{itemize}
\item \textbf{Écriture lisible} : écrivez sans vous presser ; un correcteur qui doit déchiffrer chaque mot est un correcteur moins indulgent.
\item \textbf{Alinéas visibles} : commencez chaque nouveau paragraphe en retrait (\eng{indent}) ou laissez une ligne entre les paragraphes. L'organisation doit se \emph{voir}.
\item \textbf{Pas de ratures excessives} : pour corriger, barrez proprement le mot d'un seul trait et écrivez la correction au-dessus ou à la suite. Ne noircissez jamais un passage entier.
\item \textbf{Pas d'abréviations de SMS} : écrivez \eng{because}, jamais \eng{coz} ; \eng{you}, jamais \eng{u}.
\end{itemize}

\begin{savaistu}[Un sujet classique du Baccalauréat A]
Au Baccalauréat série A au Cameroun, l'expression écrite porte très souvent sur des sujets de vie sociale et professionnelle : l'importance des langues, le monde du travail, l'éducation, la vie en famille ou en communauté. Le thème de cette leçon — anticiper ses besoins en anglais pour sa future carrière — est donc un véritable \emph{classique} de l'examen. Le maîtriser, c'est vous préparer à un sujet qui revient sous des formes variées : \eng{“The importance of English in your future job”}, \eng{“Why young Cameroonians should learn English”}, \eng{“Languages and employment in Cameroon”}. Les plans, les connecteurs et le vocabulaire appris ici sont directement réutilisables.
\end{savaistu}

\courssub{Entraînement chronométré : un sujet type Bac}

Il est temps de passer à la pratique, dans les conditions réelles de l'examen. Voici d'abord un exercice résolu qui montre comment un bon candidat analyse le sujet et bâtit son plan en dix minutes.

\begin{exoresolu}[Analyser un sujet et bâtir un plan en 10 minutes]
Énoncé. Sujet type Bac : \eng{“You are about to finish secondary school. Write a composition of about 200 words in which you explain what English language skills you will need in your future career and how you plan to acquire them.”} Construisez le plan de la réponse.
\tcblower
Solution détaillée.
\textbf{Étape 1 — Mots-clés du sujet} : \eng{English language skills} (quelles compétences ?), \eng{future career} (quel métier ?), \eng{how you plan to acquire them} (quelle méthode ?). Le sujet comporte donc \emph{trois questions} : la composition doit répondre aux trois.
\textbf{Étape 2 — Choix personnels} : métier choisi = \eng{nurse} (infirmière) à l'hôpital de Buea ; compétences = parler avec les patients anglophones, lire des documents médicaux en anglais ; méthode = cours du soir, écoute de la radio anglophone, pratique avec des amis.
\textbf{Étape 3 — Plan au brouillon} :
\begin{itemize}
\item \emph{Introduction} : \eng{English is essential for my future career as a nurse.}
\item \emph{Paragraph 1} : \eng{skills I will need — speaking with patients, reading medical documents.}
\item \emph{Paragraph 2} : \eng{how I will acquire them — evening classes, listening to English radio, practising with friends.}
\item \emph{Conclusion} : \eng{I am determined to start now, because English will open doors for me.}
\end{itemize}
\textbf{Étape 4 — Vérification anti-hors sujet} : chaque paragraphe répond à l'une des trois questions du sujet. Le plan est validé ; la rédaction peut commencer.
\end{exoresolu}

\begin{exercice}[Entraînement chronométré — 55 minutes]
Sujet : \eng{“Many jobs in Cameroon today require a good command of English. Write a composition of about 200 words in which you explain why English is important for your future professional life and what you are doing to improve your English.”}

Consignes :
\begin{enumerate}
\item Consacrez \textbf{10 minutes} à l'analyse du sujet et au plan au brouillon.
\item Rédigez votre composition en \textbf{35 minutes}, avec introduction, deux paragraphes de développement et conclusion, et au moins quatre connecteurs différents.
\item Réservez \textbf{10 minutes} à la relecture : verbes, accords, connecteurs, orthographe.
\item Échangez ensuite votre copie avec un camarade et notez-la avec la grille sur 10 points.
\end{enumerate}
\end{exercice}

\begin{corrige}[Entraînement chronométré]
Éléments de correction. Il n'existe pas une seule bonne réponse, mais toute bonne copie doit présenter les caractéristiques suivantes.
\textbf{Respect du sujet (2 pts)} : la composition répond aux deux questions (\emph{pourquoi} l'anglais est important et \emph{ce que} l'élève fait pour progresser) ; aucun paragraphe hors sujet.
\textbf{Organisation (2 pts)} : introduction qui annonce le thème, deux paragraphes distincts avec alinéas, conclusion personnelle.
\textbf{Connecteurs (2 pts)} : au moins quatre connecteurs variés et bien placés, par exemple \eng{first of all, moreover, however, as a result, in conclusion}.
\textbf{Correction grammaticale (2 pts)} : présent simple pour les vérités générales (\eng{English opens doors}), présent continu ou \eng{present perfect} pour les actions en cours (\eng{I am attending evening classes} / \eng{I have improved my speaking}), futur avec \eng{will} ou \eng{going to} pour les projets ; \cle{-s} de la troisième personne respecté.
\textbf{Vocabulaire (2 pts)} : lexique du travail et des langues correctement employé : \eng{to apply for a job, fluent, skills, training, career, to improve, employer}.
Exemple de phrase attendue dans le développement : \eng{“Moreover, most employers in Douala require candidates who can speak English fluently.”} (« De plus, la plupart des employeurs à Douala exigent des candidats capables de parler anglais couramment. »)
\end{corrige}

Vous disposez désormais de tous les outils : un plan solide, une structure claire, des connecteurs maîtrisés, une grille d'évaluation comprise et une gestion du temps efficace. À l'examen, rappelez-vous la règle d'or : \eng{plan first, write simply, check carefully} — planifiez d'abord, écrivez simplement, relisez soigneusement.

\newpage

\coursec{Activité d'intégration}

\coursec{Lesson 12 — Writing: Writes compositions about anticipating one's second language needs at work}

\courssub{Objectifs de la leçon}

À l'issue de cette leçon, l'élève sera capable de :
\begin{itemize}
\item analyser un sujet de composition et identifier les attentes du correcteur (contenu, structure, langue) ;
\item planifier sa production à l'aide d'une carte mentale (\eng{mind map}) et d'un plan détaillé en trois parties ;
\item rédiger une composition complète de 250 à 300 mots sur le thème \eng{anticipating one's second language needs at work}, structurée en introduction, développement (deux à trois paragraphes) et conclusion ;
\item employer correctement les connecteurs logiques (addition, opposition, cause, conséquence, exemple, conclusion), les temps verbaux adaptés (présent simple, futur \eng{will}/\eng{be going to}, conditionnel premier type) et le vocabulaire spécifique du monde professionnel ;
\item relire, corriger et évaluer sa propre production ou celle d'un camarade à l'aide d'une grille d'évaluation officielle (sur 20 points), puis formuler deux pistes d'amélioration concrètes.
\end{itemize}

\courssub{Analyser le sujet et planifier}

\begin{textebox}[A letter from the Guidance Counsellor — Government Bilingual High School, Yaoundé]
Dear Final Year Students,

As you prepare for the Baccalauréat and for your future careers, our school is organising a \emph{Career and Languages Week}. Many of you dream of becoming doctors, engineers, journalists, lawyers, farmers, business people or computer scientists. But have you ever asked yourselves one simple question: \emph{how much English will I need in my future job?}

Cameroon is a bilingual country, and English is also the main language of international business, science and technology. Employers in Douala, Buea, Bamenda and even abroad increasingly ask for candidates who can write emails, read reports and speak confidently in English.

For this reason, every student of Terminale is invited to write a composition entitled \eng{“My Career, My English: Anticipating My Second Language Needs at Work”}. The best compositions will be read aloud during the closing ceremony and displayed on the school notice board.

Your composition should be between 250 and 300 words. It should present your future career, explain the situations in which you will need English, and describe what you plan to do, starting today, to be ready.

Good luck!

The Guidance Counsellor
\end{textebox}

\textbf{Glossaire :} \emph{guidance counsellor} (n.) : conseiller d'orientation ; \emph{notice board} (n.) : panneau d'affichage ; \emph{increasingly} (adv.) : de plus en plus ; \emph{confidently} (adv.) : avec assurance ; \emph{abroad} (adv.) : à l'étranger.

\begin{methode}[Analyser le sujet — 3 étapes]
\begin{enumerate}
\item \textbf{Identifier le destinateur, le destinataire et le type de texte.} Ici : une lettre officielle (\eng{formal letter}) du conseiller d'orientation aux élèves de Terminale. La tâche demandée est une \eng{composition} (rédaction d'opinion/argumentative).
\item \textbf{Repérer les contraintes formelles.} Titre imposé : \eng{“My Career, My English: Anticipating My Second Language Needs at Work”}. Longueur : 250–300 mots. Contenu obligatoire : (1) futur métier, (2) situations nécessitant l'anglais, (3) plan d'action personnel dès maintenant.
\item \textbf{Dégager le registre et le ton.} Registre formel / semi-formel (\eng{no contractions} : \eng{I will} et non \eng{I'll} ; pas d'argot). Ton personnel et réfléchi (\eng{I think, I plan to, I am convinced that}).
\end{enumerate}
\end{methode}

\begin{experience}[Activité guidée : Compréhension et Mind Map]
\begin{enumerate}
\item \textbf{Questions de compréhension (écrit, 5 min).} Réponds en phrases complètes en anglais :
      \begin{itemize}
      \item Who is the author of the letter?
      \item What event is the school organising?
      \item What are the three mandatory content points for your composition?
      \end{itemize}
\item \textbf{Choix du métier (2 min).} Choisis ton futur métier (\eng{nurse, software developer, journalist, lawyer, cocoa exporter, teacher, pilot, engineer...}).
\item \textbf{Carte mentale — Mind Map (10 min).} Au centre : \eng{My future job: [ton métier]}. Quatre branches principales :
      \begin{itemize}
      \item \eng{Situations where I will need English} (ex. : \eng{writing reports, attending conferences, negotiating contracts, reading technical manuals, speaking with foreign clients/colleagues}).
      \item \eng{Skills to develop} (\eng{reading, writing, listening, speaking} — précise lesquelles sont prioritaires).
      \item \eng{People I will communicate with} (\eng{colleagues, clients, patients, partners, supervisors, international experts}).
      \item \eng{What I will do from now on} (\eng{read English news daily, join English club, take online courses, watch tutorials, practise with friends}).
      \end{itemize}
\end{enumerate}
\end{experience}

\courssub{La structure d'une composition}

\begin{propriete}[Structure type d'une composition argumentative (250–300 mots)]
\begin{center}
\begin{tabularx}{\linewidth}{|X|X|X|}
\hline
\rowcolor{popblueL} \textbf{Partie} & \textbf{Contenu et fonction} & \textbf{Repères linguistiques} \\
\hline
\textbf{Introduction} (≈ 40–50 mots) & \begin{itemize}\item Accroche (\eng{hook}) : fait marquant, citation, question, observation générale.\item Présentation du sujet : reformulation du titre.\item Annonce du plan (\eng{thesis statement}) : \eng{This composition will first examine..., then discuss..., and finally suggest...}\end{itemize} & Présent simple (vérités générales) ; \eng{In today's world...}, \eng{It is often said that...}, \eng{As a future [job], I...} \\
\hline
\textbf{Développement} (≈ 160–200 mots) & \begin{itemize}\item 2 ou 3 paragraphes distincts.\item 1 idée principale par paragraphe (\eng{topic sentence} en début de paragraphe).\item Exemples concrets, détails, arguments.\item Connecteurs logiques pour lier les paragraphes et les phrases.\end{itemize} & Futur \eng{will} (prédictions) / \eng{be going to} (projets décidés) ; Conditionnel 1er type (\eng{If I want to..., I will have to...}) ; Vocabulaire précis du métier. \\
\hline
\textbf{Conclusion} (≈ 40–50 mots) & \begin{itemize}\item Résumé synthétique des points forts (\eng{To conclude, In short...}).\item Bilan personnel / Ouverture : engagement, projection, pensée finale.\item Pas de nouvelle idée.\end{itemize} & Structure comparative \eng{The sooner..., the better...} ; \eng{I am determined to...}, \eng{Ultimately...} \\
\hline
\end{tabularx}
\end{center}
\end{propriete}

\courssub{Connecteurs, temps et vocabulaire utiles}

\begin{aretenir}[Boîte à outils linguistiques — À garder sous les yeux pendant la rédaction]
\textbf{A. Connecteurs logiques (Linking words) — au moins 5 différents dans la copie}
\begin{center}
\begin{tabularx}{\linewidth}{|l|X|X|}
\hline
\rowcolor{popblueL} \textbf{Fonction} & \textbf{Connecteurs (position variée)} & \textbf{Exemple court} \\
\hline
Addition & \eng{First of all, Moreover, Furthermore, In addition, Besides, Also} & \eng{Moreover, English is essential for emails.} \\
\hline
Opposition / Concession & \eng{However, Nevertheless, On the one hand... on the other hand, Whereas, Although} & \eng{However, speaking is harder than reading.} \\
\hline
Cause / Raison & \eng{Because, Since, As, Due to the fact that} & \eng{Since Cameroon is bilingual, I need English.} \\
\hline
Conséquence / Résultat & \eng{Therefore, Consequently, As a result, Thus, That is why} & \eng{Therefore, I must practise daily.} \\
\hline
Exemple / Illustration & \eng{For instance, For example, Such as, Namely} & \eng{For instance, writing medical reports.} \\
\hline
Conclusion / Résumé & \eng{To conclude, In conclusion, In short, Ultimately, All in all} & \eng{To conclude, English is a key asset.} \\
\hline
Ordre / Séquence & \eng{Firstly, Secondly, Finally, Next, Then} & \eng{Firstly, I will improve my vocabulary.} \\
\hline
\end{tabularx}
\end{center}

\textbf{B. Temps verbaux clés pour ce sujet}
\begin{center}
\begin{tabular}{|l|l|l|}
\hline
\rowcolor{popblueL} \textbf{Temps / Forme} & \textbf{Emploi ici} & \textbf{Exemples traduits} \\
\hline
Présent simple & Vérités générales, faits actuels, habitudes & \eng{English is the language of science.} (L'anglais est la langue de la science.) \\
\hline
\eng{Will} + BV & Prédictions, certitudes futures, promesses & \eng{I will need English for international conferences.} (J'aurai besoin de l'anglais...) \\
\hline
\eng{Be going to} + BV & Projets déjà décidés, intentions fermes & \eng{I am going to join the English club next week.} (Je vais rejoindre le club d'anglais...) \\
\hline
Conditionnel 1er type (\eng{If} + présent, \eng{will} + BV) & Hypothèse réaliste sur l'avenir & \eng{If I practise every day, I will become fluent.} (Si je pratique chaque jour, je deviendrai fluent.) \\
\hline
Modaux (\eng{have to, must, should, need to}) & Obligation, nécessité, conseil & \eng{I must improve my writing skills.} (Je dois améliorer...) \\
\hline
\end{tabular}
\end{center}

\textbf{C. Vocabulaire thématique — The World of Work / Le monde du travail}
\begin{center}
\begin{tabularx}{\linewidth}{|X|X|X|}
\hline
\rowcolor{popblueL} \textbf{Noms (Nouns)} & \textbf{Verbes (Verbs)} & \textbf{Adjectifs / Adverbes} \\
\hline
\eng{career, job, position, employer, employee, colleague, client, customer, interview, application, CV (résumé), cover letter, meeting, conference, report, email, memo, deadline, schedule, training, internship, work experience, qualification, skill, proficiency, fluency, promotion, salary, contract, negotiation, teamwork, leadership} & \eng{apply for, attend, conduct, negotiate, write, draft, proofread, communicate, collaborate, manage, lead, solve, improve, develop, master, acquire, require, involve, prepare, plan, anticipate} & \eng{fluent, proficient, bilingual, professional, formal, informal, technical, specific, essential, crucial, mandatory, optional, confident, confidently, regularly, effectively, efficiently} \\
\hline
\end{tabularx}
\end{center}
\end{aretenir}

\courssub{Modèle commenté}

\begin{exoresolu}[My Career, My English — Composition modèle (≈ 280 mots)]
Énoncé : Rédige une composition de 250 à 300 mots intitulée \eng{“My Career, My English: Anticipating My Second Language Needs at Work”}.
\tcblower
\textbf{Production modèle :}

\eng{In today's globalised world, mastering a second language is no longer an option but a necessity, especially in a bilingual country like Cameroon. As I am preparing for my future career as a software developer, I have started thinking seriously about the English skills I will need in the workplace.}

\eng{First of all, English will be indispensable in my daily professional activities. As a developer, I will have to read technical documentation, write code comments and communicate with international teams on platforms like GitHub. Moreover, most programming languages and frameworks are documented in English. Therefore, if I want to stay up to date with the latest technologies, I will have to read technical articles and watch tutorials in English regularly.}

\eng{However, technical reading is not enough. I will also need to speak confidently during online meetings with foreign clients and write clear emails to explain complex solutions. For instance, presenting a project demo to a client in London or New York requires fluent speaking and precise vocabulary. Consequently, I must develop all four language skills: reading, writing, listening and speaking.}

\eng{What am I going to do to be ready? At present, I read one tech article in English every evening and I listen to podcasts like “Software Engineering Daily” twice a week. In addition, I am going to join the coding club of my school where we collaborate in English. If I work consistently, I am convinced that I will be proficient before I enter the job market.}

\eng{To conclude, English is not merely a school subject for me; it is a vital professional tool. The earlier I invest in my language skills, the smoother my career path will be.}

\textbf{Commentaire détaillé des choix de langue :}
\begin{itemize}
\item \textbf{Introduction :} Accroche générale (\eng{In today's globalised world...}), ancrage contextuel (\eng{bilingual country like Cameroon}), annonce du métier (\eng{software developer}) et thèse implicite (nécessité de l'anglais).
\item \textbf{Connecteurs variés (7 utilisés) :} \eng{First of all} (ordre), \eng{Moreover} (addition), \eng{Therefore} (conséquence), \eng{However} (opposition), \eng{For instance} (exemple), \eng{Consequently} (conséquence), \eng{In addition} (addition), \eng{To conclude} (conclusion).
\item \textbf{Temps maîtrisés :} Présent simple (\eng{is, requires}), Futur \eng{will have to / will need} (besoins futurs), \eng{am going to join} (projet décidé), Conditionnel 1er type (\eng{If I want to stay..., I will have to read...} ; \eng{If I work consistently, I will be proficient...}), Modaux (\eng{must develop}).
\item \textbf{Vocabulaire précis :} \eng{technical documentation, code comments, frameworks, up to date, project demo, proficient, vital professional tool, invest in}.
\item \textbf{Conclusion :} Bilan (\eng{vital professional tool}) + structure comparative correlative \eng{The earlier..., the smoother...} (Plus tôt..., plus fluide...).
\end{itemize}
\end{exoresolu}

\courssub{Erreurs à éviter et relecture}

\begin{attention}[Pièges fréquents des francophones — Check-list de relecture]
\begin{itemize}
\item \textbf{\eng{Will} + infinitif SANS \eng{to}} : \eng{I will learn} (pas \eng{*I will to learn}). Idem pour \eng{can, must, should, might}.
\item \textbf{Faux ami \eng{actually}} : \eng{actually} = « en fait / en réalité » (pas « actuellement »). Pour « actuellement » : \eng{at present, currently, nowadays}.
\item \textbf{\eng{Information} indénombrable} : Jamais \eng{an information}. Dire \eng{a piece of information} ou \eng{some information}.
\item \textbf{Ordre des mots (SVOMPT)} : Sujet + Verbe + Objet + Manière + Lieu + Temps. \eng{I speak English fluently at work every day.} (Pas \eng{*I speak fluently English}).
\item \textbf{Accord 3e personne du singulier (présent)} : \eng{He needs, She works, It requires} (pas \eng{*he need, *she work}).
\item \textbf{Prépositions après verbes/adjectifs} : \eng{apply FOR a job}, \eng{good AT}, \eng{interested IN}, \eng{responsible FOR}, \eng{communicate WITH someone ABOUT something}.
\item \textbf{Majuscules} : \eng{I} (toujours), noms propres (\eng{Cameroon, English, Monday, January, BBC, GitHub}), début de phrase. Pas de majuscule à \eng{english} (langue) si ce n'est pas en début de phrase ? \textbf{SI :} \eng{English} prend toujours une majuscule (nationalité/langue).
\item \textbf{Orthographe \eng{–ing} / \eng{–ed}} : \eng{I am interested in} (pas \eng{*interrested}), \eng{developing} (double \eng{p}), \eng{writing} (perte du \eng{e} muet).
\item \textbf{Ponctuation} : Pas d'espace avant la virgule, le point, les deux-points, le point-virgule. Espace après. Guillemets anglais “ ” pour les citations/dialogues.
\end{itemize}
\end{attention}

\begin{methode}[Relire efficacement — La méthode COPS (5 minutes)]
\begin{enumerate}
\item \textbf{C — Capitals :} Majuscules (\eng{I}, noms propres, début de phrase, \eng{English}) ?
\item \textbf{O — Organisation :} 3 parties visibles (alinea sauté) ? \eng{Topic sentences} ? Connecteurs (soulignés) ? 250–300 mots (compter) ?
\item \textbf{P — Punctuation \& Grammar :} Accords SV ? Temps (\eng{will} sans \eng{to}) ? Prépositions ? \eng{–s} au présent ? \eng{Information} sans \eng{a} ?
\item \textbf{S — Spelling \& Vocabulary :} Mots difficiles (\eng{necessary, definitely, receive, achieve, fluent, colleague, environment}) ? Vocabulaire précis (pas \eng{very good} mais \eng{efficient, proficient, crucial}) ? Faux amis évités ?
\end{enumerate}
\end{methode}

\courssub{Grille d'évaluation et entraînement au Bac}

\begin{propriete}[Grille d'évaluation officielle — Composition (Total : 20 points)]
\begin{center}
\begin{tabularx}{\linewidth}{|X|c|X|}
\hline
\rowcolor{popblueL} \textbf{Critère} & \textbf{Points} & \textbf{Indicateurs de réussite (Niveau "Bien / Très bien")} \\
\hline
\textbf{1. Respect du sujet et de la structure} & 5 & Titre exact recopié. Introduction (accroche + sujet + plan). Développement : 2–3 paragraphes distincts, 1 idée/paragraphe. Conclusion (bilan + ouverture). Longueur 250–300 mots. \\
\hline
\textbf{2. Cohérence et connecteurs} & 5 & Idées logiquement liées. Minimum 5 connecteurs différents utilisés correctement (addition, opposition, cause, conséquence, exemple, conclusion). Transitions fluides entre paragraphes. \\
\hline
\textbf{3. Correction grammaticale} & 5 & Temps verbaux adaptés (présent, futur \eng{will}/\eng{be going to}, conditionnel 1er type). Accords sujet-verbe. Ordre des mots (SVOMPT). Prépositions. Modaux. Absence d'erreurs graves bloquant la compréhension. \\
\hline
\textbf{4. Vocabulaire et orthographe} & 5 & Lexique varié et précis du monde du travail (\eng{skills, proficiency, deadline, report, fluent...}). Pas de répétitions excessives. Orthographe correcte (mots usuels + lexique du thème). Registre formel/semi-formel tenu. \\
\hline
\end{tabularx}
\end{center}
\end{propriete}

\begin{experience}[Évaluation par les pairs — Peer Assessment (15 min)]
\begin{enumerate}
\item Échange ta copie avec ton voisin.
\item Note la composition sur 20 en utilisant la grille ci-dessus. Écris la note et un commentaire par critère en marge.
\item \textbf{Rédige deux pistes d'amélioration concrètes} (exemples) :
      \begin{itemize}
      \item \eng{Add a linking word at the beginning of paragraph 2 (e.g. "Moreover").}
      \item \eng{Replace "very important" with "crucial" or "essential" in paragraph 3.}
      \item \eng{Check the agreement in "She need" $\rightarrow$ "She needs".}
      \item \eng{Add a topic sentence to paragraph 1 to announce the main idea.}
      \end{itemize}
\item Remets la copie annotée à son auteur. Lis les commentaires reçus et note tes deux priorités de progression pour la prochaine composition.
\end{enumerate}
\end{experience}

\begin{exercice}[Entraînement au Bac — Sujets types]
Rédige une composition de 250 à 300 mots sur l'un des sujets suivants. Respecte la structure, le temps imparti (35 min) et la grille d'évaluation.
\begin{enumerate}
\item \eng{“My Future Career in Tourism: Why English is a Must.”} (Métier : guide touristique, agent de voyage, réceptionniste d'hôtel à Kribi ou Limbe.)
\item \eng{“English in the Medical Field: Preparing to Save Lives.”} (Métier : infirmier, médecin, pharmacien — contexte : hôpital de référence, ONG internationales.)
\item \eng{“From the Classroom to the Boardroom: My Business English Roadmap.”} (Métier : entrepreneur, exportateur de cacao/café, manager, comptable.)
\item \eng{“Engineering the Future: The Role of English in Technical Innovation.”} (Métier : ingénieur civil, informaticien, agronome, énergies renouvelables.)
\end{enumerate}
\end{exercice}

\begin{corrige}[Exercice 1 — Pistes de correction pour le sujet Tourisme]
\textbf{Structure attendue :}
\begin{itemize}
\item \textbf{Intro :} Accroche sur l'essor du tourisme au Cameroun (Kribi, Limbe, parcs nationaux). Sujet : l'anglais, langue commune des touristes. Plan : 1) Besoins quotidiens (accueil, visites), 2) Gestion des imprévus/réservations, 3) Plan de formation personnel.
\item \textbf{Dév. 1 (Accueil/Guiding) :} \eng{First of all, I will welcome English-speaking tourists... I will have to explain historical sites... Moreover, I will write brochures...}
\item \textbf{Dév. 2 (Gestion/Imprévus) :} \eng{However, tourism involves unexpected situations... Therefore, I must understand complaints... For instance, a lost booking...}
\item \textbf{Dév. 3 (Plan d'action) :} \eng{What am I going to do? At present, I watch travel vlogs in English... I am going to take a "Tourism English" online course... If I practise role-plays, I will gain confidence.}
\item \textbf{Conclusion :} \eng{To conclude, English is the passport to a successful career in tourism. The more I practise now, the better service I will provide.}
\end{itemize}
\textbf{Connecteurs ciblés :} \eng{First of all, Moreover, However, Therefore, For instance, In addition, To conclude}.
\textbf{Temps ciblés :} \eng{will have to, must, am going to take, If I practise... I will gain}.
\textbf{Vocabulaire ciblé :} \eng{brochure, booking, complaint, itinerary, guest, fluent, hospitality, landmark, reservation, check-in}.
\end{corrige}

\courssub{Bilan de la leçon}

Cette leçon a permis de mobiliser l'ensemble des savoir-faire de l'expression écrite dans une tâche authentique et personnelle, ancrée dans le contexte camerounais (bilinguisme officiel, marché du travail local et international). Chaque élève a d'abord \cle{analysé} un sujet réel, puis \cle{planifié} ses idées grâce à une carte mentale avant de \cle{rédiger} un texte structuré en trois parties. La maîtrise des \cle{connecteurs logiques}, des \cle{temps du futur et de l'hypothèse} et du \cle{vocabulaire professionnel} a été le cœur de l'apprentissage. La phase d'\cle{évaluation croisée} par grille critériée a développé l'esprit critique et l'autonomie : noter la copie d'un camarade oblige à identifier précisément les marqueurs de la qualité (structure, cohérence, correction, lexique). Enfin, la relecture systématique (méthode COPS) ancre les habitudes de correction indispensables pour le Baccalauréat.

\begin{savaistu}[Le bilinguisme au Cameroun : atout professionnel]
Le Cameroun est l'un des rares pays d'Afrique à avoir deux langues officielles : le français et l'anglais. Cette particularité historique (héritage des mandats français et britannique après la Première Guerre mondiale) est un atout majeur sur le marché du travail. Les entreprises basées à Douala (poumon économique), Yaoundé (capitale administrative), Buea et Bamenda (capitales régionales anglophones) recherchent activement des cadres \eng{bilingual} capables de basculer d'une langue à l'autre (\eng{code-switching}) selon l'interlocuteur. Maîtriser l'anglais "professionnel" (rédaction de mails, lecture de contrats, participation à des visioconférences) ouvre l'accès aux marchés du Nigeria, d'Afrique du Sud, du Kenya, du Ghana et aux organisations internationales (ONU, UA, CEMAC). Anticiper ses besoins linguistiques dès la Terminale, c'est investir dans son \cle{employability} (employabilité).
\end{savaistu}

\newpage

\coursec{Méthodes et exercices résolus}

\coursec{Lesson 12 — Writing: Writes compositions about anticipating one's second language needs at work}

\courssub{Méthode 1 — Analyser le sujet et planifier sa composition}

\begin{methode}[Analyser le sujet et construire un plan en trois parties]
\begin{enumerate}
\item \textbf{Read the topic twice.} Souligne les mots-clés du sujet (ex. \eng{anticipating}, \eng{second language needs}, \eng{at work}) et identifie le type de texte demandé : composition argumentative, descriptive ou narrative.
\item \textbf{Brainstorm.} Note en vrac, en anglais, 6 à 8 idées : pourquoi l'anglais est utile au travail, quels besoins prévoir (réunions, e-mails, présentations, clients étrangers), comment s'y préparer (cours du soir, applications, pratique quotidienne).
\item \textbf{Select and group.} Garde les 3 meilleures idées : ce seront tes 3 paragraphes de développement. Une idée = un paragraphe.
\item \textbf{Build the plan.} Introduction (accroche + annonce du plan) ; développement (3 paragraphes avec connecteurs) ; conclusion (bilan + ouverture).
\item \textbf{Check the time.} Au Bac : 5 min d'analyse, 10 min de plan, 25 min de rédaction, 5 min de relecture.
\end{enumerate}
\textbf{Structure à retenir :} \eng{Introduction — Body paragraph 1 — Body paragraph 2 — Body paragraph 3 — Conclusion.}
\end{methode}

\begin{exoresolu}[Analyse d'un sujet type Bac]
\textbf{Énoncé.} \eng{Topic: “English is becoming essential in the Cameroonian workplace. Write a composition of about 200 words explaining how a young school-leaver can anticipate his or her second language needs at work.”} — 1) De quel type de composition s'agit-il ? 2) Quels sont les mots-clés ? 3) Propose un plan en trois parties.
\tcblower
\textbf{Solution détaillée.}
\begin{enumerate}
\item \textbf{Type de composition :} il s'agit d'une composition \textbf{argumentative et explicative} : le verbe \eng{explaining} demande de donner des moyens concrets, pas de raconter une histoire.
\item \textbf{Mots-clés :} \eng{essential in the Cameroonian workplace} (contexte : le travail au Cameroun), \eng{young school-leaver} (le personnage : un jeune qui sort de l'école), \eng{anticipate his or her second language needs} (le thème : prévoir ses besoins en anglais), \eng{about 200 words} (la longueur).
\item \textbf{Plan proposé :}
\begin{itemize}
\item \emph{Introduction :} accroche — \eng{In Cameroon today, English is no longer a luxury but a professional tool.} Annonce : un jeune diplômé peut anticiper ses besoins de trois façons.
\item \emph{Body 1 :} identifier ses besoins futurs (e-mails, réunions bilingues, entretiens d'embauche).
\item \emph{Body 2 :} se former (cours du soir à Yaoundé ou Douala, applications gratuites, radio anglophone).
\item \emph{Body 3 :} pratiquer au quotidien (parler anglais avec des collègues, lire la presse anglophone).
\item \emph{Conclusion :} celui qui anticipe ses besoins linguistiques décroche les meilleurs emplois.
\end{itemize}
\end{enumerate}
\textbf{Conclusion :} un bon plan est déjà la moitié de la note : le correcteur voit immédiatement si la composition est organisée.
\end{exoresolu}

\courssub{Méthode 2 — Rédiger introduction et conclusion efficaces}

\begin{methode}[Construire l'introduction et la conclusion]
\begin{enumerate}
\item \textbf{Introduction = hook + topic + plan.} Commence par une accroche (fait général, question) : \eng{Have you ever thought of the language you will need in your future job?} Puis présente le sujet et annonce le plan : \eng{This composition will show how...}
\item \textbf{Évite de recopier le sujet mot pour mot} : reformule-le avec tes propres mots.
\item \textbf{Conclusion = summary + personal opinion + opening.} Résume l'essentiel (\eng{To sum up, ...}), donne ton avis (\eng{In my opinion, ...}) et ouvre sur une perspective (\eng{The future belongs to those who prepare for it today.}).
\item \textbf{Cohérence :} la conclusion ne doit introduire aucune idée nouvelle.
\end{enumerate}
\textbf{Expressions à retenir :} \eng{It is widely believed that...} / \eng{This essay aims at showing that...} / \eng{To sum up,...} / \eng{All in all,...}
\end{methode}

\begin{exoresolu}[Rédiger une introduction et une conclusion]
\textbf{Énoncé.} Sujet : \eng{“How can workers in Douala anticipate their English language needs?”} Rédige une introduction de 3 phrases et une conclusion de 3 phrases.
\tcblower
\textbf{Solution détaillée.}

\emph{Introduction (modèle) :} \eng{In Douala, the economic capital of Cameroon, more and more companies do business with English-speaking partners. Many workers suddenly discover that their English is too weak for meetings, e-mails or negotiations. This composition will explain how they can anticipate their second language needs before it is too late.}
« À Douala, capitale économique du Cameroun, de plus en plus d'entreprises font affaire avec des partenaires anglophones. Beaucoup de travailleurs découvrent soudain que leur anglais est trop faible pour les réunions, les e-mails ou les négociations. Cette composition expliquera comment ils peuvent anticiper leurs besoins en langue seconde avant qu'il ne soit trop tard. »

Justification : phrase 1 = accroche factuelle ; phrase 2 = problème ; phrase 3 = annonce du plan avec \eng{will explain} (futur pour annoncer ce que le texte va faire).

\emph{Conclusion (modèle) :} \eng{To sum up, anticipating one's English needs means identifying future tasks, getting trained and practising every day. In my opinion, no worker in Douala should wait for a promotion interview to start learning. Those who prepare today will be the leaders of tomorrow.}

Justification : \eng{To sum up} introduit le bilan ; \eng{In my opinion} donne l'avis personnel ; la dernière phrase ouvre sans idée nouvelle. Noter le futur \eng{will be} pour une prédiction.
\end{exoresolu}

\courssub{Méthode 3 — Organiser le développement avec des connecteurs}

\begin{methode}[Enchaîner les idées avec les bons connecteurs]
\begin{enumerate}
\item \textbf{Un paragraphe = une idée + un exemple.} Chaque paragraphe commence par une phrase d'idée (\emph{topic sentence}) : \eng{First of all, a worker must identify his future language needs.}
\item \textbf{Addition :} \eng{first of all, moreover, furthermore, in addition, besides.}
\item \textbf{Cause et conséquence :} \eng{because, since, therefore, as a result, consequently, that is why.}
\item \textbf{Opposition et concession :} \eng{however, whereas, although, even though, on the contrary.}
\item \textbf{Illustration :} \eng{for example, for instance, such as.}
\item \textbf{Conclusion :} \eng{to sum up, in conclusion, all in all.}
\item \textbf{Ponctuation :} un connecteur en début de phrase est suivi d'une virgule : \eng{Moreover, English opens doors to international companies.}
\end{enumerate}
\end{methode}

\begin{exoresolu}[Compléter un paragraphe avec des connecteurs]
\textbf{Énoncé.} Complète avec : \eng{first of all — for example — however — therefore — moreover.}

\eng{........, a young graduate should list the situations where English will be needed, ........ job interviews and video conferences. ........, he can improve his listening skills by watching English-language news every evening. ........, many workers think they have no time to study. ........, they should use short daily sessions of fifteen minutes, which are very effective.}
\tcblower
\textbf{Solution détaillée.}
\begin{enumerate}
\item \eng{First of all} : ouvre l'énumération des conseils (premier point du plan).
\item \eng{for example} : \eng{job interviews and video conferences} sont des exemples de situations ; après \eng{for example}, la virgule est correcte.
\item \eng{Moreover} : ajoute un deuxième conseil (addition d'idées).
\item \eng{However} : introduit une objection (\eng{many workers think they have no time}) qui s'oppose aux conseils précédents.
\item \eng{Therefore} : exprime la conséquence logique — puisque le temps manque, il faut des séances courtes.
\end{enumerate}
\emph{Paragraphe corrigé :} \eng{First of all, a young graduate should list the situations where English will be needed, for example, job interviews and video conferences. Moreover, he can improve his listening skills by watching English-language news every evening. However, many workers think they have no time to study. Therefore, they should use short daily sessions of fifteen minutes, which are very effective.}

\textbf{Conclusion :} le choix du connecteur dépend de la relation logique entre les phrases, jamais du hasard.
\end{exoresolu}

\courssub{Méthode 4 — Employer les bons temps verbaux et le bon vocabulaire}

\begin{methode}[Temps verbaux et lexique du thème]
\begin{enumerate}
\item \textbf{Present simple} pour les vérités générales et les habitudes : \eng{English is the language of international trade.}
\item \textbf{Futur (will / be going to)} pour les besoins et projets professionnels : \eng{I will need English for my future career.} / \eng{She is going to attend evening classes.}
\item \textbf{Present perfect} pour l'expérience acquise : \eng{He has already learnt computer vocabulary.}
\item \textbf{Modaux} pour les conseils : \eng{should, must, can, may} : \eng{Workers should practise English every day.}
\item \textbf{Vocabulaire utile :} \eng{a school-leaver} (un jeune sorti de l'école), \eng{a job interview} (un entretien d'embauche), \eng{to apply for a job} (postuler), \eng{evening classes} (cours du soir), \eng{language skills} (compétences linguistiques), \eng{a promotion} (une promotion), \eng{bilingual} (bilingue), \eng{to attend a training course} (suivre une formation).
\end{enumerate}
\end{methode}

\begin{exoresolu}[Grammaire et vocabulaire en contexte]
\textbf{Énoncé.} 1) Mets le verbe à la forme correcte : \eng{a) Every day, our manager (write) e-mails in English. b) Next year, I (attend) a training course in Buea. c) She (already / learn) the vocabulary of marketing. d) Workers (should) not wait until the last minute.} 2) Traduis : « Il postule pour un emploi dans une entreprise bilingue. »
\tcblower
\textbf{Solution détaillée.}
\begin{enumerate}
\item a) \eng{writes} : present simple (habitude, \eng{every day}), 3e personne du singulier, donc \textbf{-s}. b) \eng{will attend} ou \eng{is going to attend} : futur pour un projet (\eng{next year}). c) \eng{has already learnt} : present perfect avec \eng{already} (expérience acquise, sans date précise) ; auxiliaire \eng{has} car 3e personne. d) \eng{should not wait} : après un modal, la base verbale sans \eng{to}.
\item \eng{He is applying for a job in a bilingual company.} Attention : \eng{to apply \textbf{for} a job} (préposition \eng{for} obligatoire) ; \eng{bilingual} se place devant le nom, comme tout adjectif anglais.
\end{enumerate}
\textbf{Conclusion :} dans une composition sur les besoins futurs, alterne present simple (constats), futur (projets) et modaux (conseils) : c'est cette variété qui impressionne le correcteur.
\end{exoresolu}

\courssub{Méthode 5 — Rédiger la composition complète et la relire}

\begin{methode}[Rédiger puis relire avec une grille]
\begin{enumerate}
\item \textbf{Rédige au brouillon} en suivant ton plan, sans t'arrêter sur chaque mot.
\item \textbf{Compte les mots} : environ 200 mots au Bac ; une composition de 80 mots est hors sujet de longueur.
\item \textbf{Relis avec la grille C-O-P-S :} \textbf{C}apital letters (majuscules : \eng{I}, débuts de phrase, nationalités) ; \textbf{O}rder of words (sujet + verbe + complément) ; \textbf{P}unctuation (point, virgule après connecteur) ; \textbf{S}pelling (orthographe : \eng{necessary}, \eng{business}, \eng{which}).
\item \textbf{Vérifie les accords :} 3e personne du singulier (\eng{he works}), pluriels (\eng{skills}), temps homogènes.
\item \textbf{Recopie proprement} : écriture lisible, paragraphes visibles (alinéa ou ligne sautée).
\end{enumerate}
\textbf{Grille d'évaluation type :} respect du sujet et de la longueur (4 \%) ; organisation et connecteurs (4 \%) ; correction grammaticale (4 \%) ; vocabulaire (4 \%) ; orthographe et présentation (4 \%).
\end{methode}

\begin{exoresolu}[Composition complète corrigée]
\textbf{Énoncé.} \eng{Write a composition of about 180–200 words on: “Anticipating your English language needs at work.”}
\tcblower
\textbf{Solution détaillée (modèle commenté).}

\eng{Nowadays, English has become a key to success in the Cameroonian job market. A young worker who ignores this reality may miss many opportunities. Fortunately, it is possible to anticipate one's second language needs in three simple ways.}

\eng{First of all, everyone should identify the situations where English will be necessary, for example job interviews, e-mails and meetings with foreign partners. Moreover, it is wise to get trained before starting work: evening classes in Yaoundé or Douala and free mobile applications can help a lot. Finally, daily practice is essential because a language is quickly forgotten when it is not used.}

\eng{However, many people say they have no time or no money to learn. This is not a real excuse: fifteen minutes of practice every day costs nothing and brings great results.}

\eng{To sum up, anticipating one's English needs means planning, training and practising. In my opinion, the workers who prepare today will be the managers of tomorrow.}

\emph{Commentaires :} introduction avec accroche + annonce du plan ; trois paragraphes de développement reliés par \eng{first of all, moreover, finally, however} ; temps variés (present simple, futur \eng{will be}, modal \eng{should}) ; conclusion avec bilan et ouverture ; environ 170 mots, conforme à la consigne.
\end{exoresolu}

\begin{attention}[Les 5 erreurs qui coûtent des points au Bac]
\begin{enumerate}
\item \textbf{Faux ami :} \eng{to assist} ne veut pas dire « assister à ». Pour « assister à un cours », écris \eng{to attend a course}. \eng{I assisted to a meeting} est faux : dis \eng{I attended a meeting}.
\item \textbf{Calque du français :} « apprendre l'anglais » = \eng{to learn English}, jamais \eng{to learn myself English}. Et « j'ai besoin de » = \eng{I need}, pas \eng{I have need of}.
\item \textbf{Temps verbaux :} ne mélange pas passé et présent sans raison. Pour les projets, utilise \eng{will} ou \eng{be going to}, pas le present simple : \eng{I will need English}, pas \eng{I need English tomorrow}.
\item \textbf{Ordre des mots :} l'adjectif précède le nom : \eng{a bilingual company}, pas \eng{a company bilingual}. Pas d'inversion après \eng{I think that...} : \eng{I think that English is important}.
\item \textbf{Orthographe et majuscules :} \eng{english} avec minuscule est une faute : les langues et nationalités prennent la majuscule (\eng{English, Cameroonian}). Attention aussi à \eng{necessary} (un c, deux s), \eng{business}, \eng{which} (pas \eng{wich}).
\end{enumerate}
\end{attention}

\newpage

\coursec{Exercices}

\coursec{Lesson 12 — Writing: Anticipating One's Second Language Needs at Work}

\courssub{Je vérifie mes connaissances}

\begin{exercice}[QCM : la structure d'une composition]
Entoure la bonne réponse.
\begin{enumerate}
\item The first paragraph of a composition is called:
\begin{enumerate}
\item the body \item the introduction \item the conclusion
\end{enumerate}
\item A \cle{thesis statement} is:
\begin{enumerate}
\item a question for the reader \item the main idea of the composition \item a list of vocabulary
\end{enumerate}
\item Which connector expresses a contrast?
\begin{enumerate}
\item moreover \item however \item therefore
\end{enumerate}
\item To talk about your plans for your future career, you should mainly use:
\begin{enumerate}
\item the past simple \item the future with will or be going to \item the present perfect
\end{enumerate}
\item The conclusion of a composition should:
\begin{enumerate}
\item introduce new arguments \item sum up the main ideas \item copy the introduction word for word
\end{enumerate}
\end{enumerate}
\end{exercice}

\begin{exercice}[Vrai ou faux, justifié]
Indique si chaque affirmation est vraie ou fausse, puis justifie ta réponse en une phrase en français.
\begin{enumerate}
\item \eng{A good composition can be written without a plan.}
\item \eng{Connectors help the reader follow the writer's ideas.}
\item \eng{In English, the word “English” does not need a capital letter.}
\item \eng{Each paragraph of the body should develop one main idea.}
\item \eng{Proofreading means checking your work for mistakes before submitting it.}
\end{enumerate}
\end{exercice}

\begin{exercice}[Vocabulaire du monde du travail]
Complète chaque phrase avec le mot anglais qui convient, choisi dans la liste : \eng{career — employer — skill — fluent — interview — training}.
\begin{enumerate}
\item After my \_\_\_\_\_\_\_\_\_\_ at the university, I would like to work in a bank in Douala.
\item Speaking English is an important \_\_\_\_\_\_\_\_\_\_ for this job.
\item The \_\_\_\_\_\_\_\_\_\_ offered me a contract after the \_\_\_\_\_\_\_\_\_\_.
\item My sister is \_\_\_\_\_\_\_\_\_\_ in English because she practises every day.
\item To improve my English, I will follow a language \_\_\_\_\_\_\_\_\_\_ course in Buea.
\end{enumerate}
\end{exercice}

\begin{exercice}[Conjugaison : le futur au travail]
Mets le verbe entre parenthèses à la forme qui convient (futur avec \eng{will} ou \eng{be going to}, ou conditionnel avec \eng{would}).
\begin{enumerate}
\item I \_\_\_\_\_\_\_\_\_\_ (take) an English exam next year; I have already registered.
\item If I work for an international company, I \_\_\_\_\_\_\_\_\_\_ (need) good English.
\item We think English \_\_\_\_\_\_\_\_\_\_ (open) many doors for us in the future.
\item She \_\_\_\_\_\_\_\_\_\_ (attend) evening classes next term; it is her plan.
\item If I had more time, I \_\_\_\_\_\_\_\_\_\_ (practise) English every morning.
\end{enumerate}
\end{exercice}

\courssub{Je m'entraîne}

\begin{exercice}[Traduction]
Traduis les phrases suivantes en anglais, en utilisant le vocabulaire et les temps étudiés dans la leçon.
\begin{enumerate}
\item Je vais améliorer mon anglais avant de chercher un emploi.
\item De plus, parler anglais me permettra de travailler dans une entreprise internationale.
\item Cependant, je dois d'abord passer un examen de langue.
\item Pour conclure, l'anglais sera indispensable pour ma carrière.
\item Si je pratiquais chaque jour, je deviendrais bientôt courant en anglais.
\end{enumerate}
\end{exercice}

\begin{exercice}[Texte à trous : les connecteurs]
Complète le texte suivant avec les connecteurs appropriés, choisis dans la liste : \eng{first of all — moreover — however — therefore — to sum up}.

\medskip
\noindent
\textbf{My plan for English}\par
English is essential for my future career as a journalist. \_\_\_\_\_\_\_\_\_\_, I intend to read English newspapers every day in order to enrich my vocabulary. \_\_\_\_\_\_\_\_\_\_, I plan to listen to English radio programmes, which will improve my listening skills. \_\_\_\_\_\_\_\_\_\_, reading and listening are not enough: I must also practise speaking with my classmates. \_\_\_\_\_\_\_\_\_\_, I have decided to join the English club of my school. \_\_\_\_\_\_\_\_\_\_, with regular practice and clear goals, I am confident that I will master English before I start working.
\medskip
\end{exercice}

\begin{exercice}[Correction de phrases]
Chaque phrase contient une erreur typique des francophones. Réécris chaque phrase correctement.
\begin{enumerate}
\item \eng{I am agree with your opinion about learning english.}
\item \eng{I will to study English at the university next year.}
\item \eng{My future job will permit me to travel a lot.}
\item \eng{She has twenty years and she wants to become a nurse.}
\item \eng{I am learning English since five years.}
\item \eng{The most of employers require English today.}
\item \eng{He speaks very well English.}
\item \eng{If I will work hard, I will succeed.}
\item \eng{I must to improve my pronunciation.}
\item \eng{Actually, I am preparing my examinations.}
\end{enumerate}
\end{exercice}

\begin{exercice}[Transformation : présent, futur, conditionnel]
Réécris chaque phrase selon la consigne entre parenthèses, en adaptant le temps verbal au contexte professionnel.
\begin{enumerate}
\item \eng{I work in a hospital in Yaoundé.} (imagine ton avenir : futur avec \eng{will})
\item \eng{She practises English every evening.} (exprime un projet déjà décidé : \eng{be going to})
\item \eng{He learns technical vocabulary for his job.} (exprime une hypothèse : conditionnel avec \eng{if he had time})
\item \eng{We attend English classes after school.} (annonce pour l'année prochaine : futur)
\item \eng{They use English at work.} (exprime une conséquence probable : \eng{if they find jobs abroad})
\end{enumerate}
\end{exercice}

\courssub{Je me prépare au Bac}

\begin{exercice}[Compréhension et langue — texte type Bac]
Lis le texte suivant, puis réponds aux questions en anglais, avec tes propres mots.

\medskip
\noindent
\textbf{English, a passport to employment}\par
In Cameroon today, more and more employers ask for a good command of English, even in regions where French is the main language of administration. Job advertisements in newspapers often mention “bilingual candidates preferred”, and young graduates quickly understand that English can make the difference between two applicants with the same diploma.

Take the example of Clarisse, a final-year student in Yaoundé. She dreams of becoming an air hostess. She knows that airlines require fluent English, so she has built a personal plan: she reads one short article in English every morning, listens to podcasts during her journey to school, and speaks English with her friends twice a week. “At first it was difficult,” she says, “but now I think in English without translating.”

Teachers advise students to anticipate their language needs before they leave school. A secretary needs writing skills for emails and reports; an engineer needs technical vocabulary; a shopkeeper in Douala needs spoken English to welcome foreign customers. Each profession has its own language needs, and the earlier students identify them, the better prepared they will be.

Experts agree on one point: learning English is not a one-week effort but a long-term project. Regular practice, clear goals and self-confidence are the keys to success.
\medskip

\textbf{A. Comprehension questions}
\begin{enumerate}
\item According to the text, what do many job advertisements in Cameroon mention?
\item What is Clarisse's dream job?
\item Name two activities Clarisse does to improve her English.
\item What language skills does a secretary need, according to the text?
\item What are the three keys to success mentioned in the last paragraph?
\end{enumerate}

\textbf{B. Language work}
\begin{enumerate}
\item Find in the text: (a) a connector expressing a consequence; (b) a verb in the future; (c) an adjective meaning “able to speak easily”.
\item Put into the future: \eng{She reads one short article every morning.}
\item Put into the conditional: \eng{Regular practice is the key to success.} (begin with \eng{If students practised regularly, ...})
\item Translate into French: \eng{The earlier students identify their language needs, the better prepared they will be.}
\end{enumerate}
\end{exercice}

\begin{exercice}[Traduction]
Traduis le passage suivant en anglais. Soigne les temps verbaux, les connecteurs et le vocabulaire du monde du travail.

« Pour ma carrière future, l'anglais sera très important. D'abord, je vais passer un examen d'anglais à la fin de l'année. Ensuite, je pratiquerai la langue chaque jour en lisant des journaux et en écoutant la radio. Cependant, je sais que cela demandera beaucoup d'efforts. Si je travaille régulièrement, je serai prêt pour le monde du travail. »
\end{exercice}

\begin{exercice}[Sujet de composition type Bac]
\textbf{Topic:} \eng{You are a final-year student. Write a composition of about 250 words explaining how you plan to prepare your English for your future career.}

Consignes :
\begin{enumerate}
\item Rédige une composition complète en trois parties : introduction (avec une accroche et une phrase de thèse), développement (deux ou trois paragraphes, une idée principale par paragraphe), conclusion (bilan et ouverture).
\item Utilise au moins cinq connecteurs logiques différents (par exemple \eng{first of all, moreover, however, therefore, to sum up}).
\item Emploie le futur (\eng{will}, \eng{be going to}) et au moins une phrase au conditionnel (\eng{if... would}).
\item Compte environ 250 mots ; indique le nombre de mots à la fin.
\item Temps conseillé : 45 minutes. Relis-toi et corrige les erreurs avant de rendre ta copie.
\end{enumerate}
\end{exercice}

\newpage

\coursec{Corrigés des exercices}

\begin{corrige}[Exercice 1 — QCM : la structure d'une composition]
\begin{enumerate}
\item \textbf{b. the introduction} — Le premier paragraphe présente le sujet et annonce la thèse ; le \eng{body} (développement) vient ensuite et la \eng{conclusion} termine la composition.
\item \textbf{b. the main idea of the composition} — La \cle{thesis statement} (énoncé de thèse) est la phrase qui énonce l'idée directrice de la composition, généralement à la fin de l'introduction.
\item \textbf{b. however} — \eng{However} (« cependant ») exprime le contraste. \eng{Moreover} (« de plus ») exprime l'addition et \eng{therefore} (« par conséquent ») la conséquence.
\item \textbf{b. the future with will or be going to} — Pour parler de projets professionnels futurs, on utilise \eng{will} (prévisions, promesses) ou \eng{be going to} (projets déjà décidés).
\item \textbf{b. sum up the main ideas} — La conclusion résume les idées principales et peut s'ouvrir sur une perspective ; elle n'introduit jamais de nouveaux arguments et ne recopie pas l'introduction.
\end{enumerate}
\end{corrige}

\begin{corrige}[Exercice 2 — Vrai ou faux, justifié]
\begin{enumerate}
\item \textbf{Faux.} Un plan est indispensable : il permet d'organiser les idées en introduction, développement et conclusion avant de rédiger, ce qui garantit la cohérence de la composition.
\item \textbf{Vrai.} Les connecteurs (\eng{first of all, moreover, however, therefore}) relient les phrases et les paragraphes ; ils guident le lecteur dans l'enchaînement des idées.
\item \textbf{Faux.} En anglais, les noms de langues et de nationalités prennent toujours une majuscule : \eng{English}, \eng{French}, \eng{Cameroonian}.
\item \textbf{Vrai.} La règle est « une idée principale par paragraphe » : chaque paragraphe du développement s'ouvre sur une phrase d'idée (\eng{topic sentence}) puis l'illustre.
\item \textbf{Vrai.} \eng{To proofread} signifie relire sa production pour repérer et corriger les erreurs de grammaire, d'orthographe et de ponctuation avant de la rendre.
\end{enumerate}
\end{corrige}

\begin{corrige}[Exercice 3 — Vocabulaire du monde du travail]
\begin{enumerate}
\item \eng{After my \textbf{training} at university, I would like to work in a bank in Douala.} — \eng{training} = « formation ». (On accepte aussi \eng{studies}, mais ici le mot de la liste est \eng{training}.)
\item \eng{Speaking English is an important \textbf{skill} for this job.} — \eng{skill} = « compétence ».
\item \eng{The \textbf{employer} offered me a contract after the \textbf{interview}.} — \eng{employer} = « employeur » ; \eng{interview} = « entretien (d'embauche) ».
\item \eng{My sister is \textbf{fluent} in English because she practises every day.} — \eng{fluent in} = « courant en » ; attention à la préposition \eng{in}.
\item \eng{To improve my English, I will follow a language \textbf{training} course in Buea.} — \eng{a training course} = « un stage / une formation ».
\end{enumerate}
\textbf{Remarque :} la phrase 1 admet aussi \eng{career} dans d'autres contextes, mais ici seul \eng{training} convient grammaticalement et logiquement ; \eng{career} (« carrière professionnelle ») est un faux ami de « carrière » au sens de « parcours professionnel » seulement.
\end{corrige}

\begin{corrige}[Exercice 4 — Conjugaison : le futur au travail]
\begin{enumerate}
\item \eng{I \textbf{am going to take} an English exam next year; I have already registered.} — Projet déjà décidé et préparé (\eng{I have already registered}) : \eng{be going to} + base verbale.
\item \eng{If I work for an international company, I \textbf{will need} good English.} — Phrase hypothétique du premier type : \eng{if} + présent simple, \eng{will} + base verbale dans la principale. Jamais \eng{will} après \eng{if}.
\item \eng{We think English \textbf{will open} many doors for us in the future.} — Prévision, opinion sur l'avenir (\eng{we think...}) : \eng{will} + base verbale.
\item \eng{She \textbf{is going to attend} evening classes next term; it is her plan.} — Intention, plan arrêté : \eng{be going to}.
\item \eng{If I had more time, I \textbf{would practise} English every morning.} — Hypothèse irréelle au présent : \eng{if} + past simple, \eng{would} + base verbale.
\end{enumerate}
\textbf{Point de grammaire :} après \eng{will}, \eng{would} et \eng{going to}, le verbe est toujours à la base verbale, sans \eng{to} et sans \eng{-s}. Rappel : \eng{practise} (verbe, UK) s'écrit avec un \eng{s}, \eng{practice} (nom) avec un \eng{c}.
\end{corrige}

\begin{corrige}[Exercice 5 — Traduction]
\begin{enumerate}
\item \eng{I am going to improve my English before looking for a job.} — Projet décidé : \eng{be going to} ; après la préposition \eng{before}, le verbe prend la forme en \eng{-ing}.
\item \eng{Moreover, speaking English will allow me to work for an international company.} — Sujet nominalisé en \eng{-ing} ; \eng{to allow somebody to do something} = « permettre à quelqu'un de faire quelque chose ».
\item \eng{However, I must first take a language exam.} — \eng{must} + base verbale ; \eng{to take an exam} = « passer un examen » (et non \eng{pass an exam} dans le sens « se présenter à »).
\item \eng{To sum up, English will be essential for my career.} — \eng{to sum up} = « pour conclure » ; \eng{essential} = « indispensable ».
\item \eng{If I practised every day, I would soon become fluent in English.} — Conditionnel irréel : \eng{if} + past simple, \eng{would} + base verbale ; \eng{fluent in English} = « courant en anglais ».
\end{enumerate}
\end{corrige}

\begin{corrige}[Exercice 6 — Texte à trous : les connecteurs]
\textbf{Texte complété :}

\eng{English is essential for my future career as a journalist. First of all, I intend to read English newspapers every day in order to enrich my vocabulary. Moreover, I plan to listen to English radio programmes, which will improve my listening skills. However, reading and listening are not enough: I must also practise speaking with my classmates. Therefore, I have decided to join the English club of my school. To sum up, with regular practice and clear goals, I am confident that I will master English before I start working.}

\textbf{Justifications :}
\begin{itemize}
\item \eng{First of all} ouvre l'énumération des actions prévues (premier point du plan).
\item \eng{Moreover} ajoute une deuxième action (addition d'idées).
\item \eng{However} introduit une nuance : lire et écouter ne suffisent pas (contraste).
\item \eng{Therefore} exprime la conséquence logique : puisqu'il faut parler, l'élève rejoint le club d'anglais.
\item \eng{To sum up} annonce le bilan final (conclusion).
\end{itemize}
\textbf{Remarque :} après un connecteur placé en tête de phrase, on met une virgule en anglais. Orthographe UK : \eng{programmes}, \eng{practise} (verbe).
\end{corrige}

\begin{corrige}[Exercice 7 — Correction de phrases]
\begin{enumerate}
\item \eng{I \textbf{agree} with your opinion about learning \textbf{English}.} — \eng{agree} est un verbe (pas d'auxiliaire \eng{be}) ; majuscule obligatoire à \eng{English}.
\item \eng{I will \textbf{study} English at university next year.} — Après \eng{will}, base verbale sans \eng{to} ; \eng{at university} (UK) sans article.
\item \eng{My future job will \textbf{allow} me to travel a lot.} — Faux ami : \eng{permit} existe mais est soutenu ; en anglais courant on dit \eng{allow me to travel}. (\eng{permit me to travel} est toléré mais lourd.)
\item \eng{She \textbf{is twenty} and she wants to become a nurse.} — L'âge se dit avec \eng{be}, pas \eng{have} ; \eng{twenty years old} est possible, jamais \eng{has twenty years}.
\item \eng{I \textbf{have been learning} English \textbf{for} five years.} — Action commencée dans le passé qui continue : present perfect continuous ; \eng{for} + durée, \eng{since} + point de départ.
\item \eng{\textbf{Most} employers require English today.} — « La plupart de » se dit \eng{most} + nom, sans \eng{the} ni \eng{of} (ou \eng{most of the employers} avec article).
\item \eng{He speaks English very \textbf{well}.} — Ordre des mots : verbe + complément + adverbe ; on ne sépare pas le verbe de son complément.
\item \eng{If I \textbf{work} hard, I will succeed.} — Après \eng{if}, présent simple ; jamais \eng{will} dans la subordonnée.
\item \eng{I must \textbf{improve} my pronunciation.} — \eng{must} est un modal : base verbale sans \eng{to}.
\item \eng{\textbf{At present}, I am preparing \textbf{for} my examinations.} — Faux ami : \eng{actually} = « en fait » ; « actuellement » = \eng{at present / currently} ; \eng{to prepare for an exam}.
\end{enumerate}
\end{corrige}

\begin{corrige}[Exercice 8 — Transformation : présent, futur, conditionnel]
\begin{enumerate}
\item \eng{I \textbf{will work} in a hospital in Yaoundé.} — Futur simple : \eng{will} + base verbale.
\item \eng{She \textbf{is going to practise} English every evening.} — Projet décidé : \eng{be going to} + base verbale ; \eng{is} car le sujet est \eng{she}.
\item \eng{\textbf{If he had time}, he \textbf{would learn} technical vocabulary for his job.} — Hypothèse irréelle : \eng{if} + past simple (\eng{had}), \eng{would} + base verbale.
\item \eng{We \textbf{will attend} English classes after school next year.} — Annonce pour l'année prochaine : \eng{will} + base verbale ; ajout de \eng{next year} pour marquer le futur.
\item \eng{\textbf{If they find jobs abroad}, they \textbf{will use} English at work.} — Conséquence probable : \eng{if} + présent simple, \eng{will} + base verbale.
\end{enumerate}
\textbf{Point de vigilance :} dans les phrases 3 et 5, le verbe de la subordonnée en \eng{if} ne prend jamais \eng{will} ni \eng{would} ; c'est la principale qui porte la marque du futur ou du conditionnel.
\end{corrige}

\begin{corrige}[Exercice 9 — Compréhension et langue — texte type Bac]
\textbf{A. Comprehension questions — réponses attendues (formulations personnelles acceptées) :}
\begin{enumerate}
\item \eng{Many job advertisements mention that bilingual candidates are preferred.} — Reformulation de “bilingual candidates preferred”.
\item \eng{Clarisse dreams of becoming an air hostess.} — Attention : \eng{to dream of + -ing}. (Terme moderne : \eng{flight attendant}).
\item \eng{She reads a short article in English every morning, listens to podcasts on her way to school, and speaks English with her friends twice a week.} — Deux de ces trois activités suffisent.
\item \eng{A secretary needs writing skills for emails and reports.}
\item \eng{The three keys to success are regular practice, clear goals and self-confidence.}
\end{enumerate}

\textbf{B. Language work :}
\begin{enumerate}
\item (a) \eng{so} (« She knows that airlines require fluent English, \emph{so} she has built a personal plan ») — conséquence ; (b) \eng{they will be} (« the better prepared they \emph{will be} ») — futur ; (c) \eng{fluent} — « capable de parler facilement, courant ».
\item \eng{She \textbf{will read} one short article every morning.} — Futur simple : \eng{will} + base verbale ; attention, \eng{read} est irrégulier mais sa base verbale reste \eng{read}.
\item \eng{\textbf{If students practised regularly}, regular practice \textbf{would be} the key to success.} — Conditionnel irréel : past simple après \eng{if}, \eng{would} + \eng{be} dans la principale.
\item Traduction : « Plus tôt les élèves identifieront leurs besoins en langue, mieux ils seront préparés. » — Structure comparative proportionnelle \eng{the earlier..., the better prepared...} = « plus tôt..., mieux... ».
\end{enumerate}

\textbf{Barème indicatif (sur 20) :} A. Compréhension : 5 × 2 = 10 points (réponse correcte et en anglais : 2 pts ; réponse partielle ou maladroite : 1 pt). B. Langue : question 1 = 3 pts (1 pt par élément), questions 2 et 3 = 2 pts chacune, question 4 = 3 pts.

\textbf{Points de vigilance :} répondre par des phrases complètes, pas par un seul mot ; ne pas recopier intégralement le texte sans reformuler ; vérifier la concordance des temps dans les transformations.
\end{corrige}

\begin{corrige}[Exercice 10 — Traduction]
\textbf{Proposition de traduction :}

\eng{For my future career, English will be very important. First of all, I am going to take an English exam at the end of the year. Then, I will practise the language every day by reading newspapers and listening to the radio. However, I know that it will require a lot of effort. If I work regularly, I will be ready for the world of work.}

\textbf{Justifications et points de vigilance :}
\begin{itemize}
\item « je vais passer un examen » : projet décidé $\rightarrow$ \eng{I am going to take an English exam} ; \eng{to take an exam} = « passer un examen » (faux ami : \eng{pass an exam} = « réussir un examen »).
\item « en lisant... en écoutant... » : \eng{by + -ing} exprime le moyen.
\item « cela demandera beaucoup d'efforts » : \eng{it will require a lot of effort} ; \eng{effort} est ici indénombrable (singulier).
\item « Si je travaille régulièrement, je serai prêt » : condition réelle $\rightarrow$ \eng{if} + présent simple, \eng{will be} dans la principale ; jamais \eng{if I will work}.
\item Connecteurs : \eng{first of all} (« d'abord »), \eng{then} (« ensuite »), \eng{however} (« cependant ») suivis d'une virgule.
\end{itemize}
\textbf{Barème indicatif (sur 10) :} 2 pts par phrase (1 pt pour le sens, 1 pt pour la correction grammaticale : temps, connecteurs, vocabulaire).
\end{corrige}

\begin{corrige}[Exercice 11 — Sujet de composition type Bac]
\textbf{Proposition de rédaction modèle (environ 250 mots) :}

\eng{In today's world, English is no longer a luxury but a necessity. In Cameroon, most job advertisements now ask for bilingual candidates. As a final-year student, I believe that I must anticipate my second language needs before I enter the job market. This is why I have drawn up a clear personal plan.}

\eng{First of all, I am going to improve my vocabulary and my reading skills. I intend to read one English article every morning and to note down ten new words a day in a notebook. Moreover, I will listen to English radio programmes and podcasts on my way to school, because understanding spoken English is essential for interviews and meetings.}

\eng{However, reading and listening are not enough. My dream is to become a flight attendant, so I must speak fluently. Therefore, I have decided to join the English club of my school and to speak English with my classmates twice a week. If I practised every day, I would become fluent before the end of the year.}

\eng{Finally, I am going to take an international English exam next year. A certificate will prove my level to future employers and will make my application stronger.}

\eng{To sum up, anticipating my language needs means setting clear goals and practising regularly. I am confident that with effort and discipline, English will open many doors for my future career.}

\eng{(about 240 words)}

\textbf{Analyse du modèle :}
\begin{itemize}
\item \textbf{Introduction :} accroche (\eng{In today's world...}) + contexte camerounais + thesis statement (\eng{I have drawn up a clear personal plan}).
\item \textbf{Développement :} trois paragraphes, une idée par paragraphe (lecture/écoute ; expression orale ; certification), chacun ouvert par un connecteur.
\item \textbf{Connecteurs utilisés :} \eng{first of all, moreover, however, therefore, finally, to sum up} — six connecteurs, soit plus que le minimum exigé.
\item \textbf{Temps :} futur \eng{will} (\eng{I will listen}), projet \eng{be going to} (\eng{I am going to improve / to take}), conditionnel irréel (\eng{If I practised..., I would become...}).
\item \textbf{Conclusion :} bilan + ouverture optimiste, sans nouvel argument.
\end{itemize}

\textbf{Barème indicatif (sur 20) :}
\begin{center}
\begin{tabularx}{\linewidth}{|X|l|}
\hline
\rowcolor{popblueL} Critère & Points \\
\hline
Respect du sujet et de la tâche (anticipation des besoins en anglais au travail) & 4 \\
\hline
Structure : introduction, développement, conclusion & 4 \\
\hline
Connecteurs logiques (au moins 5, correctement employés) & 3 \\
\hline
Correction grammaticale (futur, conditionnel, ordre des mots) & 4 \\
\hline
Vocabulaire du monde du travail, richesse lexicale & 3 \\
\hline
Orthographe, ponctuation, présentation, nombre de mots & 2 \\
\hline
\end{tabularx}
\end{center}

\textbf{Points de vigilance pour le jour de l'examen :}
\begin{enumerate}
\item Lis le sujet deux fois et souligne les mots-clés ; fais un plan au brouillon (5 minutes).
\item Une idée principale par paragraphe, annoncée par une \eng{topic sentence}.
\item Vérifie chaque \eng{will} et \eng{would} : base verbale sans \eng{to} derrière.
\item Majuscules : \eng{English}, \eng{I}, débuts de phrases.
\item Garde 5 minutes pour la relecture : accords sujet-verbe, \eng{-s} à la 3\textsuperscript{e} personne du présent, orthographe (\eng{career}, \eng{because}, \eng{which}).
\item Indique le nombre de mots à la fin de la copie.
\end{enumerate}
\end{corrige}

\newpage

\coursec{Fiche bilan}

\begin{popfigure}
\begin{tikzpicture}[mindmap, grow cyclic, every node/.style={concept}, concept color=popblue, text=white,
root concept/.append style={concept color=popdark, font=\bfseries\small},
level 1/.append style={level distance=3.4cm, sibling angle=51.4, font=\small},
level 2/.append style={level distance=1.9cm, font=\footnotesize}]
\node[root concept, align=center]{Writing:\\ anticipating my\\ English needs\\ at work}
  child[concept color=popblue] { node[align=center]{1. Analyse\\ du sujet}
    child { node[align=center]{brainstorm} }
    child { node[align=center]{plan en\\ 3 parties} }
  }
  child[concept color=poporange] { node[align=center]{2. Structure}
    child { node[align=center]{intro} }
    child { node[align=center]{body} }
    child { node[align=center]{conclusion} }
  }
  child[concept color=popgreen] { node[align=center]{3. Connecteurs}
    child { node[align=center]{first, then,\\ finally} }
    child { node[align=center]{however,\\ therefore} }
  }
  child[concept color=popgold] { node[align=center]{4. Temps\\ verbaux}
    child { node[align=center]{future\\ will / going to} }
    child { node[align=center]{present\\ simple} }
  }
  child[concept color=poppurple] { node[align=center]{5. Vocabulaire\\ du travail}
    child { node[align=center]{skills,\\ training} }
    child { node[align=center]{job,\\ career} }
  }
  child[concept color=poppink] { node[align=center]{6. Modèle\\ commenté}
    child { node[align=center]{lire,\\ repérer,\\ imiter} }
  }
  child[concept color=popmuted] { node[align=center]{7. Relecture\\ et grille}
    child { node[align=center]{corriger\\ les erreurs} }
    child { node[align=center]{critères\\ du Bac} }
  };
\end{tikzpicture}
\legende{Carte mentale de la leçon : rédiger une composition sur ses besoins futurs en anglais au travail.}
\end{popfigure}

\begin{bilan}
\textbf{Lexique clé (English -- Français)}

\begin{itemize}
\item \eng{need} (n.) : besoin ; \eng{to anticipate} : anticiper, prévoir
\item \eng{second language} : langue seconde ; \eng{workplace} : lieu de travail
\item \eng{skill} : compétence ; \eng{training} : formation ; \eng{to improve} : améliorer
\item \eng{employer / employee} : employeur / employé ; \eng{career} : carrière
\item \eng{job interview} : entretien d'embauche ; \eng{promotion} : promotion, avancement
\item \eng{fluency} : aisance (à l'oral) ; \eng{communication} : communication
\item \eng{opportunity} : opportunité ; \eng{requirement} : exigence
\item \eng{to apply for a job} : postuler à un emploi ; \eng{to attend a course} : suivre un cours
\end{itemize}

\textbf{Règles de grammaire à connaître par cœur}

\begin{center}
\begin{tabularx}{\linewidth}{|X|X|X|}
\hline
\rowcolor{popblueL} Point & Forme & Exemple \\
\hline
Futur avec \eng{will} & will + base verbale (prédiction, décision) & \eng{I will need English at work.} \\
\hline
Futur avec \eng{going to} & be + going to + BV (projet, intention) & \eng{I am going to take English courses.} \\
\hline
Présent simple & vérités générales, faits & \eng{English opens many doors.} \\
\hline
\eng{need to} + BV & exprimer la nécessité & \eng{Workers need to speak English.} \\
\hline
\eng{in order to / so as to} & exprimer le but & \eng{I study hard in order to succeed.} \\
\hline
\eng{if} + présent, will + BV & conditionnel de type 1 & \eng{If I master English, I will get promoted.} \\
\hline
\end{tabularx}
\end{center}

\textbf{Connecteurs logiques essentiels}

\begin{itemize}
\item Ajout : \eng{first(ly), moreover, in addition, besides}
\item Opposition : \eng{however, whereas, although, on the other hand}
\item Cause/conséquence : \eng{because, since, therefore, as a result, so}
\item But : \eng{in order to, so that, for this purpose}
\item Conclusion : \eng{to conclude, in conclusion, all in all, to sum up}
\end{itemize}

\textbf{Méthode express : les 5 étapes de la composition}

\begin{enumerate}
\item \textbf{Lire} le sujet deux fois et souligner les mots-clés.
\item \textbf{Brainstormer} : noter 6 à 8 idées en vrac, puis en choisir 3 ou 4.
\item \textbf{Planifier} : introduction (sujet + problématique), 2 ou 3 paragraphes de développement (une idée + un exemple chacun), conclusion (bilan + ouverture).
\item \textbf{Rédiger} avec connecteurs, temps corrects et vocabulaire du thème.
\item \textbf{Relire} : accords, temps, orthographe, ponctuation, majuscules.
\end{enumerate}
\end{bilan}

\begin{aretenir}[Checklist avant le Bac]
Avant de rendre ma copie, je vérifie que :
\begin{itemize}
\item[$\square$] J'ai bien compris et traité le sujet posé (pas de hors-sujet).
\item[$\square$] Ma composition a une introduction, un développement et une conclusion clairement séparés.
\item[$\square$] Chaque paragraphe contient une idée principale et au moins un exemple.
\item[$\square$] J'ai utilisé au moins cinq connecteurs logiques variés (\eng{first, moreover, however, therefore, in conclusion}).
\item[$\square$] J'ai employé correctement \eng{will} et \eng{going to} pour parler de mes projets professionnels.
\item[$\square$] J'ai utilisé le vocabulaire du travail (\eng{skills, training, career, employer, opportunity}).
\item[$\square$] Mes verbes sont conjugués au bon temps et accordés avec leur sujet (\eng{he needs}, pas \eng{he need}).
\item[$\square$] Je n'ai pas traduit mot à mot le français (\eng{I have 18 years} est faux : \eng{I am 18}).
\item[$\square$] Mes phrases commencent par une majuscule et se terminent par un point.
\item[$\square$] J'ai vérifié l'orthographe des mots difficiles (\eng{business, necessary, opportunity}).
\item[$\square$] J'ai relu ma copie au moins une fois en entier.
\item[$\square$] Ma présentation est soignée : paragraphes visibles, écriture lisible, pas de ratures excessives.
\end{itemize}
\end{aretenir}

\findecours

\end{document}
